
\documentclass[aspectratio=169,11pt]{beamer}
\usepackage[utf8]{inputenc}
\usepackage[T1]{fontenc}
\usepackage{lmodern}
\usepackage{graphicx,amsfonts,amsmath,amssymb,bm}
\usetheme{Madrid}
\setbeamertemplate{navigation symbols}{}
\setbeamertemplate{footline}[frame number]
\setbeamercovered{transparent}
\usecolortheme{default}
\definecolor{AccentBlue}{RGB}{28,76,160}
\setbeamercolor{structure}{fg=AccentBlue}
\setbeamercolor{frametitle}{fg=white,bg=AccentBlue}
\setbeamercolor{title}{fg=AccentBlue}
\setbeamertemplate{itemize item}{\color{AccentBlue}$\blacktriangleright$}
\setbeamertemplate{itemize subitem}{\color{AccentBlue}\scriptsize$\blacktriangleright$}

\title{Artificial Intelligence:\\A Mathematical and Engineering Perspective}
\subtitle{Slide-Style Lecture Notes}
\author{}
\date{March 2026}

\newcommand{\SectionDivider}[1]{
\begin{frame}[plain,noframenumbering]
\centering
\vfill
{\usebeamerfont{title}\color{AccentBlue}\LARGE #1\par}
\vspace{0.6em}
{\large Slide-style lecture note conversion}
\vfill
\end{frame}
}

\begin{document}

\begin{frame}[plain,noframenumbering]
\centering
\vfill
{\usebeamerfont{title}\color{AccentBlue}\LARGE Artificial Intelligence:\\A Mathematical and Engineering Perspective\par}
\vspace{0.8em}
{\Large Slide-Style Lecture Notes\par}
\vspace{1.2em}
{March 2026\par}
\vfill
\end{frame}

\begin{frame}{How to use these slides}
\small
\begin{itemize}
    \item Each lecture or tutorial starts with a divider slide and a short outline.
    \item Most original subsections were converted directly into one or more Beamer frames.
    \item Longer content uses automatic frame breaks so equations and lists stay readable.
    \item The source material was preserved as much as possible while changing the layout to a slide format.
\end{itemize}
\end{frame}
\section{Lecture 1: Learning as Function Approximation}
\SectionDivider{Lecture 1: Learning as Function Approximation}
\begin{frame}[t,allowframebreaks]{Outline}
\small
\begin{itemize}
  \item Motivation
  \item The Learning Problem
  \item Examples of Function Approximation
  \item Approximation, Estimation, and Optimization
  \item Geometry of Function Approximation
  \item Limitations of Finite Data
  \item A Unifying Perspective
  \item Outlook
\end{itemize}
\end{frame}

\subsection{Motivation}
\begin{frame}[t,allowframebreaks]{Motivation}
\small
At its core, machine learning seeks to extract patterns from data in order to make predictions about unseen instances. This goal can be formulated in a unified and mathematically precise way by viewing learning as a problem of function approximation.

In many applications, we are given observations of the form:
\[
(x_i, y_i), \quad i = 1, \dots, n,
\]
where $x_i \in \mathcal{X}$ represents an input and $y_i \in \mathcal{Y}$ represents a corresponding output.

The central assumption is that there exists an (unknown) relationship between inputs and outputs:
\[
y = f(x),
\]
for some function $f : \mathcal{X} \to \mathcal{Y}$.

The task of learning is to approximate this unknown function using the observed data.

\textbf{Guiding idea.}\\ Machine learning is the problem of constructing an approximation to an unknown function based on finite data.
\end{frame}

\subsection{The Learning Problem}
\begin{frame}[t,allowframebreaks]{The Learning Problem}
\small
We formalize the learning problem by introducing a parameterized family of functions:
\[
\{ f_\theta : \mathcal{X} \to \mathcal{Y} \mid \theta \in \Theta \},
\]
called a \emph{model class} or \emph{hypothesis space}.

Given data $\{(x_i, y_i)\}_{i=1}^n$, we seek parameters $\theta$ such that $f_\theta$ approximates the unknown function $f$.

To quantify approximation quality, we introduce a loss function:
\[
\ell : \mathcal{Y} \times \mathcal{Y} \to \mathbb{R}.
\]

The empirical learning problem becomes:
\[
\min_{\theta \in \Theta} \; \frac{1}{n} \sum_{i=1}^n \ell\big(f_\theta(x_i), y_i\big).
\]

\textbf{Interpretation.}\\ \begin{itemize}
    \item $f_\theta$ is our approximation to the unknown function $f$
    \item The loss measures discrepancy between predictions and observations
    \item Learning corresponds to selecting the best approximation within the model class
\end{itemize}
\end{frame}

\subsection{Examples of Function Approximation}
\begin{frame}[t,allowframebreaks]{Examples of Function Approximation}{Linear Regression}
\small
In linear regression, we assume:
\[
f_\theta(x) = w^\top x,
\]
where $\theta = w \in \mathbb{R}^d$.

Using squared loss:
\[
\ell(f(x), y) = (f(x) - y)^2,
\]
we obtain the objective:
\[
\min_w \sum_{i=1}^n (w^\top x_i - y_i)^2.
\]

This corresponds to approximating $f$ by a linear function.
\end{frame}

\begin{frame}[t,allowframebreaks]{Examples of Function Approximation}{Classification}
\small
In classification, outputs belong to discrete classes:
\[
y \in \{1, \dots, K\}.
\]

We may model:
\[
f_\theta(x) = \arg\max_k p_\theta(y = k \mid x).
\]

Here, the function approximation perspective remains valid: we approximate the mapping from inputs to labels.
\end{frame}

\begin{frame}[t,allowframebreaks]{Examples of Function Approximation}{Nonlinear Models}
\small
More expressive models introduce nonlinear transformations:
\[
f_\theta(x) = \sigma(Wx + b),
\]
or compositions of such transformations (as in neural networks).

These allow approximation of more complex functions.
\end{frame}

\subsection{Approximation, Estimation, and Optimization}
\begin{frame}[t,allowframebreaks]{Approximation, Estimation, and Optimization}
\small
The learning problem can be decomposed into three interacting components:

\textbf{1. Approximation.}\\ The hypothesis space determines which functions can be represented.

\textbf{2. Estimation.}\\ Given finite data, we estimate parameters $\theta$.

\textbf{3. Optimization.}\\ We solve a numerical problem to find the best parameters.

\textbf{Insight.}\\ Machine learning is not a single problem, but the interaction of approximation, estimation, and optimization.
\end{frame}

\subsection{Geometry of Function Approximation}
\begin{frame}[t,allowframebreaks]{Geometry of Function Approximation}
\small
It is often useful to interpret learning geometrically.

Let $\mathcal{H} = \{f_\theta\}$ denote the hypothesis space. Then learning can be viewed as finding:
\[
f_\theta \in \mathcal{H}
\]
that best matches observed data.

In linear regression, this corresponds to projecting the target vector $y$ onto the subspace spanned by input features.

\textbf{Geometric picture.}\\ \begin{itemize}
    \item Data defines a target point
    \item Model class defines a set (or manifold)
    \item Learning finds the closest point in that set
\end{itemize}
\end{frame}

\subsection{Limitations of Finite Data}
\begin{frame}[t,allowframebreaks]{Limitations of Finite Data}
\small
Since we only observe finitely many samples, we cannot recover the true function $f$ exactly.

Two fundamental challenges arise:
\begin{itemize}
    \item \textbf{Underdetermination:}\\ many functions may fit the data
    \item \textbf{Generalization:}\\ performance on unseen data is not guaranteed
\end{itemize}

This motivates the need for:
\begin{itemize}
    \item appropriate model classes
    \item regularization
    \item probabilistic reasoning
\end{itemize}
\end{frame}

\subsection{A Unifying Perspective}
\begin{frame}[t,allowframebreaks]{A Unifying Perspective}
\small
Learning as function approximation provides a unifying framework for machine learning:

\begin{itemize}
    \item \textbf{Supervised learning:}\\ approximate input-output mappings
    \item \textbf{Classification:}\\ approximate decision functions
    \item \textbf{Regression:}\\ approximate real-valued functions
    \item \textbf{Deep learning:}\\ approximate complex compositions of functions
\end{itemize}

\textbf{Core abstraction.}\\ All learning problems can be viewed as selecting a function from a hypothesis space to approximate an unknown target.

\textbf{Key insight.}\\ The choice of function class determines what can be learned, while the data determines what is actually learned.
\end{frame}

\subsection{Outlook}
\begin{frame}[t,allowframebreaks]{Outlook}
\small
In this lecture, we introduced the central paradigm of learning as function approximation.

However, this formulation raises several important questions:

\begin{itemize}
    \item How is the data generated?
    \item What does it mean for a model to generalize?
    \item How should we choose the loss function?
\end{itemize}

In the next lecture, we address these questions by adopting a probabilistic perspective, where data is modeled as random variables and learning is framed as statistical estimation.

\textbf{Guiding principle.}\\ Understanding learning requires combining approximation theory with probabilistic reasoning.
\end{frame}

\section{Lecture 2: A Probabilistic View of Learning}
\SectionDivider{Lecture 2: A Probabilistic View of Learning}
\begin{frame}[t,allowframebreaks]{Outline}
\small
\begin{itemize}
  \item Motivation
  \item Data as Random Variables
  \item Risk and Expected Loss
  \item Empirical Risk Minimization
  \item Conditional Distributions and Prediction
  \item Bayes’ Theorem
  \item Loss Functions and Probabilistic Modeling
  \item Uncertainty and Noise
  \item A Unifying Perspective
  \item Outlook
\end{itemize}
\end{frame}

\subsection{Motivation}
\begin{frame}[t,allowframebreaks]{Motivation}
\small
In the previous lecture, we formulated learning as the problem of approximating an unknown function from data. However, this formulation leaves open a fundamental question:

\textit{Where does the data come from?}

In practice, data is not deterministic but arises from a stochastic process. Even if there exists an underlying functional relationship, observations are typically corrupted by noise or variability.

To capture this, we adopt a probabilistic framework in which both inputs and outputs are modeled as random variables.

\textbf{Guiding idea.}\\ Learning is the problem of estimating relationships between random variables based on finite samples.
\end{frame}

\subsection{Data as Random Variables}
\begin{frame}[t,allowframebreaks]{Data as Random Variables}
\small
We assume that data points are drawn independently from an unknown joint distribution:
\[
(x, y) \sim P,
\]
where $P$ is a probability distribution on $\mathcal{X} \times \mathcal{Y}$.

The dataset $\{(x_i, y_i)\}_{i=1}^n$ is thus viewed as an i.i.d.\ sample from $P$.

\textbf{Interpretation.}\\ \begin{itemize}
    \item The true relationship between $x$ and $y$ is encoded in $P$
    \item The function $f$ from the previous lecture is now understood as a conditional expectation or decision rule derived from $P$
\end{itemize}
\end{frame}

\subsection{Risk and Expected Loss}
\begin{frame}[t,allowframebreaks]{Risk and Expected Loss}
\small
Given a model $f : \mathcal{X} \to \mathcal{Y}$, we define its \emph{risk} as the expected loss:
\[
R(f) = \mathbb{E}_{(x,y) \sim P} \left[ \ell(f(x), y) \right].
\]

\textbf{Goal.}\\ Find a function $f$ that minimizes the risk:
\[
f^* = \arg\min_f R(f).
\]

\textbf{Bayes optimal predictor.}\\ The function $f^*$ is called the \emph{Bayes predictor}.

\textbf{Example (squared loss).}\\ If $\ell(f(x), y) = (f(x) - y)^2$, then:
\[
f^*(x) = \mathbb{E}[y \mid x].
\]

\textbf{Interpretation.}\\ \begin{itemize}
    \item The optimal prediction is the conditional expectation
    \item Learning seeks to approximate this conditional structure
\end{itemize}
\end{frame}

\subsection{Empirical Risk Minimization}
\begin{frame}[t,allowframebreaks]{Empirical Risk Minimization}
\small
Since the true distribution $P$ is unknown, we cannot compute $R(f)$ directly.

Instead, we approximate it using the empirical risk:
\[
\hat{R}_n(f) = \frac{1}{n} \sum_{i=1}^n \ell(f(x_i), y_i).
\]

\textbf{Empirical Risk Minimization (ERM).}\\ \[
\hat{f} = \arg\min_{f \in \mathcal{H}} \hat{R}_n(f),
\]
where $\mathcal{H}$ is the hypothesis space.

\textbf{Justification.}\\ \begin{itemize}
    \item By the law of large numbers, $\hat{R}_n(f)$ approximates $R(f)$
    \item Minimizing empirical risk approximates minimizing true risk
\end{itemize}

\textbf{Key challenge.}\\ The empirical minimizer may not generalize well to unseen data.
\end{frame}

\subsection{Conditional Distributions and Prediction}
\begin{frame}[t,allowframebreaks]{Conditional Distributions and Prediction}
\small
The probabilistic framework emphasizes modeling the conditional distribution:
\[
p(y \mid x).
\]

\textbf{Prediction strategies.}\\ \begin{itemize}
    \item Regression: predict $\mathbb{E}[y \mid x]$
    \item Classification: predict $\arg\max_y p(y \mid x)$
\end{itemize}

\textbf{Example (classification).}\\ For $y \in \{1, \dots, K\}$:
\[
f^*(x) = \arg\max_k p(y = k \mid x).
\]

\textbf{Insight.}\\ Learning can be viewed as estimating conditional distributions or functions derived from them.
\end{frame}

\subsection{Bayes’ Theorem}
\begin{frame}[t,allowframebreaks]{Bayes’ Theorem}
\small
Bayes’ theorem provides a fundamental relationship between conditional and marginal probabilities:
\[
p(y \mid x) = \frac{p(x \mid y)\, p(y)}{p(x)}.
\]

\textbf{Components.}\\ \begin{itemize}
    \item $p(y)$: prior distribution
    \item $p(x \mid y)$: likelihood
    \item $p(y \mid x)$: posterior
    \item $p(x)$: evidence
\end{itemize}

\textbf{Interpretation.}\\ \begin{itemize}
    \item Prior encodes initial beliefs
    \item Likelihood captures how data is generated
    \item Posterior updates beliefs given observations
\end{itemize}

\textbf{Bayesian viewpoint.}\\ Learning is the process of updating beliefs about unknown quantities using observed data.
\end{frame}

\subsection{Loss Functions and Probabilistic Modeling}
\begin{frame}[t,allowframebreaks]{Loss Functions and Probabilistic Modeling}
\small
Loss functions often arise naturally from probabilistic assumptions.

\textbf{Negative log-likelihood.}\\ If we model $p_\theta(y \mid x)$, we may define:
\[
\ell(f_\theta(x), y) = -\log p_\theta(y \mid x).
\]

\textbf{Result.}\\ \[
\min_\theta \sum_i -\log p_\theta(y_i \mid x_i)
\]
corresponds to maximum likelihood estimation.

\textbf{Examples.}\\ \begin{itemize}
    \item Gaussian noise $\rightarrow$ squared loss
    \item Bernoulli model $\rightarrow$ cross-entropy loss
\end{itemize}

\textbf{Insight.}\\ Loss functions encode assumptions about the data-generating process.
\end{frame}

\subsection{Uncertainty and Noise}
\begin{frame}[t,allowframebreaks]{Uncertainty and Noise}
\small
The probabilistic framework explicitly accounts for uncertainty.

\textbf{Sources of uncertainty.}\\ \begin{itemize}
    \item Noise in observations
    \item Incomplete information
    \item Model limitations
\end{itemize}

\textbf{Prediction with uncertainty.}\\ \begin{itemize}
    \item Instead of predicting a single value, we may predict a distribution
    \item Confidence can be quantified
\end{itemize}

\textbf{Insight.}\\ Probability provides a natural language for reasoning under uncertainty.
\end{frame}

\subsection{A Unifying Perspective}
\begin{frame}[t,allowframebreaks]{A Unifying Perspective}
\small
The probabilistic view reframes learning as statistical estimation:

\begin{itemize}
    \item \textbf{Data:}\\ samples from an unknown distribution
    \item \textbf{Model:}\\ approximation of conditional structure
    \item \textbf{Objective:}\\ minimize expected loss
\end{itemize}

\textbf{Core abstraction.}\\ Learning is the estimation of functions or distributions from random data.

\textbf{Key insight.}\\ The deterministic function approximation view is a special case of a broader probabilistic framework.
\end{frame}

\subsection{Outlook}
\begin{frame}[t,allowframebreaks]{Outlook}
\small
In this lecture, we introduced a probabilistic formulation of learning, including:

\begin{itemize}
    \item data as random variables,
    \item risk and empirical risk minimization,
    \item conditional distributions,
    \item Bayes’ theorem.
\end{itemize}

This perspective raises further questions:

\begin{itemize}
    \item How should we choose the hypothesis space?
    \item What assumptions guide learning?
    \item How do we balance flexibility and generalization?
\end{itemize}

In the next lecture, we study hypothesis spaces and inductive bias, which determine the structure of the functions we learn.

\textbf{Guiding principle.}\\ Learning is fundamentally a statistical problem: we infer structure from randomness using limited data.
\end{frame}

\section{Tutorial 1: Foundations}
\SectionDivider{Tutorial 1: Foundations}
\begin{frame}[t,allowframebreaks]{Outline}
\small
\begin{itemize}
  \item Objectives
  \item Exercise 1: Empirical Risk vs True Risk
  \item Exercise 2: Bias--Variance Intuition
  \item Exercise 3: Loss Functions and Probabilistic Interpretation
  \item Exercise 4: Simple Optimization Problem
  \item Exercise 5 (Discussion): What Enables Generalization?
  \item Summary
\end{itemize}
\end{frame}

\subsection{Objectives}
\begin{frame}[t,allowframebreaks]{Objectives}
\small
This tutorial reinforces the key ideas from Lectures 1 and 2:
\begin{itemize}
    \item Learning as function approximation
    \item Probabilistic formulation of learning
    \item Empirical risk minimization (ERM)
    \item Bias--variance tradeoff (intuition)
\end{itemize}

\textbf{Focus.}\\ Develop intuition through concrete examples and simple derivations.
\end{frame}

\subsection{Exercise 1: Empirical Risk vs True Risk}
\begin{frame}[t,allowframebreaks]{Exercise 1: Empirical Risk vs True Risk}
\small
Let $(x,y) \sim P$, and let $f$ be a fixed model.

\textbf{Tasks.}\\ \begin{enumerate}
    \item Define the true risk $R(f)$.
    \item Define the empirical risk $\hat{R}_n(f)$.
    \item Explain why $\hat{R}_n(f)$ is a good approximation of $R(f)$.
    \item Give an example where $\hat{R}_n(f)$ is misleading.
\end{enumerate}

\textbf{Solution.}\\ \begin{enumerate}
    \item 
    \[
    R(f) = \mathbb{E}_{(x,y)\sim P}[\ell(f(x),y)].
    \]

    \item 
    \[
    \hat{R}_n(f) = \frac{1}{n} \sum_{i=1}^n \ell(f(x_i), y_i).
    \]

    \item By the law of large numbers:
    \[
    \hat{R}_n(f) \to R(f) \quad \text{as } n \to \infty.
    \]

    \item If $f$ overfits the training data (e.g., memorizes noise), then:
    \[
    \hat{R}_n(f) \ll R(f).
    \]
\end{enumerate}

\textbf{Key takeaway.}\\ Empirical risk approximates true risk, but only when the model generalizes.
\end{frame}

\subsection{Exercise 2: Bias--Variance Intuition}
\begin{frame}[t,allowframebreaks]{Exercise 2: Bias--Variance Intuition}
\small
Consider fitting a model to noisy data:
\[
y = f(x) + \varepsilon, \quad \varepsilon \sim \mathcal{N}(0, \sigma^2).
\]

\textbf{Tasks.}\\ \begin{enumerate}
    \item Describe what happens when the model is too simple.
    \item Describe what happens when the model is too complex.
    \item Explain the bias--variance tradeoff qualitatively.
\end{enumerate}

\textbf{Solution.}\\ \begin{enumerate}
    \item \textbf{Too simple (high bias):}\\ \begin{itemize}
        \item Model cannot capture true structure
        \item Systematic error remains
    \end{itemize}

    \item \textbf{Too complex (high variance):}\\ \begin{itemize}
        \item Model fits noise
        \item Predictions vary significantly with data
    \end{itemize}

    \item \textbf{Tradeoff:}\\ \begin{itemize}
        \item Increasing complexity reduces bias
        \item Increasing complexity increases variance
    \end{itemize}
\end{enumerate}

\textbf{Key takeaway.}\\ Good models balance approximation accuracy and stability.
\end{frame}

\subsection{Exercise 3: Loss Functions and Probabilistic Interpretation}
\begin{frame}[t,allowframebreaks]{Exercise 3: Loss Functions and Probabilistic Interpretation}
\small
Suppose we assume:
\[
y = f_\theta(x) + \varepsilon, \quad \varepsilon \sim \mathcal{N}(0, \sigma^2).
\]

\textbf{Tasks.}\\ \begin{enumerate}
    \item Write down the likelihood $p(y \mid x, \theta)$.
    \item Show that maximizing likelihood is equivalent to minimizing squared loss.
\end{enumerate}

\textbf{Solution.}\\ \begin{enumerate}
    \item 
    \[
    p(y \mid x, \theta)
    =
    \frac{1}{\sqrt{2\pi \sigma^2}}
    \exp\left(-\frac{(y - f_\theta(x))^2}{2\sigma^2}\right).
    \]

    \item Taking negative log-likelihood:
    \[
    -\log p(y \mid x, \theta)
    =
    \frac{(y - f_\theta(x))^2}{2\sigma^2} + C.
    \]

    Minimizing this is equivalent to minimizing:
    \[
    (y - f_\theta(x))^2.
    \]
\end{enumerate}

\textbf{Key takeaway.}\\ Loss functions encode assumptions about the data-generating process.
\end{frame}

\subsection{Exercise 4: Simple Optimization Problem}
\begin{frame}[t,allowframebreaks]{Exercise 4: Simple Optimization Problem}
\small
Consider linear regression:
\[
L(w) = \sum_{i=1}^n (w^\top x_i - y_i)^2.
\]

\textbf{Tasks.}\\ \begin{enumerate}
    \item Compute the gradient $\nabla_w L(w)$.
    \item Write the gradient descent update rule.
\end{enumerate}

\textbf{Solution.}\\ \begin{enumerate}
    \item 
    \[
    \nabla_w L(w) = 2 \sum_{i=1}^n (w^\top x_i - y_i) x_i.
    \]

    In matrix form:
    \[
    \nabla_w L(w) = 2 X^\top (Xw - y).
    \]

    \item Gradient descent update:
    \[
    w_{t+1} = w_t - \eta \nabla_w L(w_t).
    \]
\end{enumerate}

\textbf{Key takeaway.}\\ Learning reduces to solving an optimization problem.
\end{frame}

\subsection{Exercise 5 (Discussion): What Enables Generalization?}
\begin{frame}[t,allowframebreaks]{Exercise 5 (Discussion): What Enables Generalization?}
\small
\textbf{Question.}\\ If many functions fit the data perfectly, why do some generalize better than others?

\textbf{Discussion points.}\\ \begin{itemize}
    \item Hypothesis space restrictions
    \item Inductive bias
    \item Simplicity of solutions
\end{itemize}

\textbf{Insight.}\\ Generalization depends on assumptions beyond data.
\end{frame}

\subsection{Summary}
\begin{frame}[t,allowframebreaks]{Summary}
\small
\begin{itemize}
    \item ERM approximates true risk but may fail under overfitting
    \item Bias--variance tradeoff governs model complexity
    \item Loss functions arise from probabilistic assumptions
    \item Learning involves optimization over parameters
\end{itemize}

\textbf{Preparation for next lecture.}\\ We now examine how hypothesis spaces and inductive bias formalize these ideas.
\end{frame}

\section{Lecture 3: Hypothesis Spaces and Inductive Bias}
\SectionDivider{Lecture 3: Hypothesis Spaces and Inductive Bias}
\begin{frame}[t,allowframebreaks]{Outline}
\small
\begin{itemize}
  \item Motivation
  \item Function Classes and Model Families
  \item Inductive Bias and Prior Structure
  \item Bias--Variance Tradeoff (Intuition)
  \item Choosing Representations
  \item A Unifying Perspective
  \item Outlook
\end{itemize}
\end{frame}

\subsection{Motivation}
\begin{frame}[t,allowframebreaks]{Motivation}
\small
In the previous lectures, we formulated learning as:
\begin{itemize}
    \item approximating an unknown function,
    \item minimizing expected loss under a data distribution.
\end{itemize}

However, a fundamental question remains:

\textit{Among all possible functions, which ones should we consider?}

Since we only observe finite data, many functions can fit the data equally well. Without further assumptions, learning is ill-posed.

\textbf{Guiding idea.}\\ Learning requires restricting the space of possible functions and introducing assumptions about the structure of the problem.
\end{frame}

\subsection{Function Classes and Model Families}
\begin{frame}[t,allowframebreaks]{Function Classes and Model Families}
\small
We define a \emph{hypothesis space}:
\[
\mathcal{H} = \{ f_\theta : \mathcal{X} \to \mathcal{Y} \mid \theta \in \Theta \}.
\]

\textbf{Examples.}\\ \begin{itemize}
    \item \textbf{Linear models:}\\ \[
    f(x) = w^\top x
    \]

    \item \textbf{Polynomial models:}\\ \[
    f(x) = \sum_{k=0}^d a_k x^k
    \]

    \item \textbf{Neural networks:}\\ \[
    f(x) = \sigma(W_2 \sigma(W_1 x))
    \]
\end{itemize}

\textbf{Interpretation.}\\ \begin{itemize}
    \item $\mathcal{H}$ defines which functions are representable
    \item Learning selects a function within $\mathcal{H}$
\end{itemize}

\textbf{Expressivity.}\\ A richer hypothesis space can approximate more complex functions.

\textbf{Tradeoff.}\\ \begin{itemize}
    \item Small $\mathcal{H}$: limited expressivity
    \item Large $\mathcal{H}$: risk of overfitting
\end{itemize}
\end{frame}

\subsection{Inductive Bias and Prior Structure}
\begin{frame}[t,allowframebreaks]{Inductive Bias and Prior Structure}
\small
The restriction to a hypothesis space introduces an \emph{inductive bias}.

\textbf{Definition.}\\ Inductive bias is the set of assumptions that a learning algorithm uses to generalize beyond observed data.

\textbf{Sources of inductive bias.}\\ \begin{itemize}
    \item Choice of hypothesis space
    \item Regularization
    \item Optimization algorithm
\end{itemize}

\textbf{Example.}\\ \begin{itemize}
    \item Linear models assume linear relationships
    \item Convolutional networks assume translation invariance
\end{itemize}

\textbf{Prior structure.}\\ Inductive bias can be interpreted probabilistically as a prior over functions:
\[
p(f).
\]

\textbf{Insight.}\\ Learning is not purely data-driven; it depends critically on assumptions encoded in the model.
\end{frame}

\subsection{Bias--Variance Tradeoff (Intuition)}
\begin{frame}[t,allowframebreaks]{Bias--Variance Tradeoff (Intuition)}
\small
A central concept in learning is the tradeoff between bias and variance.

\textbf{Bias.}\\ \begin{itemize}
    \item Error due to model assumptions
    \item High when the model is too simple
\end{itemize}

\textbf{Variance.}\\ \begin{itemize}
    \item Sensitivity to data fluctuations
    \item High when the model is too complex
\end{itemize}

\textbf{Qualitative decomposition.}\\ \[
\text{Error} = \text{Bias}^2 + \text{Variance} + \text{Noise}.
\]

\textbf{Interpretation.}\\ \begin{itemize}
    \item Increasing model complexity reduces bias
    \item Increasing model complexity increases variance
\end{itemize}

\textbf{Example.}\\ \begin{itemize}
    \item Linear model on nonlinear data: high bias
    \item Highly flexible model on small dataset: high variance
\end{itemize}

\textbf{Insight.}\\ Good generalization requires balancing simplicity and flexibility.
\end{frame}

\subsection{Choosing Representations}
\begin{frame}[t,allowframebreaks]{Choosing Representations}
\small
The representation of data plays a central role in learning.

\textbf{Feature representation.}\\ \begin{itemize}
    \item Input $x$ may be transformed into features $\phi(x)$
    \item Learning is performed on $\phi(x)$ rather than $x$
\end{itemize}

\textbf{Examples.}\\ \begin{itemize}
    \item Polynomial features
    \item Kernel mappings
    \item Learned representations (neural networks)
\end{itemize}

\textbf{Key idea.}\\ The difficulty of learning depends on the representation.

\textbf{Linearization principle.}\\ A complex problem may become simple under the right representation.

\textbf{Modern perspective.}\\ \begin{itemize}
    \item Deep learning automatically learns representations
    \item Representation learning replaces manual feature design
\end{itemize}

\textbf{Insight.}\\ Choosing or learning representations is often more important than the choice of model.
\end{frame}

\subsection{A Unifying Perspective}
\begin{frame}[t,allowframebreaks]{A Unifying Perspective}
\small
Hypothesis spaces and inductive bias define the structure of learning:

\begin{itemize}
    \item \textbf{Hypothesis space:}\\ what functions are possible
    \item \textbf{Inductive bias:}\\ which functions are preferred
    \item \textbf{Data:}\\ which functions are consistent with observations
\end{itemize}

\textbf{Core idea.}\\ Learning is the process of selecting a function that balances data fit with prior assumptions.

\textbf{Key insight.}\\ Generalization is possible only because of inductive bias; without it, learning from finite data would be impossible.
\end{frame}

\subsection{Outlook}
\begin{frame}[t,allowframebreaks]{Outlook}
\small
In this lecture, we examined:

\begin{itemize}
    \item hypothesis spaces and model families,
    \item inductive bias and prior structure,
    \item the bias--variance tradeoff,
    \item the role of representations.
\end{itemize}

These ideas naturally lead to the next question:

\textit{How do we actually find the best function within a hypothesis space?}

In the next lecture, we study learning as an optimization problem, where parameter estimation is framed as minimizing an objective function.

\textbf{Guiding principle.}\\ The success of learning depends as much on the choice of representation and inductive bias as on the data itself.
\end{frame}

\section{Lecture 4: Learning as Optimization}
\SectionDivider{Lecture 4: Learning as Optimization}
\begin{frame}[t,allowframebreaks]{Outline}
\small
\begin{itemize}
  \item Motivation
  \item Optimization Formulations of Learning
  \item Gradient-Based Methods
  \item Stochastic Optimization
  \item Geometry of Loss Landscapes
  \item A Unifying Perspective
  \item Outlook
\end{itemize}
\end{frame}

\subsection{Motivation}
\begin{frame}[t,allowframebreaks]{Motivation}
\small
In previous lectures, we developed three key ideas:
\begin{itemize}
    \item Learning as function approximation
    \item Learning as statistical estimation
    \item The role of hypothesis spaces and inductive bias
\end{itemize}

These ideas naturally lead to a central computational question:

\textit{How do we actually find the best function within a hypothesis space?}

The answer is to formulate learning as an optimization problem.

\textbf{Guiding idea.}\\ Learning is the process of minimizing an objective function defined by data and model assumptions.
\end{frame}

\subsection{Optimization Formulations of Learning}
\begin{frame}[t,allowframebreaks]{Optimization Formulations of Learning}
\small
Given a hypothesis space $\mathcal{H} = \{f_\theta\}$ and a loss function $\ell$, we define the empirical objective:
\[
L(\theta) = \frac{1}{n} \sum_{i=1}^n \ell(f_\theta(x_i), y_i).
\]

\textbf{Learning problem.}\\ \[
\min_{\theta \in \Theta} L(\theta).
\]

\textbf{Regularized formulation.}\\ \[
\min_\theta \; L(\theta) + \lambda \Omega(\theta),
\]
where $\Omega(\theta)$ is a regularization term.

\textbf{Examples.}\\ \begin{itemize}
    \item Linear regression: quadratic objective
    \item Logistic regression: convex but nonlinear objective
    \item Neural networks: highly non-convex objective
\end{itemize}

\textbf{Interpretation.}\\ \begin{itemize}
    \item Objective function encodes both data fit and inductive bias
    \item Learning reduces to searching for optimal parameters
\end{itemize}
\end{frame}

\subsection{Gradient-Based Methods}
\begin{frame}[t,allowframebreaks]{Gradient-Based Methods}
\small
Many learning problems are solved using gradient-based optimization.

\textbf{Gradient descent.}\\ \[
\theta_{t+1} = \theta_t - \eta \nabla L(\theta_t),
\]
where $\eta > 0$ is the step size.

\textbf{Interpretation.}\\ \begin{itemize}
    \item The gradient points in the direction of steepest increase
    \item We move in the opposite direction to decrease the loss
\end{itemize}

\textbf{First-order approximation.}\\ \[
L(\theta + \delta) \approx L(\theta) + \nabla L(\theta)^\top \delta.
\]

\textbf{Choice of step size.}\\ \begin{itemize}
    \item Too large: instability
    \item Too small: slow convergence
\end{itemize}

\textbf{Variants.}\\ \begin{itemize}
    \item Momentum methods
    \item Adaptive methods (e.g., scaling by past gradients)
\end{itemize}

\textbf{Insight.}\\ Gradient methods exploit local information to iteratively improve solutions.
\end{frame}

\subsection{Stochastic Optimization}
\begin{frame}[t,allowframebreaks]{Stochastic Optimization}
\small
In large-scale problems, computing the full gradient is often expensive.

\textbf{Stochastic gradient descent (SGD).}\\ \[
\theta_{t+1} = \theta_t - \eta \nabla \ell(f_\theta(x_i), y_i),
\]
where $(x_i, y_i)$ is a randomly selected data point.

\textbf{Mini-batch SGD.}\\ \[
\theta_{t+1} = \theta_t - \eta \frac{1}{|B|} \sum_{i \in B} \nabla \ell_i(\theta),
\]
where $B$ is a subset of the data.

\textbf{Properties.}\\ \begin{itemize}
    \item Lower computational cost per step
    \item Introduces noise into updates
\end{itemize}

\textbf{Effect of noise.}\\ \begin{itemize}
    \item Helps escape poor local minima
    \item Influences generalization behavior
\end{itemize}

\textbf{Tradeoff.}\\ \begin{itemize}
    \item Small batches: noisy but fast updates
    \item Large batches: stable but expensive
\end{itemize}

\textbf{Insight.}\\ Stochasticity is not merely a computational trick; it plays a role in learning dynamics.
\end{frame}

\subsection{Geometry of Loss Landscapes}
\begin{frame}[t,allowframebreaks]{Geometry of Loss Landscapes}
\small
The objective function $L(\theta)$ defines a landscape over parameter space.

\textbf{Parameter space.}\\ \[
\theta \in \mathbb{R}^d
\]

\textbf{Loss landscape.}\\ \[
L : \mathbb{R}^d \to \mathbb{R}
\]

\textbf{Critical points.}\\ \begin{itemize}
    \item Local minima: $\nabla L(\theta) = 0$, positive curvature
    \item Saddle points: mixed curvature directions
\end{itemize}

\textbf{Convex vs non-convex.}\\ \begin{itemize}
    \item Convex problems: global minimum is unique
    \item Non-convex problems: many local minima and saddle points
\end{itemize}

\textbf{High-dimensional effects.}\\ \begin{itemize}
    \item Saddle points are more common than poor local minima
    \item Many directions may be nearly flat
\end{itemize}

\textbf{Flat vs sharp minima.}\\ \begin{itemize}
    \item Flat minima: stable under perturbations
    \item Sharp minima: sensitive to parameter changes
\end{itemize}

\textbf{Geometric intuition.}\\ \begin{itemize}
    \item Optimization follows paths along the landscape
    \item Curvature affects convergence speed
\end{itemize}

\textbf{Insight.}\\ The geometry of the loss landscape determines both optimization behavior and generalization properties.
\end{frame}

\subsection{A Unifying Perspective}
\begin{frame}[t,allowframebreaks]{A Unifying Perspective}
\small
Learning as optimization brings together multiple aspects:

\begin{itemize}
    \item \textbf{Objective:}\\ encodes data fit and regularization
    \item \textbf{Algorithm:}\\ determines how solutions are found
    \item \textbf{Geometry:}\\ shapes optimization trajectories
    \item \textbf{Stochasticity:}\\ influences dynamics and generalization
\end{itemize}

\textbf{Core idea.}\\ Learning is a dynamical process in parameter space driven by gradients.

\textbf{Key insight.}\\ The outcome of learning depends not only on the objective, but also on the optimization algorithm used.
\end{frame}

\subsection{Outlook}
\begin{frame}[t,allowframebreaks]{Outlook}
\small
In this lecture, we formulated learning as an optimization problem and studied:

\begin{itemize}
    \item objective functions for learning,
    \item gradient-based optimization,
    \item stochastic methods,
    \item geometric properties of loss landscapes.
\end{itemize}

These ideas lead to deeper questions:

\begin{itemize}
    \item Why do some models generalize better than others?
    \item How does optimization interact with model complexity?
    \item What role does the structure of the loss landscape play?
\end{itemize}

In the next lecture, we study generalization and overfitting, focusing on how models perform on unseen data.

\textbf{Guiding principle.}\\ Learning is not only about finding a solution, but about finding a solution that generalizes.
\end{frame}

\section{Tutorial 2: Gradient Descent Exercises}
\SectionDivider{Tutorial 2: Gradient Descent Exercises}
\begin{frame}[t,allowframebreaks]{Outline}
\small
\begin{itemize}
  \item Objectives
  \item Exercise 1: Gradient of a Quadratic Loss
  \item Exercise 2: One-Dimensional Gradient Descent
  \item Exercise 3: Effect of Conditioning
  \item Exercise 4: Stochastic Gradient Descent
  \item Exercise 5: Gradient Descent as Projection
  \item Exercise 6 (Discussion): Why Does GD Work in High Dimensions?
  \item Summary
\end{itemize}
\end{frame}

\subsection{Objectives}
\begin{frame}[t,allowframebreaks]{Objectives}
\small
This tutorial develops intuition and technical skill for gradient-based optimization:
\begin{itemize}
    \item Compute gradients for common losses
    \item Analyze convergence of gradient descent
    \item Understand the role of step size and conditioning
    \item Connect algebraic updates with geometric intuition
\end{itemize}

\textbf{Focus.}\\ Work through simple but revealing cases where everything can be computed explicitly.
\end{frame}

\subsection{Exercise 1: Gradient of a Quadratic Loss}
\begin{frame}[t,allowframebreaks]{Exercise 1: Gradient of a Quadratic Loss}
\small
Let
\[
L(w) = \frac{1}{2}\|Xw - y\|^2, \quad X \in \mathbb{R}^{n \times d}, \; y \in \mathbb{R}^n.
\]

\textbf{Tasks.}\\ \begin{enumerate}
    \item Compute $\nabla_w L(w)$.
    \item Write the gradient descent update.
\end{enumerate}

\textbf{Solution.}\\ \begin{enumerate}
\item Using matrix calculus:
\[
\nabla_w L(w) = X^\top (Xw - y).
\]

\item Gradient descent:
\[
w_{t+1} = w_t - \eta X^\top (Xw_t - y).
\]
\end{enumerate}

\textbf{Key takeaway.}\\ Quadratic objectives lead to linear gradient updates.
\end{frame}

\subsection{Exercise 2: One-Dimensional Gradient Descent}
\begin{frame}[t,allowframebreaks]{Exercise 2: One-Dimensional Gradient Descent}
\small
Consider
\[
L(\theta) = \frac{1}{2} a \theta^2, \quad a > 0.
\]

\textbf{Tasks.}\\ \begin{enumerate}
    \item Compute the gradient.
    \item Write the update rule.
    \item Solve the recursion for $\theta_t$.
    \item Determine when the method converges.
\end{enumerate}

\textbf{Solution.}\\ \begin{enumerate}
\item 
\[
\nabla L(\theta) = a\theta.
\]

\item 
\[
\theta_{t+1} = \theta_t - \eta a \theta_t = (1 - \eta a)\theta_t.
\]

\item 
\[
\theta_t = (1 - \eta a)^t \theta_0.
\]

\item Convergence requires:
\[
|1 - \eta a| < 1 \quad \Rightarrow \quad 0 < \eta < \frac{2}{a}.
\]
\end{enumerate}

\textbf{Key takeaway.}\\ Step size controls convergence speed and stability.
\end{frame}

\subsection{Exercise 3: Effect of Conditioning}
\begin{frame}[t,allowframebreaks]{Exercise 3: Effect of Conditioning}
\small
Let
\[
L(w) = \frac{1}{2} w^\top A w,
\]
where $A$ is symmetric positive definite.

\textbf{Tasks.}\\ \begin{enumerate}
    \item Compute the gradient.
    \item Explain how eigenvalues of $A$ affect convergence.
\end{enumerate}

\textbf{Solution.}\\ \begin{enumerate}
\item 
\[
\nabla L(w) = A w.
\]

\item Let eigenvalues of $A$ be $\lambda_1 \leq \cdots \leq \lambda_d$.

\begin{itemize}
    \item Convergence rate depends on condition number:
    \[
    \kappa = \frac{\lambda_{\max}}{\lambda_{\min}}.
    \]
    \item Large $\kappa$ $\Rightarrow$ slow convergence
\end{itemize}
\end{enumerate}

\textbf{Geometric interpretation.}\\ Level sets are ellipses; gradient descent zig-zags along narrow directions.

\textbf{Key takeaway.}\\ Poor conditioning slows optimization.
\end{frame}

\subsection{Exercise 4: Stochastic Gradient Descent}
\begin{frame}[t,allowframebreaks]{Exercise 4: Stochastic Gradient Descent}
\small
Consider empirical loss:
\[
L(\theta) = \frac{1}{n} \sum_{i=1}^n \ell_i(\theta).
\]

\textbf{Tasks.}\\ \begin{enumerate}
    \item Define the SGD update.
    \item Compare SGD with full gradient descent.
    \item Explain the effect of stochastic noise.
\end{enumerate}

\textbf{Solution.}\\ \begin{enumerate}
\item SGD:
\[
\theta_{t+1} = \theta_t - \eta \nabla \ell_i(\theta_t),
\]
where $i$ is sampled randomly.

\item Comparison:
\begin{itemize}
    \item GD: exact gradient, expensive
    \item SGD: noisy gradient, cheap
\end{itemize}

\item Noise effects:
\begin{itemize}
    \item Helps escape saddle points
    \item Introduces fluctuations near minima
\end{itemize}
\end{enumerate}

\textbf{Key takeaway.}\\ SGD trades accuracy for efficiency and exploration.
\end{frame}

\subsection{Exercise 5: Gradient Descent as Projection}
\begin{frame}[t,allowframebreaks]{Exercise 5: Gradient Descent as Projection}
\small
Consider linear regression:
\[
\min_w \|Xw - y\|^2.
\]

\textbf{Tasks.}\\ \begin{enumerate}
    \item Describe the solution geometrically.
    \item Explain how gradient descent approaches this solution.
\end{enumerate}

\textbf{Solution.}\\ \begin{enumerate}
\item The solution is the projection of $y$ onto the column space of $X$.

\item Gradient descent:
\begin{itemize}
    \item Iteratively reduces residual $Xw - y$
    \item Moves toward projection point
\end{itemize}
\end{enumerate}

\textbf{Key takeaway.}\\ Optimization can be interpreted as a geometric projection process.
\end{frame}

\subsection{Exercise 6 (Discussion): Why Does GD Work in High Dimensions?}
\begin{frame}[t,allowframebreaks]{Exercise 6 (Discussion): Why Does GD Work in High Dimensions?}
\small
\textbf{Question.}\\ Neural networks are highly non-convex. Why does gradient descent often work well?

\textbf{Discussion points.}\\ \begin{itemize}
    \item Many local minima are similarly good
    \item High-dimensional landscapes dominated by saddle points
    \item Overparameterization simplifies optimization
\end{itemize}

\textbf{Insight.}\\ Optimization in modern machine learning behaves differently from classical low-dimensional intuition.
\end{frame}

\subsection{Summary}
\begin{frame}[t,allowframebreaks]{Summary}
\small
\begin{itemize}
    \item Gradient descent follows local gradients to minimize loss
    \item Step size controls convergence behavior
    \item Conditioning affects speed of convergence
    \item Stochasticity introduces noise but aids exploration
    \item Optimization has a geometric interpretation
\end{itemize}

\textbf{Preparation for next lecture.}\\ We now study how optimization interacts with generalization and overfitting.
\end{frame}

\section{Lecture 5: Generalization and Overfitting}
\SectionDivider{Lecture 5: Generalization and Overfitting}
\begin{frame}[t,allowframebreaks]{Outline}
\small
\begin{itemize}
  \item Motivation
  \item Training and Test Performance
  \item Overfitting and Underfitting
  \item Model Capacity and Complexity
  \item Informal Generalization Principles
  \item A Unifying Perspective
  \item Outlook
\end{itemize}
\end{frame}

\subsection{Motivation}
\begin{frame}[t,allowframebreaks]{Motivation}
\small
In previous lectures, we formulated learning as:
\begin{itemize}
    \item approximation of an unknown function,
    \item estimation from random data,
    \item optimization of an objective function.
\end{itemize}

However, a central challenge remains:

\textit{Why should a model that fits the training data perform well on unseen data?}

This question lies at the heart of machine learning and is addressed by the concept of \emph{generalization}.

\textbf{Guiding idea.}\\ Learning is successful only if performance extends beyond the training data.
\end{frame}

\subsection{Training and Test Performance}
\begin{frame}[t,allowframebreaks]{Training and Test Performance}
\small
Let $\mathcal{D}_{\text{train}}$ denote the training dataset and $\mathcal{D}_{\text{test}}$ denote unseen data drawn from the same distribution.

\textbf{Training error.}\\ \[
\hat{R}_{\text{train}}(f) = \frac{1}{n} \sum_{(x_i,y_i)\in \mathcal{D}_{\text{train}}} \ell(f(x_i), y_i).
\]

\textbf{Test error.}\\ \[
\hat{R}_{\text{test}}(f) = \frac{1}{m} \sum_{(x_j,y_j)\in \mathcal{D}_{\text{test}}} \ell(f(x_j), y_j).
\]

\textbf{True risk.}\\ \[
R(f) = \mathbb{E}_{(x,y)\sim P}[\ell(f(x), y)].
\]

\textbf{Generalization gap.}\\ \[
\hat{R}_{\text{test}}(f) - \hat{R}_{\text{train}}(f).
\]

\textbf{Interpretation.}\\ \begin{itemize}
    \item Training error measures fit to observed data
    \item Test error measures predictive performance
    \item The gap reflects overfitting
\end{itemize}

\textbf{Insight.}\\ Low training error alone is not sufficient; generalization is the true objective.
\end{frame}

\subsection{Overfitting and Underfitting}
\begin{frame}[t,allowframebreaks]{Overfitting and Underfitting}
\small
\textbf{Overfitting.}\\ \begin{itemize}
    \item Model fits training data extremely well
    \item Captures noise rather than underlying structure
    \item High variance, poor test performance
\end{itemize}

\textbf{Underfitting.}\\ \begin{itemize}
    \item Model is too simple
    \item Fails to capture underlying patterns
    \item High bias, poor performance on both training and test data
\end{itemize}

\textbf{Typical behavior.}\\ \begin{itemize}
    \item Increasing model complexity decreases training error
    \item Test error decreases initially, then increases (classical view)
\end{itemize}

\textbf{Illustration.}\\ \begin{itemize}
    \item Simple model: smooth but inaccurate
    \item Complex model: oscillatory, fits noise
\end{itemize}

\textbf{Insight.}\\ There exists an optimal level of complexity balancing bias and variance.
\end{frame}

\subsection{Model Capacity and Complexity}
\begin{frame}[t,allowframebreaks]{Model Capacity and Complexity}
\small
\textbf{Definition.}\\ Model capacity refers to the ability of a hypothesis space to fit a wide range of functions.

\textbf{Examples.}\\ \begin{itemize}
    \item Linear models: low capacity
    \item High-degree polynomials: high capacity
    \item Deep neural networks: extremely high capacity
\end{itemize}

\textbf{Measures of capacity.}\\ \begin{itemize}
    \item Number of parameters
    \item Flexibility of function class
    \item Ability to fit arbitrary labels
\end{itemize}

\textbf{Interpolation phenomenon.}\\ High-capacity models can fit training data exactly:
\[
\hat{R}_{\text{train}}(f) = 0.
\]

\textbf{Classical view.}\\ \begin{itemize}
    \item High capacity $\Rightarrow$ overfitting
\end{itemize}

\textbf{Modern observation.}\\ \begin{itemize}
    \item Overparameterized models can still generalize well
\end{itemize}

\textbf{Insight.}\\ Capacity alone does not determine generalization; the interaction with optimization is crucial.
\end{frame}

\subsection{Informal Generalization Principles}
\begin{frame}[t,allowframebreaks]{Informal Generalization Principles}
\small
Although a complete theory of generalization is complex, several guiding principles are useful.

\textbf{1. Simplicity (Occam's principle).}\\ \begin{itemize}
    \item Simpler models tend to generalize better
    \item Regularization promotes simplicity
\end{itemize}

\textbf{2. Stability.}\\ \begin{itemize}
    \item Models that are insensitive to small data changes generalize better
\end{itemize}

\textbf{3. Margin and robustness.}\\ \begin{itemize}
    \item Larger margins in classification often improve generalization
\end{itemize}

\textbf{4. Implicit bias of optimization.}\\ \begin{itemize}
    \item Optimization algorithms favor certain solutions
    \item Gradient descent often finds structured or “simple” solutions
\end{itemize}

\textbf{5. Data quality and representation.}\\ \begin{itemize}
    \item Better features lead to better generalization
\end{itemize}

\textbf{Insight.}\\ Generalization depends on a combination of model structure, data, and optimization dynamics.
\end{frame}

\subsection{A Unifying Perspective}
\begin{frame}[t,allowframebreaks]{A Unifying Perspective}
\small
Generalization can be understood as the relationship between:

\begin{itemize}
    \item \textbf{Data:}\\ finite samples from a distribution
    \item \textbf{Model:}\\ hypothesis space and capacity
    \item \textbf{Optimization:}\\ how solutions are selected
\end{itemize}

\textbf{Core idea.}\\ Learning selects one solution among many that fit the data; generalization depends on which solution is chosen.

\textbf{Key insight.}\\ The generalization behavior of modern machine learning systems cannot be explained by model capacity alone.
\end{frame}

\subsection{Outlook}
\begin{frame}[t,allowframebreaks]{Outlook}
\small
In this lecture, we studied:

\begin{itemize}
    \item training and test performance,
    \item overfitting and underfitting,
    \item model capacity and complexity,
    \item informal principles of generalization.
\end{itemize}

These ideas motivate the need for methods that control complexity and improve generalization.

In the next lecture, we study \emph{regularization}, which introduces explicit constraints and penalties to guide learning.

\textbf{Guiding principle.}\\ The goal of learning is not to fit the data perfectly, but to capture the underlying structure that generates it.
\end{frame}

\section{Lecture 6: Regularization and Structure}
\SectionDivider{Lecture 6: Regularization and Structure}
\begin{frame}[t,allowframebreaks]{Outline}
\small
\begin{itemize}
  \item Motivation
  \item Norm-Based Regularization
  \item Constraints and Penalization
  \item Implicit Regularization
  \item Early Stopping and Practical Methods
  \item A Unifying Perspective
  \item Outlook
\end{itemize}
\end{frame}

\subsection{Motivation}
\begin{frame}[t,allowframebreaks]{Motivation}
\small
In the previous lecture, we studied generalization and observed that:

\begin{itemize}
    \item Models with high capacity can overfit
    \item Low training error does not guarantee low test error
\end{itemize}

This raises a central question:

\textit{How can we guide learning toward solutions that generalize well?}

The answer lies in \emph{regularization}, which introduces additional structure or constraints into the learning process.

\textbf{Guiding idea.}\\ Regularization shapes the solution space to favor functions with desirable properties.
\end{frame}

\subsection{Norm-Based Regularization}
\begin{frame}[t,allowframebreaks]{Norm-Based Regularization}
\small
A common approach is to penalize the size of the parameters.

\textbf{Regularized objective.}\\ \[
\min_\theta \; L(\theta) + \lambda \|\theta\|_p,
\]
where $\lambda > 0$ controls the strength of regularization.

\textbf{Examples.}\\ \begin{itemize}
    \item \textbf{$\ell_2$ regularization (ridge):}\\ \[
    \Omega(\theta) = \|\theta\|_2^2
    \]

    \item \textbf{$\ell_1$ regularization (lasso):}\\ \[
    \Omega(\theta) = \|\theta\|_1
    \]
\end{itemize}

\textbf{Effects.}\\ \begin{itemize}
    \item $\ell_2$: discourages large weights, promotes smooth solutions
    \item $\ell_1$: encourages sparsity
\end{itemize}

\textbf{Interpretation.}\\ \begin{itemize}
    \item Penalizing large parameters reduces model complexity
    \item Encourages simpler functions
\end{itemize}

\textbf{Insight.}\\ Norm-based regularization controls capacity by constraining parameter magnitude.
\end{frame}

\subsection{Constraints and Penalization}
\begin{frame}[t,allowframebreaks]{Constraints and Penalization}
\small
Regularization can also be formulated as a constrained optimization problem.

\textbf{Constrained formulation.}\\ \[
\min_\theta \; L(\theta) \quad \text{subject to} \quad \|\theta\|_p \leq C.
\]

\textbf{Penalized formulation.}\\ \[
\min_\theta \; L(\theta) + \lambda \|\theta\|_p.
\]

\textbf{Equivalence.}\\ \begin{itemize}
    \item Under suitable conditions, these formulations are equivalent
    \item $\lambda$ corresponds to a constraint level $C$
\end{itemize}

\textbf{Geometric interpretation.}\\ \begin{itemize}
    \item Optimization occurs within a constraint set (e.g., a ball)
    \item The solution lies where level sets of $L(\theta)$ touch the constraint boundary
\end{itemize}

\textbf{Example.}\\ \begin{itemize}
    \item $\ell_2$ constraint: spherical region
    \item $\ell_1$ constraint: diamond-shaped region (promotes sparsity)
\end{itemize}

\textbf{Insight.}\\ Regularization can be viewed as restricting the feasible set of solutions.
\end{frame}

\subsection{Implicit Regularization}
\begin{frame}[t,allowframebreaks]{Implicit Regularization}
\small
Regularization is not only imposed explicitly; it can also arise implicitly through the learning process.

\textbf{Definition.}\\ Implicit regularization refers to biases introduced by the optimization algorithm or parameterization, even without explicit penalties.

\textbf{Examples.}\\ \begin{itemize}
    \item Gradient descent in linear models tends to find minimum-norm solutions
    \item Optimization dynamics favor certain regions of parameter space
\end{itemize}

\textbf{Overparameterized models.}\\ \begin{itemize}
    \item Many solutions achieve zero training error
    \item Optimization selects one among them
\end{itemize}

\textbf{Key phenomenon.}\\ \begin{itemize}
    \item Even without explicit regularization, learned solutions can generalize well
\end{itemize}

\textbf{Insight.}\\ The optimization algorithm itself acts as a form of inductive bias.
\end{frame}

\subsection{Early Stopping and Practical Methods}
\begin{frame}[t,allowframebreaks]{Early Stopping and Practical Methods}
\small
Regularization can also be implemented through practical training strategies.

\textbf{Early stopping.}\\ \begin{itemize}
    \item Stop training before convergence
    \item Prevents overfitting to noise
\end{itemize}

\textbf{Interpretation.}\\ \begin{itemize}
    \item Early iterations capture dominant patterns
    \item Later iterations fit noise
\end{itemize}

\textbf{Other methods.}\\ \begin{itemize}
    \item Noise injection
    \item Data augmentation
    \item Dropout (in neural networks)
\end{itemize}

\textbf{Connection to optimization.}\\ \begin{itemize}
    \item Training trajectory influences final solution
    \item Regularization is tied to dynamics, not just objectives
\end{itemize}

\textbf{Insight.}\\ Regularization can be achieved through both objective design and training procedures.
\end{frame}

\subsection{A Unifying Perspective}
\begin{frame}[t,allowframebreaks]{A Unifying Perspective}
\small
Regularization integrates several aspects of learning:

\begin{itemize}
    \item \textbf{Explicit:}\\ penalties and constraints
    \item \textbf{Implicit:}\\ optimization dynamics
    \item \textbf{Practical:}\\ training strategies
\end{itemize}

\textbf{Core idea.}\\ Regularization shapes the set of solutions considered by the learning algorithm.

\textbf{Key insight.}\\ Generalization depends not only on model capacity, but on how solutions are selected within that capacity.
\end{frame}

\subsection{Outlook}
\begin{frame}[t,allowframebreaks]{Outlook}
\small
In this lecture, we studied:

\begin{itemize}
    \item norm-based regularization,
    \item constrained and penalized formulations,
    \item implicit regularization,
    \item practical regularization methods.
\end{itemize}

These ideas provide tools for controlling generalization in learning systems.

In the next part of the course, we turn to concrete models, beginning with linear regression, where these concepts can be analyzed explicitly.

\textbf{Guiding principle.}\\ Effective learning requires not only fitting data, but shaping the solution toward structure and simplicity.
\end{frame}

\section{Tutorial 3: Regularization and Overfitting}
\SectionDivider{Tutorial 3: Regularization and Overfitting}
\begin{frame}[t,allowframebreaks]{Outline}
\small
\begin{itemize}
  \item Objectives
  \item Exercise 1: Overfitting in Polynomial Regression
  \item Exercise 2: Ridge Regression Effect
  \item Exercise 3: Geometry of Regularization
  \item Exercise 4: Implicit Regularization
  \item Exercise 5: Early Stopping
  \item Exercise 6 (Discussion): Double Descent (Preview)
  \item Summary
\end{itemize}
\end{frame}

\subsection{Objectives}
\begin{frame}[t,allowframebreaks]{Objectives}
\small
This tutorial develops intuition for generalization and regularization:
\begin{itemize}
    \item Understand overfitting and underfitting through examples
    \item Analyze the effect of model complexity
    \item Study the role of regularization
    \item Interpret learning dynamics (early stopping, implicit bias)
\end{itemize}

\textbf{Focus.}\\ Bridge theory and practice through concrete scenarios and reasoning.
\end{frame}

\subsection{Exercise 1: Overfitting in Polynomial Regression}
\begin{frame}[t,allowframebreaks]{Exercise 1: Overfitting in Polynomial Regression}
\small
Suppose we fit a polynomial of degree $d$ to noisy data:
\[
y = f(x) + \varepsilon.
\]

\textbf{Tasks.}\\ \begin{enumerate}
    \item Describe what happens when $d$ is very small.
    \item Describe what happens when $d$ is very large.
    \item Sketch (qualitatively) training and test error as a function of $d$.
\end{enumerate}

\textbf{Solution.}\\ \begin{enumerate}
\item \textbf{Small $d$:}\\ \begin{itemize}
    \item Model too simple
    \item Cannot capture structure
    \item High bias
\end{itemize}

\item \textbf{Large $d$:}\\ \begin{itemize}
    \item Model highly flexible
    \item Fits noise
    \item High variance
\end{itemize}

\item \textbf{Behavior:}\\ \begin{itemize}
    \item Training error decreases monotonically
    \item Test error decreases then increases (U-shape)
\end{itemize}
\end{enumerate}

\textbf{Key takeaway.}\\ Model complexity must be matched to data.
\end{frame}

\subsection{Exercise 2: Ridge Regression Effect}
\begin{frame}[t,allowframebreaks]{Exercise 2: Ridge Regression Effect}
\small
Consider ridge regression:
\[
\min_w \|Xw - y\|^2 + \lambda \|w\|^2.
\]

\textbf{Tasks.}\\ \begin{enumerate}
    \item What happens when $\lambda = 0$?
    \item What happens when $\lambda \to \infty$?
    \item Explain how $\lambda$ affects bias and variance.
\end{enumerate}

\textbf{Solution.}\\ \begin{enumerate}
\item $\lambda = 0$:
\begin{itemize}
    \item Standard least squares
    \item Potential overfitting
\end{itemize}

\item $\lambda \to \infty$:
\begin{itemize}
    \item $w \to 0$
    \item Model predicts near-constant output
    \item High bias
\end{itemize}

\item Effect:
\begin{itemize}
    \item Increasing $\lambda$ increases bias
    \item Increasing $\lambda$ reduces variance
\end{itemize}
\end{enumerate}

\textbf{Key takeaway.}\\ Regularization controls the bias–variance tradeoff.
\end{frame}

\subsection{Exercise 3: Geometry of Regularization}
\begin{frame}[t,allowframebreaks]{Exercise 3: Geometry of Regularization}
\small
Consider:
\[
\min_w \|Xw - y\|^2 \quad \text{subject to} \quad \|w\|_2 \leq C.
\]

\textbf{Tasks.}\\ \begin{enumerate}
    \item Describe the feasible set.
    \item Explain where the solution lies geometrically.
\end{enumerate}

\textbf{Solution.}\\ \begin{enumerate}
\item The feasible set is a ball in $\mathbb{R}^d$.

\item The solution lies where:
\begin{itemize}
    \item level sets of the loss intersect the constraint boundary
\end{itemize}
\end{enumerate}

\textbf{Insight.}\\ Regularization restricts the search space to structured regions.
\end{frame}

\subsection{Exercise 4: Implicit Regularization}
\begin{frame}[t,allowframebreaks]{Exercise 4: Implicit Regularization}
\small
Consider overparameterized linear regression:
\[
\min_w \|Xw - y\|^2,
\]
where infinitely many solutions exist.

\textbf{Tasks.}\\ \begin{enumerate}
    \item Explain why multiple solutions exist.
    \item Which solution does gradient descent find (starting from $w=0$)?
\end{enumerate}

\textbf{Solution.}\\ \begin{enumerate}
\item If $X$ has more parameters than constraints:
\[
Xw = y
\]
has infinitely many solutions.

\item Gradient descent finds the minimum-norm solution:
\[
w^* = \arg\min \|w\| \quad \text{subject to } Xw = y.
\]
\end{enumerate}

\textbf{Key takeaway.}\\ Optimization induces implicit regularization.
\end{frame}

\subsection{Exercise 5: Early Stopping}
\begin{frame}[t,allowframebreaks]{Exercise 5: Early Stopping}
\small
Suppose we train a model using gradient descent.

\textbf{Tasks.}\\ \begin{enumerate}
    \item What happens if we train for too long?
    \item Why can stopping early improve generalization?
\end{enumerate}

\textbf{Solution.}\\ \begin{enumerate}
\item Training too long:
\begin{itemize}
    \item Model may fit noise
    \item Overfitting occurs
\end{itemize}

\item Early stopping:
\begin{itemize}
    \item Limits effective complexity
    \item Prevents fitting fine-scale noise
\end{itemize}
\end{enumerate}

\textbf{Insight.}\\ Training time acts as a form of regularization.
\end{frame}

\subsection{Exercise 6 (Discussion): Double Descent (Preview)}
\begin{frame}[t,allowframebreaks]{Exercise 6 (Discussion): Double Descent (Preview)}
\small
\textbf{Question.}\\ Classically, test error is U-shaped. Modern deep learning sometimes shows a second descent after interpolation.

\textbf{Discussion points.}\\ \begin{itemize}
    \item Overparameterization
    \item Implicit bias of optimization
    \item Role of noise and structure
\end{itemize}

\textbf{Insight.}\\ Generalization behavior in modern models differs from classical intuition.
\end{frame}

\subsection{Summary}
\begin{frame}[t,allowframebreaks]{Summary}
\small
\begin{itemize}
    \item Overfitting arises from excessive model flexibility
    \item Regularization controls complexity explicitly and implicitly
    \item Geometry provides insight into constrained optimization
    \item Optimization dynamics influence generalization
\end{itemize}

\textbf{Preparation for next lectures.}\\ We now turn to linear regression, where these ideas can be analyzed exactly and explicitly.
\end{frame}

\section{Lecture 7: Linear Regression and Geometry}
\SectionDivider{Lecture 7: Linear Regression and Geometry}
\begin{frame}[t,allowframebreaks]{Outline}
\small
\begin{itemize}
  \item Motivation
  \item Linear Models and Projections
  \item Least Squares as Optimization
  \item Conditioning and Numerical Stability
  \item Ridge Regression
  \item A Unifying Perspective
  \item Outlook
\end{itemize}
\end{frame}

\subsection{Motivation}
\begin{frame}[t,allowframebreaks]{Motivation}
\small
Linear regression is one of the simplest yet most fundamental models in machine learning. Despite its apparent simplicity, it illustrates many core ideas:

\begin{itemize}
    \item function approximation,
    \item optimization,
    \item geometry of solutions,
    \item regularization and generalization.
\end{itemize}

\textbf{Guiding idea.}\\ Linear regression provides a setting where learning can be understood exactly, both algebraically and geometrically.
\end{frame}

\subsection{Linear Models and Projections}
\begin{frame}[t,allowframebreaks]{Linear Models and Projections}
\small
We consider a linear model:
\[
f_w(x) = w^\top x, \quad w \in \mathbb{R}^d.
\]

Given data $(x_i, y_i)$, define the design matrix:
\[
X =
\begin{bmatrix}
x_1^\top \\
x_2^\top \\
\vdots \\
x_n^\top
\end{bmatrix}
\in \mathbb{R}^{n \times d}, \quad
y =
\begin{bmatrix}
y_1 \\
y_2 \\
\vdots \\
y_n
\end{bmatrix}.
\]

Predictions are given by:
\[
\hat{y} = Xw.
\]

\textbf{Geometric interpretation.}\\ \begin{itemize}
    \item The column space of $X$ is a subspace of $\mathbb{R}^n$
    \item Predictions $Xw$ lie in this subspace
\end{itemize}

\textbf{Projection viewpoint.}\\ The optimal prediction is the projection of $y$ onto the column space of $X$.

\textbf{Orthogonality condition.}\\ \[
X^\top (Xw - y) = 0.
\]

\textbf{Insight.}\\ Learning corresponds to projecting observed outputs onto the space of representable predictions.
\end{frame}

\subsection{Least Squares as Optimization}
\begin{frame}[t,allowframebreaks]{Least Squares as Optimization}
\small
We minimize the squared loss:
\[
L(w) = \|Xw - y\|^2.
\]

\textbf{Normal equations.}\\ \[
X^\top X w = X^\top y.
\]

\textbf{Solution (full rank case).}\\ \[
w^* = (X^\top X)^{-1} X^\top y.
\]

\textbf{Interpretation.}\\ \begin{itemize}
    \item Closed-form solution exists
    \item Equivalent to minimizing empirical risk with squared loss
\end{itemize}

\textbf{Overparameterized case ($d > n$).}\\ \begin{itemize}
    \item Infinitely many solutions
    \item Minimum-norm solution:
    \[
    w^* = X^\top (XX^\top)^{-1} y
    \]
\end{itemize}

\textbf{Connection to optimization.}\\ Gradient descent converges to this minimum-norm solution when initialized at $w=0$.

\textbf{Insight.}\\ Even simple models reveal the role of implicit regularization.
\end{frame}

\subsection{Conditioning and Numerical Stability}
\begin{frame}[t,allowframebreaks]{Conditioning and Numerical Stability}
\small
The matrix $X^\top X$ plays a central role.

\textbf{Condition number.}\\ \[
\kappa = \frac{\lambda_{\max}}{\lambda_{\min}},
\]
where $\lambda_i$ are eigenvalues of $X^\top X$.

\textbf{Ill-conditioning.}\\ \begin{itemize}
    \item Small eigenvalues $\Rightarrow$ unstable inversion
    \item Solutions sensitive to noise
\end{itemize}

\textbf{Effects.}\\ \begin{itemize}
    \item Slow convergence of gradient descent
    \item Large variance in estimated parameters
\end{itemize}

\textbf{Geometric view.}\\ \begin{itemize}
    \item Level sets are elongated ellipsoids
    \item Optimization proceeds slowly along narrow directions
\end{itemize}

\textbf{Insight.}\\ Numerical properties of the data matrix strongly affect learning behavior.
\end{frame}

\subsection{Ridge Regression}
\begin{frame}[t,allowframebreaks]{Ridge Regression}
\small
To address instability and overfitting, we introduce $\ell_2$ regularization:
\[
\min_w \|Xw - y\|^2 + \lambda \|w\|^2.
\]

\textbf{Solution.}\\ \[
w^* = (X^\top X + \lambda I)^{-1} X^\top y.
\]

\textbf{Interpretation.}\\ \begin{itemize}
    \item Adds stability to inversion
    \item Shrinks parameter values
\end{itemize}

\textbf{Geometric view.}\\ \begin{itemize}
    \item Solution lies in intersection of:
    \begin{itemize}
        \item loss level sets
        \item $\ell_2$ ball
    \end{itemize}
\end{itemize}

\textbf{Effect on eigenvalues.}\\ \[
\lambda_i \mapsto \lambda_i + \lambda,
\]
which improves conditioning.

\textbf{Bias--variance tradeoff.}\\ \begin{itemize}
    \item Increased bias
    \item Reduced variance
\end{itemize}

\textbf{Insight.}\\ Ridge regression stabilizes learning and improves generalization.
\end{frame}

\subsection{A Unifying Perspective}
\begin{frame}[t,allowframebreaks]{A Unifying Perspective}
\small
Linear regression illustrates key principles of learning:

\begin{itemize}
    \item \textbf{Approximation:}\\ linear hypothesis space
    \item \textbf{Optimization:}\\ least squares objective
    \item \textbf{Geometry:}\\ projection onto subspaces
    \item \textbf{Regularization:}\\ explicit and implicit
\end{itemize}

\textbf{Core idea.}\\ Learning is the projection of data onto a structured space, shaped by both the model and the optimization process.

\textbf{Key insight.}\\ Even the simplest models reveal deep connections between geometry, optimization, and generalization.
\end{frame}

\subsection{Outlook}
\begin{frame}[t,allowframebreaks]{Outlook}
\small
In this lecture, we developed a complete understanding of linear regression, including:

\begin{itemize}
    \item geometric interpretation via projections,
    \item closed-form solutions,
    \item numerical stability and conditioning,
    \item regularization through ridge regression.
\end{itemize}

These ideas extend naturally to classification problems.

In the next lecture, we study logistic regression, where probabilistic modeling and optimization come together in classification tasks.

\textbf{Guiding principle.}\\ Simple models provide a foundation for understanding complex learning systems.
\end{frame}

\section{Lecture 8: Logistic Regression and Classification}
\SectionDivider{Lecture 8: Logistic Regression and Classification}
\begin{frame}[t,allowframebreaks]{Outline}
\small
\begin{itemize}
  \item Motivation
  \item Probabilistic Interpretation
  \item Cross-Entropy Loss
  \item Maximum Likelihood Perspective
  \item Decision Boundaries
  \item A Unifying Perspective
  \item Outlook
\end{itemize}
\end{frame}

\subsection{Motivation}
\begin{frame}[t,allowframebreaks]{Motivation}
\small
In the previous lecture, we studied linear regression, where outputs are continuous. We now turn to \emph{classification}, where outputs are discrete.

\textbf{Goal.}\\ Given input $x \in \mathcal{X}$, predict a label:
\[
y \in \{0,1\}.
\]

A naive approach would be to apply linear regression and threshold the output. However, this lacks a probabilistic interpretation and can produce invalid predictions.

\textbf{Guiding idea.}\\ Model the probability of class membership directly.
\end{frame}

\subsection{Probabilistic Interpretation}
\begin{frame}[t,allowframebreaks]{Probabilistic Interpretation}
\small
We model the conditional probability:
\[
p(y = 1 \mid x) = \sigma(w^\top x),
\]
where $\sigma$ is the sigmoid function:
\[
\sigma(z) = \frac{1}{1 + e^{-z}}.
\]

\textbf{Properties.}\\ \begin{itemize}
    \item $\sigma(z) \in (0,1)$
    \item Monotonically increasing
    \item Maps linear scores to probabilities
\end{itemize}

\textbf{Odds interpretation.}\\ \[
\frac{p(y=1 \mid x)}{p(y=0 \mid x)} = e^{w^\top x}.
\]

Taking logarithms:
\[
\log \frac{p(y=1 \mid x)}{p(y=0 \mid x)} = w^\top x.
\]

\textbf{Insight.}\\ Logistic regression models the log-odds as a linear function of the input.
\end{frame}

\subsection{Cross-Entropy Loss}
\begin{frame}[t,allowframebreaks]{Cross-Entropy Loss}
\small
Given predictions $p_i = \sigma(w^\top x_i)$, we define the loss:
\[
\ell(w) = - \sum_{i=1}^n \left[
y_i \log p_i + (1 - y_i)\log(1 - p_i)
\right].
\]

\textbf{Interpretation.}\\ \begin{itemize}
    \item Penalizes confident incorrect predictions
    \item Encourages correct probability estimates
\end{itemize}

\textbf{Gradient.}\\ \[
\nabla_w \ell(w) = \sum_{i=1}^n (p_i - y_i)x_i.
\]

\textbf{Properties.}\\ \begin{itemize}
    \item Convex objective
    \item No closed-form solution
    \item Solved using iterative methods
\end{itemize}

\textbf{Insight.}\\ Cross-entropy aligns optimization with probabilistic modeling.
\end{frame}

\subsection{Maximum Likelihood Perspective}
\begin{frame}[t,allowframebreaks]{Maximum Likelihood Perspective}
\small
We assume:
\[
y_i \sim \text{Bernoulli}(p_i), \quad p_i = \sigma(w^\top x_i).
\]

\textbf{Likelihood.}\\ \[
p(y_i \mid x_i, w) = p_i^{y_i}(1 - p_i)^{1 - y_i}.
\]

\textbf{Log-likelihood.}\\ \[
\sum_{i=1}^n \left[
y_i \log p_i + (1 - y_i)\log(1 - p_i)
\right].
\]

\textbf{Result.}\\ Maximizing likelihood is equivalent to minimizing cross-entropy loss.

\textbf{Interpretation.}\\ \begin{itemize}
    \item Learning corresponds to estimating model parameters under a probabilistic model
\end{itemize}

\textbf{Insight.}\\ Logistic regression arises naturally from maximum likelihood estimation.
\end{frame}

\subsection{Decision Boundaries}
\begin{frame}[t,allowframebreaks]{Decision Boundaries}
\small
The predicted class is:
\[
\hat{y} =
\begin{cases}
1 & \text{if } p(y=1 \mid x) \geq 0.5 \\
0 & \text{otherwise}
\end{cases}
\]

Since $\sigma(z) \geq 0.5$ when $z \geq 0$, this becomes:
\[
w^\top x = 0.
\]

\textbf{Decision boundary.}\\ \begin{itemize}
    \item A hyperplane in $\mathbb{R}^d$
\end{itemize}

\textbf{Geometry.}\\ \begin{itemize}
    \item Points are classified based on which side of the hyperplane they lie
\end{itemize}

\textbf{Extension.}\\ \begin{itemize}
    \item Nonlinear boundaries can be obtained using feature maps $\phi(x)$
\end{itemize}

\textbf{Insight.}\\ Logistic regression produces linear decision boundaries in feature space.
\end{frame}

\subsection{A Unifying Perspective}
\begin{frame}[t,allowframebreaks]{A Unifying Perspective}
\small
Logistic regression integrates key ideas from previous lectures:

\begin{itemize}
    \item \textbf{Function approximation:}\\ linear model in feature space
    \item \textbf{Probability:}\\ modeling conditional distributions
    \item \textbf{Optimization:}\\ minimizing cross-entropy
    \item \textbf{Geometry:}\\ linear decision boundaries
\end{itemize}

\textbf{Core idea.}\\ Classification can be framed as probabilistic modeling combined with optimization.

\textbf{Key insight.}\\ Logistic regression is a bridge between linear models and more complex classification methods.
\end{frame}

\subsection{Outlook}
\begin{frame}[t,allowframebreaks]{Outlook}
\small
In this lecture, we studied:

\begin{itemize}
    \item probabilistic modeling of classification,
    \item cross-entropy loss,
    \item maximum likelihood estimation,
    \item geometric interpretation via decision boundaries.
\end{itemize}

However, linear models are limited in their expressivity.

In the next lecture, we introduce feature maps and kernel methods, which enable nonlinear modeling while retaining linear structure.

\textbf{Guiding principle.}\\ By combining representations with linear models, we can capture complex patterns in data.
\end{frame}

\section{Tutorial 4: Regression vs Classification}
\SectionDivider{Tutorial 4: Regression vs Classification}
\begin{frame}[t,allowframebreaks]{Outline}
\small
\begin{itemize}
  \item Objectives
  \item Exercise 1: Linear Regression for Classification
  \item Exercise 2: Logistic Regression Interpretation
  \item Exercise 3: Loss Function Comparison
  \item Exercise 4: Decision Boundaries
  \item Exercise 5: Effect of Feature Mapping
  \item Exercise 6 (Discussion): When to Use Regression vs Classification
  \item Summary
\end{itemize}
\end{frame}

\subsection{Objectives}
\begin{frame}[t,allowframebreaks]{Objectives}
\small
This tutorial develops a clear understanding of the differences and connections between regression and classification:

\begin{itemize}
    \item Compare linear regression and logistic regression
    \item Understand loss functions and probabilistic interpretations
    \item Analyze decision boundaries
    \item Explore when each formulation is appropriate
\end{itemize}

\textbf{Focus.}\\ Unify regression and classification under the framework of function approximation and probabilistic modeling.
\end{frame}

\subsection{Exercise 1: Linear Regression for Classification}
\begin{frame}[t,allowframebreaks]{Exercise 1: Linear Regression for Classification}
\small
Suppose we use linear regression for binary classification:
\[
f(x) = w^\top x, \quad \hat{y} = \text{sign}(f(x)).
\]

\textbf{Tasks.}\\ \begin{enumerate}
    \item What loss is being minimized?
    \item What are the potential issues with this approach?
\end{enumerate}

\textbf{Solution.}\\ \begin{enumerate}
\item The squared loss:
\[
\sum_i (w^\top x_i - y_i)^2.
\]

\item Issues:
\begin{itemize}
    \item Predictions are not probabilities
    \item Outputs can lie outside $[0,1]$
    \item Sensitive to outliers
    \item Does not align well with classification objectives
\end{itemize}
\end{enumerate}

\textbf{Key takeaway.}\\ Regression loss is not well-suited for classification.
\end{frame}

\subsection{Exercise 2: Logistic Regression Interpretation}
\begin{frame}[t,allowframebreaks]{Exercise 2: Logistic Regression Interpretation}
\small
Recall:
\[
p(y=1 \mid x) = \sigma(w^\top x).
\]

\textbf{Tasks.}\\ \begin{enumerate}
    \item Why is the sigmoid function appropriate here?
    \item What does $w^\top x$ represent?
\end{enumerate}

\textbf{Solution.}\\ \begin{enumerate}
\item Sigmoid maps $\mathbb{R} \to (0,1)$, making it suitable for probabilities.

\item $w^\top x$ represents the log-odds:
\[
\log \frac{p(y=1 \mid x)}{p(y=0 \mid x)}.
\]
\end{enumerate}

\textbf{Key takeaway.}\\ Logistic regression provides a principled probabilistic interpretation.
\end{frame}

\subsection{Exercise 3: Loss Function Comparison}
\begin{frame}[t,allowframebreaks]{Exercise 3: Loss Function Comparison}
\small
Compare:

\begin{itemize}
    \item Squared loss (regression)
    \item Cross-entropy loss (classification)
\end{itemize}

\textbf{Tasks.}\\ \begin{enumerate}
    \item How do they penalize errors differently?
    \item Which is more appropriate for classification and why?
\end{enumerate}

\textbf{Solution.}\\ \begin{enumerate}
\item 
\begin{itemize}
    \item Squared loss penalizes distance symmetrically
    \item Cross-entropy penalizes confident wrong predictions heavily
\end{itemize}

\item Cross-entropy is more appropriate because:
\begin{itemize}
    \item It arises from likelihood
    \item It reflects uncertainty
    \item It aligns with classification objectives
\end{itemize}
\end{enumerate}

\textbf{Key takeaway.}\\ Choice of loss reflects modeling assumptions.
\end{frame}

\subsection{Exercise 4: Decision Boundaries}
\begin{frame}[t,allowframebreaks]{Exercise 4: Decision Boundaries}
\small
Consider:
\[
f(x) = w^\top x.
\]

\textbf{Tasks.}\\ \begin{enumerate}
    \item Describe the decision boundary for regression (if thresholded).
    \item Describe the decision boundary for logistic regression.
\end{enumerate}

\textbf{Solution.}\\ \begin{enumerate}
\item Regression with thresholding:
\begin{itemize}
    \item Boundary: $w^\top x = c$ (depends on threshold)
\end{itemize}

\item Logistic regression:
\begin{itemize}
    \item Boundary: $w^\top x = 0$
\end{itemize}
\end{enumerate}

\textbf{Insight.}\\ Both produce linear boundaries, but logistic regression has a probabilistic interpretation.
\end{frame}

\subsection{Exercise 5: Effect of Feature Mapping}
\begin{frame}[t,allowframebreaks]{Exercise 5: Effect of Feature Mapping}
\small
Suppose we use features $\phi(x)$ instead of $x$.

\textbf{Tasks.}\\ \begin{enumerate}
    \item What happens to the decision boundary?
    \item How does this affect expressivity?
\end{enumerate}

\textbf{Solution.}\\ \begin{enumerate}
\item Decision boundary becomes:
\[
w^\top \phi(x) = 0,
\]
which can be nonlinear in the original input space.

\item Expressivity increases:
\begin{itemize}
    \item More complex patterns can be captured
\end{itemize}
\end{enumerate}

\textbf{Key takeaway.}\\ Nonlinearity arises through representation, not necessarily through the model.
\end{frame}

\subsection{Exercise 6 (Discussion): When to Use Regression vs Classification}
\begin{frame}[t,allowframebreaks]{Exercise 6 (Discussion): When to Use Regression vs Classification}
\small
\textbf{Question.}\\ When should we use regression, and when should we use classification?

\textbf{Discussion points.}\\ \begin{itemize}
    \item Regression:
    \begin{itemize}
        \item Continuous outputs
        \item Predicting quantities
    \end{itemize}

    \item Classification:
    \begin{itemize}
        \item Discrete labels
        \item Decision-making tasks
    \end{itemize}
\end{itemize}

\textbf{Subtle point.}\\ \begin{itemize}
    \item Classification models often estimate probabilities
    \item Final decisions arise from thresholds
\end{itemize}

\textbf{Insight.}\\ Classification can be viewed as structured regression on probabilities.
\end{frame}

\subsection{Summary}
\begin{frame}[t,allowframebreaks]{Summary}
\small
\begin{itemize}
    \item Regression predicts continuous values; classification predicts labels
    \item Logistic regression provides probabilistic classification
    \item Cross-entropy aligns with likelihood-based modeling
    \item Both models can produce linear decision boundaries
    \item Nonlinearity arises through feature transformations
\end{itemize}

\textbf{Preparation for next lecture.}\\ We now generalize linear models using feature maps and kernel methods to capture nonlinear relationships.
\end{frame}

\section{Lecture 9: Feature Maps and Nonlinear Models}
\SectionDivider{Lecture 9: Feature Maps and Nonlinear Models}
\begin{frame}[t,allowframebreaks]{Outline}
\small
\begin{itemize}
  \item Motivation
  \item Basis Expansions
  \item Linear Models in Feature Space
  \item The Kernel Idea (Intuition)
  \item Expressivity through Features
  \item A Unifying Perspective
  \item Outlook
\end{itemize}
\end{frame}

\subsection{Motivation}
\begin{frame}[t,allowframebreaks]{Motivation}
\small
In previous lectures, we studied linear models such as linear and logistic regression.

While these models are simple and interpretable, they are limited:

\begin{itemize}
    \item They can only represent linear relationships
    \item Many real-world patterns are inherently nonlinear
\end{itemize}

\textbf{Key question.}\\ How can we model nonlinear relationships while retaining the simplicity of linear models?

\textbf{Guiding idea.}\\ Transform the input representation so that nonlinear patterns become linear.
\end{frame}

\subsection{Basis Expansions}
\begin{frame}[t,allowframebreaks]{Basis Expansions}
\small
We introduce a feature map:
\[
\phi : \mathcal{X} \to \mathbb{R}^D,
\]
which transforms inputs into a higher-dimensional feature space.

\textbf{Examples.}\\ \begin{itemize}
    \item Polynomial features:
    \[
    \phi(x) = (1, x, x^2, \dots, x^d)
    \]

    \item Interaction terms:
    \[
    \phi(x_1, x_2) = (x_1, x_2, x_1 x_2)
    \]
\end{itemize}

\textbf{Model.}\\ \[
f(x) = w^\top \phi(x).
\]

\textbf{Interpretation.}\\ \begin{itemize}
    \item Model is linear in $\phi(x)$
    \item May be nonlinear in original input $x$
\end{itemize}

\textbf{Insight.}\\ Nonlinearity can be introduced through feature transformations.
\end{frame}

\subsection{Linear Models in Feature Space}
\begin{frame}[t,allowframebreaks]{Linear Models in Feature Space}
\small
We now apply linear regression or logistic regression to $\phi(x)$:

\[
f(x) = w^\top \phi(x).
\]

\textbf{Key observation.}\\ \begin{itemize}
    \item The learning algorithm remains unchanged
    \item Only the representation changes
\end{itemize}

\textbf{Decision boundaries.}\\ \begin{itemize}
    \item Linear in feature space
    \item Nonlinear in input space
\end{itemize}

\textbf{Example.}\\ \begin{itemize}
    \item A quadratic feature map produces curved decision boundaries
\end{itemize}

\textbf{Insight.}\\ Feature maps allow us to trade simplicity of models for richness of representation.
\end{frame}

\subsection{The Kernel Idea (Intuition)}
\begin{frame}[t,allowframebreaks]{The Kernel Idea (Intuition)}
\small
Working with high-dimensional feature maps can be computationally expensive.

\textbf{Key observation.}\\ Many learning algorithms depend only on inner products:
\[
\langle \phi(x), \phi(x') \rangle.
\]

\textbf{Kernel function.}\\ \[
k(x, x') = \langle \phi(x), \phi(x') \rangle.
\]

\textbf{Examples.}\\ \begin{itemize}
    \item Polynomial kernel:
    \[
    k(x, x') = (x^\top x' + c)^d
    \]

    \item Gaussian kernel:
    \[
    k(x, x') = \exp(-\gamma \|x - x'\|^2)
    \]
\end{itemize}

\textbf{Kernel trick.}\\ \begin{itemize}
    \item Compute inner products without explicitly computing $\phi(x)$
    \item Enables work in very high-dimensional spaces
\end{itemize}

\textbf{Insight.}\\ We can operate in rich feature spaces implicitly through kernels.
\end{frame}

\subsection{Expressivity through Features}
\begin{frame}[t,allowframebreaks]{Expressivity through Features}
\small
The expressive power of a model depends heavily on the choice of features.

\textbf{Key idea.}\\ \begin{itemize}
    \item A linear model with rich features can approximate complex functions
\end{itemize}

\textbf{Tradeoffs.}\\ \begin{itemize}
    \item More expressive features:
    \begin{itemize}
        \item Lower bias
        \item Higher risk of overfitting
    \end{itemize}
\end{itemize}

\textbf{Connection to inductive bias.}\\ \begin{itemize}
    \item Feature map defines the hypothesis space
\end{itemize}

\textbf{Modern perspective.}\\ \begin{itemize}
    \item Deep learning learns feature maps automatically
\end{itemize}

\textbf{Insight.}\\ Representation is a central driver of model performance.
\end{frame}

\subsection{A Unifying Perspective}
\begin{frame}[t,allowframebreaks]{A Unifying Perspective}
\small
Feature maps unify linear and nonlinear modeling:

\begin{itemize}
    \item \textbf{Linear model:}\\ simple structure
    \item \textbf{Feature map:}\\ expressive representation
    \item \textbf{Kernel:}\\ implicit computation in feature space
\end{itemize}

\textbf{Core idea.}\\ Complex learning problems can be solved by combining simple models with powerful representations.

\textbf{Key insight.}\\ The boundary between linear and nonlinear models is often a matter of representation.
\end{frame}

\subsection{Outlook}
\begin{frame}[t,allowframebreaks]{Outlook}
\small
In this lecture, we introduced:

\begin{itemize}
    \item basis expansions and feature maps,
    \item linear models in feature space,
    \item the kernel idea,
    \item expressivity through representations.
\end{itemize}

These ideas lead naturally to kernel-based methods.

In the next lecture, we study support vector machines, where kernels play a central role.

\textbf{Guiding principle.}\\ Representation transforms the problem; learning solves it.
\end{frame}

\section{Lecture 10: Support Vector Machines (Conceptual)}
\SectionDivider{Lecture 10: Support Vector Machines (Conceptual)}
\begin{frame}[t,allowframebreaks]{Outline}
\small
\begin{itemize}
  \item Motivation
  \item Margin Maximization
  \item Hard and Soft Margins
  \item Dual Formulation (Conceptual)
  \item Kernel Methods Revisited
  \item A Unifying Perspective
  \item Outlook
\end{itemize}
\end{frame}

\subsection{Motivation}
\begin{frame}[t,allowframebreaks]{Motivation}
\small
In previous lectures, we studied:
\begin{itemize}
    \item Linear and logistic regression
    \item Feature maps and kernel methods
\end{itemize}

These methods aim to separate data using linear decision boundaries.

However, a key question remains:

\textit{Among all separating hyperplanes, which one should we choose?}

\textbf{Guiding idea.}\\ Choose the decision boundary that is most robust to perturbations.
\end{frame}

\subsection{Margin Maximization}
\begin{frame}[t,allowframebreaks]{Margin Maximization}
\small
Consider binary classification with labels $y_i \in \{-1, +1\}$.

A linear classifier is:
\[
f(x) = w^\top x + b.
\]

\textbf{Decision boundary.}\\ \[
w^\top x + b = 0.
\]

\textbf{Margin.}\\ \begin{itemize}
    \item Distance between the decision boundary and the closest data points
\end{itemize}

\textbf{Geometric margin.}\\ \[
\gamma = \min_i \frac{y_i (w^\top x_i + b)}{\|w\|}.
\]

\textbf{Optimization problem.}\\ \[
\max_{w,b} \; \gamma.
\]

\textbf{Equivalent formulation.}\\ \[
\min_{w,b} \; \|w\|^2 \quad \text{subject to} \quad y_i(w^\top x_i + b) \geq 1.
\]

\textbf{Interpretation.}\\ \begin{itemize}
    \item Maximizing margin improves robustness
    \item Only closest points matter
\end{itemize}

\textbf{Support vectors.}\\ \begin{itemize}
    \item Points that lie on the margin boundary
    \item Determine the classifier
\end{itemize}

\textbf{Insight.}\\ SVM focuses on the most informative data points.
\end{frame}

\subsection{Hard and Soft Margins}
\begin{frame}[t,allowframebreaks]{Hard and Soft Margins}
\small
\textbf{Hard-margin SVM.}\\ \begin{itemize}
    \item Assumes perfectly separable data
    \item Constraints:
    \[
    y_i(w^\top x_i + b) \geq 1
    \]
\end{itemize}

\textbf{Limitation.}\\ \begin{itemize}
    \item Sensitive to noise
    \item Often infeasible in real data
\end{itemize}

\textbf{Soft-margin SVM.}\\ Introduce slack variables $\xi_i \geq 0$:
\[
y_i(w^\top x_i + b) \geq 1 - \xi_i.
\]

\textbf{Objective.}\\ \[
\min_{w,b,\xi} \; \|w\|^2 + C \sum_i \xi_i.
\]

\textbf{Interpretation.}\\ \begin{itemize}
    \item Allows violations of margin constraints
    \item Balances margin size and classification error
\end{itemize}

\textbf{Connection to regularization.}\\ \begin{itemize}
    \item $\|w\|^2$ acts as regularizer
    \item $C$ controls tradeoff between fit and complexity
\end{itemize}

\textbf{Insight.}\\ Soft margins extend SVM to noisy, real-world data.
\end{frame}

\subsection{Dual Formulation (Conceptual)}
\begin{frame}[t,allowframebreaks]{Dual Formulation (Conceptual)}
\small
Rather than solving directly for $w$, we can express the solution in terms of training data.

\textbf{Key idea.}\\ \begin{itemize}
    \item The optimal $w$ can be written as:
    \[
    w = \sum_{i=1}^n \alpha_i y_i x_i.
    \]
\end{itemize}

\textbf{Interpretation.}\\ \begin{itemize}
    \item Only support vectors have $\alpha_i > 0$
    \item Solution depends on a subset of data points
\end{itemize}

\textbf{Inner products.}\\ \[
w^\top x = \sum_i \alpha_i y_i (x_i^\top x).
\]

\textbf{Insight.}\\ The classifier depends only on inner products between data points.
\end{frame}

\subsection{Kernel Methods Revisited}
\begin{frame}[t,allowframebreaks]{Kernel Methods Revisited}
\small
From the dual perspective, predictions depend on:
\[
x_i^\top x.
\]

We replace this with a kernel:
\[
k(x_i, x) = \langle \phi(x_i), \phi(x) \rangle.
\]

\textbf{Kernel SVM.}\\ \[
f(x) = \sum_i \alpha_i y_i k(x_i, x) + b.
\]

\textbf{Effect.}\\ \begin{itemize}
    \item Linear classifier in feature space
    \item Nonlinear classifier in input space
\end{itemize}

\textbf{Key observation.}\\ \begin{itemize}
    \item Kernel trick enables high-dimensional modeling efficiently
\end{itemize}

\textbf{Insight.}\\ SVM combines margin maximization with kernel-based representation.
\end{frame}

\subsection{A Unifying Perspective}
\begin{frame}[t,allowframebreaks]{A Unifying Perspective}
\small
Support vector machines unify multiple ideas:

\begin{itemize}
    \item \textbf{Geometry:}\\ margin maximization
    \item \textbf{Optimization:}\\ constrained quadratic problem
    \item \textbf{Regularization:}\\ control of model complexity
    \item \textbf{Representation:}\\ kernel feature spaces
\end{itemize}

\textbf{Core idea.}\\ A good classifier is one that separates data with maximum robustness.

\textbf{Key insight.}\\ The combination of margin and representation leads to powerful nonlinear models.
\end{frame}

\subsection{Outlook}
\begin{frame}[t,allowframebreaks]{Outlook}
\small
In this lecture, we introduced:

\begin{itemize}
    \item margin maximization,
    \item hard and soft margin SVM,
    \item dual formulation (conceptually),
    \item kernel methods in SVM.
\end{itemize}

These ideas illustrate how geometry, optimization, and representation interact.

In the next lecture, we study model evaluation and selection, focusing on how to assess and compare models in practice.

\textbf{Guiding principle.}\\ Effective learning balances data fit, robustness, and representation.
\end{frame}

\section{Tutorial 5: Kernel Tricks}
\SectionDivider{Tutorial 5: Kernel Tricks}
\begin{frame}[t,allowframebreaks]{Outline}
\small
\begin{itemize}
  \item Objectives
  \item Exercise 1: Explicit Feature Map
  \item Exercise 2: Kernel Without Explicit Mapping
  \item Exercise 3: Linear vs Kernel Classifier
  \item Exercise 4: Decision Boundary Interpretation
  \item Exercise 5: Valid Kernels (Conceptual)
  \item Exercise 6: Kernel Choice and Inductive Bias
  \item Summary
\end{itemize}
\end{frame}

\subsection{Objectives}
\begin{frame}[t,allowframebreaks]{Objectives}
\small
This tutorial develops intuition and technical understanding of kernel methods:

\begin{itemize}
    \item Understand feature maps and implicit representations
    \item Compute and interpret kernel functions
    \item Connect kernels to inner products in feature space
    \item Analyze how kernels induce nonlinear decision boundaries
\end{itemize}

\textbf{Focus.}\\ Understand kernels as a computational and conceptual tool, not just a formula.
\end{frame}

\subsection{Exercise 1: Explicit Feature Map}
\begin{frame}[t,allowframebreaks]{Exercise 1: Explicit Feature Map}
\small
Consider the polynomial kernel:
\[
k(x, x') = (x^\top x' + 1)^2, \quad x, x' \in \mathbb{R}^2.
\]

\textbf{Tasks.}\\ \begin{enumerate}
    \item Expand $k(x,x')$ explicitly.
    \item Find a feature map $\phi(x)$ such that:
    \[
    k(x,x') = \phi(x)^\top \phi(x').
    \]
\end{enumerate}

\textbf{Solution.}\\ Let $x = (x_1, x_2)$, $x' = (x_1', x_2')$.

\begin{enumerate}
\item Expand:
\[
(x^\top x' + 1)^2 =
(x_1 x_1' + x_2 x_2' + 1)^2.
\]

\[
= x_1^2 x_1'^2 + x_2^2 x_2'^2 + 1
+ 2 x_1 x_1' x_2 x_2'
+ 2 x_1 x_1' + 2 x_2 x_2'.
\]

\item One valid feature map:
\[
\phi(x) =
\begin{pmatrix}
x_1^2 \\
x_2^2 \\
\sqrt{2} x_1 x_2 \\
\sqrt{2} x_1 \\
\sqrt{2} x_2 \\
1
\end{pmatrix}.
\]
\end{enumerate}

\textbf{Key takeaway.}\\ Kernels correspond to inner products in higher-dimensional spaces.
\end{frame}

\subsection{Exercise 2: Kernel Without Explicit Mapping}
\begin{frame}[t,allowframebreaks]{Exercise 2: Kernel Without Explicit Mapping}
\small
Consider the Gaussian kernel:
\[
k(x,x') = \exp(-\gamma \|x - x'\|^2).
\]

\textbf{Tasks.}\\ \begin{enumerate}
    \item Explain why computing $\phi(x)$ explicitly is difficult.
    \item Why is the kernel trick useful here?
\end{enumerate}

\textbf{Solution.}\\ \begin{enumerate}
\item The feature map $\phi(x)$ is infinite-dimensional.

\item The kernel allows us to compute inner products directly without constructing $\phi(x)$.
\end{enumerate}

\textbf{Key takeaway.}\\ The kernel trick enables computation in high-dimensional (even infinite) spaces efficiently.
\end{frame}

\subsection{Exercise 3: Linear vs Kernel Classifier}
\begin{frame}[t,allowframebreaks]{Exercise 3: Linear vs Kernel Classifier}
\small
Consider two classifiers:

\begin{itemize}
    \item Linear: $f(x) = w^\top x$
    \item Kernel: $f(x) = \sum_i \alpha_i y_i k(x_i, x)$
\end{itemize}

\textbf{Tasks.}\\ \begin{enumerate}
    \item What is the key difference between these models?
    \item Why can the kernel classifier represent more complex boundaries?
\end{enumerate}

\textbf{Solution.}\\ \begin{enumerate}
\item The kernel classifier depends on training points and pairwise similarities.

\item Because kernels implicitly map data into higher-dimensional spaces, allowing nonlinear separation.
\end{enumerate}

\textbf{Key takeaway.}\\ Kernel methods extend linear models to nonlinear settings.
\end{frame}

\subsection{Exercise 4: Decision Boundary Interpretation}
\begin{frame}[t,allowframebreaks]{Exercise 4: Decision Boundary Interpretation}
\small
Consider a kernel SVM with Gaussian kernel.

\textbf{Tasks.}\\ \begin{enumerate}
    \item Describe qualitatively the decision boundary.
    \item How do support vectors influence the boundary?
\end{enumerate}

\textbf{Solution.}\\ \begin{enumerate}
\item The boundary can be highly nonlinear and adapt to local structure.

\item Each support vector contributes a localized influence via the kernel function.
\end{enumerate}

\textbf{Insight.}\\ Kernel methods produce flexible, data-adaptive decision boundaries.
\end{frame}

\subsection{Exercise 5: Valid Kernels (Conceptual)}
\begin{frame}[t,allowframebreaks]{Exercise 5: Valid Kernels (Conceptual)}
\small
\textbf{Question.}\\ What properties must a function $k(x,x')$ satisfy to be a valid kernel?

\textbf{Discussion points.}\\ \begin{itemize}
    \item Symmetry:
    \[
    k(x,x') = k(x',x)
    \]
    \item Positive semi-definiteness:
    \[
    \sum_{i,j} c_i c_j k(x_i, x_j) \geq 0
    \]
\end{itemize}

\textbf{Insight.}\\ Valid kernels correspond to inner products in some feature space.
\end{frame}

\subsection{Exercise 6: Kernel Choice and Inductive Bias}
\begin{frame}[t,allowframebreaks]{Exercise 6: Kernel Choice and Inductive Bias}
\small
\textbf{Question.}\\ How does the choice of kernel affect learning?

\textbf{Discussion points.}\\ \begin{itemize}
    \item Polynomial kernel:
    \begin{itemize}
        \item Captures global interactions
    \end{itemize}

    \item Gaussian kernel:
    \begin{itemize}
        \item Captures local similarity
    \end{itemize}
\end{itemize}

\textbf{Insight.}\\ Kernels encode inductive bias through similarity structure.
\end{frame}

\subsection{Summary}
\begin{frame}[t,allowframebreaks]{Summary}
\small
\begin{itemize}
    \item Kernels represent inner products in feature space
    \item The kernel trick avoids explicit high-dimensional computation
    \item Kernel methods extend linear models to nonlinear problems
    \item Choice of kernel determines inductive bias and model behavior
\end{itemize}

\textbf{Preparation for next lecture.}\\ We now study how to evaluate and compare models, selecting the best one for a given task.
\end{frame}

\section{Lecture 11: Model Evaluation and Selection}
\SectionDivider{Lecture 11: Model Evaluation and Selection}
\begin{frame}[t,allowframebreaks]{Outline}
\small
\begin{itemize}
  \item Motivation
  \item Training, Validation, and Testing
  \item Cross-Validation
  \item Performance Metrics
  \item ROC Curves and AUC
  \item Model Comparison
  \item A Unifying Perspective
  \item Outlook
\end{itemize}
\end{frame}

\subsection{Motivation}
\begin{frame}[t,allowframebreaks]{Motivation}
\small
In previous lectures, we developed models for prediction:
\begin{itemize}
    \item linear and logistic regression,
    \item kernel methods and support vector machines.
\end{itemize}

However, a fundamental question remains:

\textit{How do we know whether a model is good?}

\textbf{Guiding idea.}\\ Model evaluation estimates how well a model generalizes to unseen data.
\end{frame}

\subsection{Training, Validation, and Testing}
\begin{frame}[t,allowframebreaks]{Training, Validation, and Testing}
\small
We typically split data into three parts:

\textbf{Training set.}\\ \begin{itemize}
    \item Used to fit model parameters
\end{itemize}

\textbf{Validation set.}\\ \begin{itemize}
    \item Used to tune hyperparameters
    \item Guides model selection
\end{itemize}

\textbf{Test set.}\\ \begin{itemize}
    \item Used only for final evaluation
    \item Provides an unbiased estimate of performance
\end{itemize}

\textbf{Workflow.}\\ \begin{enumerate}
    \item Train models on training data
    \item Select best model using validation data
    \item Report final performance on test data
\end{enumerate}

\textbf{Key principle.}\\ The test set must not influence model selection.

\textbf{Insight.}\\ Proper data splitting is essential for reliable evaluation.
\end{frame}

\subsection{Cross-Validation}
\begin{frame}[t,allowframebreaks]{Cross-Validation}
\small
When data is limited, a single split may be unreliable.

\textbf{$k$-fold cross-validation.}\\ \begin{itemize}
    \item Partition data into $k$ subsets (folds)
    \item Train on $k-1$ folds, validate on remaining fold
    \item Repeat for all folds
\end{itemize}

\textbf{Estimate.}\\ \[
\text{CV error} = \frac{1}{k} \sum_{j=1}^k \text{error}_j.
\]

\textbf{Benefits.}\\ \begin{itemize}
    \item More stable estimate of generalization error
    \item Efficient use of data
\end{itemize}

\textbf{Tradeoffs.}\\ \begin{itemize}
    \item Higher computational cost
    \item Choice of $k$ affects bias–variance of estimate
\end{itemize}

\textbf{Insight.}\\ Cross-validation approximates test performance without requiring a large held-out set.
\end{frame}

\subsection{Performance Metrics}
\begin{frame}[t,allowframebreaks]{Performance Metrics}
\small
The choice of metric depends on the task.

\textbf{Regression metrics.}\\ \begin{itemize}
    \item Mean squared error (MSE):
    \[
    \frac{1}{n} \sum (y_i - \hat{y}_i)^2
    \]
    \item Mean absolute error (MAE)
\end{itemize}

\textbf{Classification metrics.}\\ \begin{itemize}
    \item Accuracy:
    \[
    \frac{\text{correct predictions}}{\text{total}}
    \]
    \item Precision, recall
    \item F1-score
\end{itemize}

\textbf{Probabilistic metrics.}\\ \begin{itemize}
    \item Cross-entropy loss
    \item Log-likelihood
\end{itemize}

\textbf{Important consideration.}\\ \begin{itemize}
    \item Metrics should reflect the goal of the task
    \item Different metrics emphasize different aspects of performance
\end{itemize}

\textbf{Insight.}\\ Choosing the right metric is part of defining the learning problem.
\end{frame}

\subsection{ROC Curves and AUC}
\begin{frame}[t,allowframebreaks]{ROC Curves and AUC}
\small
In many classification problems, especially those involving probabilistic predictions, evaluating performance using a single threshold (e.g., accuracy) can be misleading.

\textbf{Key idea.}\\ Instead of fixing a threshold, evaluate performance across all possible thresholds.
\end{frame}

\begin{frame}[t,allowframebreaks]{ROC Curves and AUC}{True and False Positive Rates}
\small
Consider binary classification with labels $y \in \{0,1\}$.

\textbf{Definitions.}\\ \begin{itemize}
    \item \textbf{True Positive Rate (TPR)}\\ (Recall):
    \[
    \text{TPR} = \frac{\text{True Positives}}{\text{Actual Positives}}
    \]

    \item \textbf{False Positive Rate (FPR)}\\ :
    \[
    \text{FPR} = \frac{\text{False Positives}}{\text{Actual Negatives}}
    \]
\end{itemize}
\end{frame}

\begin{frame}[t,allowframebreaks]{ROC Curves and AUC}{ROC Curve}
\small
The \emph{Receiver Operating Characteristic (ROC)} curve plots:
\[
\text{TPR} \quad \text{vs} \quad \text{FPR}
\]
as the classification threshold varies.

\textbf{Interpretation.}\\ \begin{itemize}
    \item Each point corresponds to a different threshold
    \item The curve shows the tradeoff between sensitivity and false alarms
\end{itemize}

\textbf{Special cases.}\\ \begin{itemize}
    \item Perfect classifier: reaches $(0,1)$ (top-left corner)
    \item Random classifier: diagonal line $\text{TPR} = \text{FPR}$
\end{itemize}

\textbf{Insight.}\\ ROC curves evaluate ranking quality rather than fixed-threshold decisions.
\end{frame}

\begin{frame}[t,allowframebreaks]{ROC Curves and AUC}{Area Under the Curve (AUC)}
\small
The \emph{Area Under the ROC Curve (AUC)} is defined as:
\[
\text{AUC} = \int_0^1 \text{TPR}(\text{FPR}) \, d(\text{FPR}).
\]

\textbf{Interpretation.}\\ \begin{itemize}
    \item $\text{AUC} = 1$: perfect classifier
    \item $\text{AUC} = 0.5$: random guessing
\end{itemize}

\textbf{Probabilistic interpretation.}\\ AUC equals the probability that:
\[
\text{score}(x^+) > \text{score}(x^-),
\]
where $x^+$ is a positive example and $x^-$ is a negative example.

\textbf{Insight.}\\ AUC measures how well the model ranks positive examples above negative ones.
\end{frame}

\begin{frame}[t,allowframebreaks]{ROC Curves and AUC}{When to Use ROC and AUC}
\small
\textbf{Advantages.}\\ \begin{itemize}
    \item Threshold-independent evaluation
    \item Useful for imbalanced datasets
    \item Captures ranking performance
\end{itemize}

\textbf{Limitations.}\\ \begin{itemize}
    \item Does not reflect specific operating points
    \item May be overly optimistic in highly skewed settings
\end{itemize}

\textbf{Practical guidance.}\\ \begin{itemize}
    \item Use ROC/AUC when comparing models broadly
    \item Use precision/recall when class imbalance is critical
\end{itemize}
\end{frame}

\begin{frame}[t,allowframebreaks]{ROC Curves and AUC}{Connection to Previous Concepts}
\small
\begin{itemize}
    \item Logistic regression outputs probabilities used for ranking
    \item Cross-entropy optimizes likelihood, not directly AUC
    \item Evaluation metrics define what ``good performance'' means
\end{itemize}

\textbf{Key insight.}\\ Different metrics emphasize different aspects of prediction quality; choosing the right one is essential.
\end{frame}

\subsection{Model Comparison}
\begin{frame}[t,allowframebreaks]{Model Comparison}
\small
We often need to choose between multiple models or hyperparameters.

\textbf{Approach.}\\ \begin{itemize}
    \item Train multiple models
    \item Evaluate using validation or cross-validation
    \item Select best-performing model
\end{itemize}

\textbf{Hyperparameters.}\\ \begin{itemize}
    \item Regularization strength
    \item Kernel parameters
    \item Model complexity
\end{itemize}

\textbf{Overfitting to validation data.}\\ \begin{itemize}
    \item Excessive tuning can overfit validation set
    \item Requires careful experimental design
\end{itemize}

\textbf{Statistical considerations.}\\ \begin{itemize}
    \item Performance differences may be due to noise
    \item Multiple evaluations increase risk of overfitting
\end{itemize}

\textbf{Insight.}\\ Model selection is itself a learning problem and must be treated carefully.
\end{frame}

\subsection{A Unifying Perspective}
\begin{frame}[t,allowframebreaks]{A Unifying Perspective}
\small
Model evaluation connects several aspects of learning:

\begin{itemize}
    \item \textbf{Generalization:}\\ performance on unseen data
    \item \textbf{Optimization:}\\ training minimizes empirical loss
    \item \textbf{Regularization:}\\ controls model complexity
\end{itemize}

\textbf{Core idea.}\\ Evaluation estimates how well learned models approximate the true underlying function.

\textbf{Key insight.}\\ Reliable evaluation requires separating model fitting from model assessment.
\end{frame}

\subsection{Outlook}
\begin{frame}[t,allowframebreaks]{Outlook}
\small
In this lecture, we studied:

\begin{itemize}
    \item data splitting strategies,
    \item cross-validation,
    \item performance metrics,
    \item model comparison.
\end{itemize}

These tools are essential for practical machine learning.

In the next lecture, we return to optimization, focusing on practical aspects such as stochastic gradient descent and modern optimization methods.

\textbf{Guiding principle.}\\ A model is only as good as its performance on unseen data, and evaluation must reflect this.
\end{frame}

\section{Lecture 12: Optimization in Practice}
\SectionDivider{Lecture 12: Optimization in Practice}
\begin{frame}[t,allowframebreaks]{Outline}
\small
\begin{itemize}
  \item Motivation
  \item Stochastic Gradient Descent
  \item Mini-Batching and Efficiency
  \item Momentum and Adaptive Methods
  \item Practical Convergence Behavior
  \item A Unifying Perspective
  \item Outlook
\end{itemize}
\end{frame}

\subsection{Motivation}
\begin{frame}[t,allowframebreaks]{Motivation}
\small
In Lecture 4, we studied gradient-based optimization in an idealized setting.

However, modern machine learning involves:

\begin{itemize}
    \item large datasets,
    \item high-dimensional parameter spaces,
    \item non-convex objectives.
\end{itemize}

This raises practical questions:

\textit{How do we efficiently optimize large-scale models?}

\textbf{Guiding idea.}\\ Efficient optimization relies on stochastic approximations and carefully designed update rules.
\end{frame}

\subsection{Stochastic Gradient Descent}
\begin{frame}[t,allowframebreaks]{Stochastic Gradient Descent}
\small
Recall the empirical objective:
\[
L(\theta) = \frac{1}{n} \sum_{i=1}^n \ell_i(\theta).
\]

\textbf{Gradient descent.}\\ \[
\theta_{t+1} = \theta_t - \eta \nabla L(\theta_t).
\]

\textbf{Stochastic gradient descent (SGD).}\\ \[
\theta_{t+1} = \theta_t - \eta \nabla \ell_i(\theta_t),
\]
where $i$ is sampled randomly.

\textbf{Properties.}\\ \begin{itemize}
    \item Low computational cost per update
    \item Noisy gradient estimates
\end{itemize}

\textbf{Expectation.}\\ \[
\mathbb{E}[\nabla \ell_i(\theta)] = \nabla L(\theta).
\]

\textbf{Insight.}\\ SGD approximates true gradients using random samples.
\end{frame}

\subsection{Mini-Batching and Efficiency}
\begin{frame}[t,allowframebreaks]{Mini-Batching and Efficiency}
\small
Instead of a single sample, we often use a mini-batch $B$:

\[
\theta_{t+1} = \theta_t - \eta \frac{1}{|B|} \sum_{i \in B} \nabla \ell_i(\theta_t).
\]

\textbf{Advantages.}\\ \begin{itemize}
    \item Reduces variance of gradient estimate
    \item Efficient use of parallel hardware (e.g., GPUs)
\end{itemize}

\textbf{Tradeoffs.}\\ \begin{itemize}
    \item Small batches:
    \begin{itemize}
        \item Noisy but fast
    \end{itemize}
    \item Large batches:
    \begin{itemize}
        \item Stable but computationally expensive
    \end{itemize}
\end{itemize}

\textbf{Insight.}\\ Mini-batching balances efficiency and stability.
\end{frame}

\subsection{Momentum and Adaptive Methods}
\begin{frame}[t,allowframebreaks]{Momentum and Adaptive Methods}
\small
\textbf{Momentum.}\\ Introduce a velocity variable:
\[
v_{t+1} = \beta v_t + \nabla L(\theta_t),
\]
\[
\theta_{t+1} = \theta_t - \eta v_{t+1}.
\]

\textbf{Interpretation.}\\ \begin{itemize}
    \item Accumulates past gradients
    \item Smooths updates
\end{itemize}

\textbf{Effect.}\\ \begin{itemize}
    \item Accelerates convergence in consistent directions
    \item Reduces oscillations in ill-conditioned problems
\end{itemize}

\textbf{Adaptive methods.}\\ Adjust learning rates per parameter:
\[
\theta_{t+1} = \theta_t - \eta D_t^{-1} \nabla L(\theta_t),
\]
where $D_t$ depends on past gradients.

\textbf{Examples (conceptually).}\\ \begin{itemize}
    \item Scale updates based on gradient magnitude
    \item Larger steps for infrequent features
\end{itemize}

\textbf{Tradeoffs.}\\ \begin{itemize}
    \item Faster initial progress
    \item Sometimes poorer generalization
\end{itemize}

\textbf{Insight.}\\ Optimization methods influence both convergence and generalization.
\end{frame}

\subsection{Practical Convergence Behavior}
\begin{frame}[t,allowframebreaks]{Practical Convergence Behavior}
\small
\textbf{Non-convexity.}\\ \begin{itemize}
    \item Many local minima and saddle points
    \item Exact convergence guarantees are limited
\end{itemize}

\textbf{Empirical observations.}\\ \begin{itemize}
    \item SGD often finds good solutions
    \item Noise helps escape saddle points
\end{itemize}

\textbf{Learning rate.}\\ \begin{itemize}
    \item Critical hyperparameter
    \item Often decreased over time
\end{itemize}

\textbf{Generalization effects.}\\ \begin{itemize}
    \item Noisy updates can favor flatter minima
    \item Optimization dynamics act as implicit regularization
\end{itemize}

\textbf{Stopping criteria.}\\ \begin{itemize}
    \item Fixed number of iterations
    \item Convergence of loss
    \item Validation performance (early stopping)
\end{itemize}

\textbf{Insight.}\\ Optimization behavior in practice differs significantly from classical theory.
\end{frame}

\subsection{A Unifying Perspective}
\begin{frame}[t,allowframebreaks]{A Unifying Perspective}
\small
Optimization in modern machine learning involves:

\begin{itemize}
    \item \textbf{Stochasticity:}\\ approximate gradients
    \item \textbf{Dynamics:}\\ iterative updates in parameter space
    \item \textbf{Adaptivity:}\\ adjusting updates based on history
    \item \textbf{Regularization:}\\ implicit effects of optimization
\end{itemize}

\textbf{Core idea.}\\ Optimization is not just a computational tool, but a fundamental component of learning.

\textbf{Key insight.}\\ The choice of optimization algorithm affects both training efficiency and generalization performance.
\end{frame}

\subsection{Outlook}
\begin{frame}[t,allowframebreaks]{Outlook}
\small
In this lecture, we studied:

\begin{itemize}
    \item stochastic gradient descent,
    \item mini-batching,
    \item momentum and adaptive methods,
    \item practical convergence behavior.
\end{itemize}

These ideas are essential for training modern machine learning models.

In the next lecture, we examine numerical and statistical stability, focusing on how data scaling and noise affect learning.

\textbf{Guiding principle.}\\ Effective learning depends on both the objective and the optimization dynamics used to minimize it.
\end{frame}

\section{Tutorial 6: Optimization Behavior}
\SectionDivider{Tutorial 6: Optimization Behavior}
\begin{frame}[t,allowframebreaks]{Outline}
\small
\begin{itemize}
  \item Objectives
  \item Exercise 1: Gradient Descent on a Quadratic
  \item Exercise 2: Effect of Learning Rate
  \item Exercise 3: SGD vs Full Gradient
  \item Exercise 4: Role of Batch Size
  \item Exercise 5: Momentum Interpretation
  \item Exercise 6: Implicit Regularization of SGD
  \item Exercise 7: Early Stopping as Regularization
  \item Exercise 8 (Discussion): Flat vs Sharp Minima
  \item Summary
\end{itemize}
\end{frame}

\subsection{Objectives}
\begin{frame}[t,allowframebreaks]{Objectives}
\small
This tutorial develops intuition for how optimization behaves in practice:

\begin{itemize}
    \item Understand the behavior of gradient descent vs SGD
    \item Analyze the effect of learning rate and batch size
    \item Study conditioning and convergence speed
    \item Connect optimization dynamics to generalization
\end{itemize}

\textbf{Focus.}\\ Go beyond formulas and understand how optimization behaves in real training.
\end{frame}

\subsection{Exercise 1: Gradient Descent on a Quadratic}
\begin{frame}[t,allowframebreaks]{Exercise 1: Gradient Descent on a Quadratic}
\small
Consider:
\[
f(w) = \frac{1}{2} w^\top A w,
\]
where $A$ is symmetric positive definite.

\textbf{Tasks.}\\ \begin{enumerate}
    \item Write the gradient descent update.
    \item Describe convergence behavior in terms of eigenvalues of $A$.
\end{enumerate}

\textbf{Solution.}\\ \begin{enumerate}
\item Update:
\[
w_{t+1} = w_t - \eta A w_t.
\]

\item Let eigenvalues of $A$ be $\lambda_1 \leq \cdots \leq \lambda_d$.

\begin{itemize}
    \item Convergence rate depends on condition number:
    \[
    \kappa = \frac{\lambda_{\max}}{\lambda_{\min}}.
    \]
    \item Large $\kappa$ $\Rightarrow$ slow convergence
\end{itemize}
\end{enumerate}

\textbf{Insight.}\\ Ill-conditioning leads to slow and anisotropic convergence.
\end{frame}

\subsection{Exercise 2: Effect of Learning Rate}
\begin{frame}[t,allowframebreaks]{Exercise 2: Effect of Learning Rate}
\small
Consider gradient descent with step size $\eta$.

\textbf{Tasks.}\\ \begin{enumerate}
    \item What happens if $\eta$ is too small?
    \item What happens if $\eta$ is too large?
\end{enumerate}

\textbf{Solution.}\\ \begin{enumerate}
\item Small $\eta$:
\begin{itemize}
    \item Slow convergence
\end{itemize}

\item Large $\eta$:
\begin{itemize}
    \item Oscillations or divergence
\end{itemize}
\end{enumerate}

\textbf{Key takeaway.}\\ Learning rate critically controls stability and speed.
\end{frame}

\subsection{Exercise 3: SGD vs Full Gradient}
\begin{frame}[t,allowframebreaks]{Exercise 3: SGD vs Full Gradient}
\small
Consider minimizing:
\[
L(\theta) = \frac{1}{n} \sum_{i=1}^n \ell_i(\theta).
\]

\textbf{Tasks.}\\ \begin{enumerate}
    \item Compare full gradient descent and SGD updates.
    \item Why does SGD introduce noise?
\end{enumerate}

\textbf{Solution.}\\ \begin{enumerate}
\item 
\begin{itemize}
    \item GD uses full gradient
    \item SGD uses a single sample or mini-batch
\end{itemize}

\item Noise arises because:
\[
\nabla \ell_i(\theta) \neq \nabla L(\theta)
\]
for a single sample.
\end{enumerate}

\textbf{Insight.}\\ SGD trades accuracy of gradients for efficiency.
\end{frame}

\subsection{Exercise 4: Role of Batch Size}
\begin{frame}[t,allowframebreaks]{Exercise 4: Role of Batch Size}
\small
\textbf{Tasks.}\\ \begin{enumerate}
    \item What happens as batch size increases?
    \item How does batch size affect generalization?
\end{enumerate}

\textbf{Solution.}\\ \begin{enumerate}
\item Larger batch:
\begin{itemize}
    \item Lower variance gradient
    \item More stable updates
\end{itemize}

\item Smaller batch:
\begin{itemize}
    \item More noise
    \item Often better generalization
\end{itemize}
\end{enumerate}

\textbf{Insight.}\\ Noise in SGD can act as implicit regularization.
\end{frame}

\subsection{Exercise 5: Momentum Interpretation}
\begin{frame}[t,allowframebreaks]{Exercise 5: Momentum Interpretation}
\small
Consider momentum updates:
\[
v_{t+1} = \beta v_t + \nabla L(\theta_t).
\]

\textbf{Tasks.}\\ \begin{enumerate}
    \item Explain the role of $\beta$.
    \item Why does momentum help in ill-conditioned problems?
\end{enumerate}

\textbf{Solution.}\\ \begin{enumerate}
\item $\beta$ controls how much past gradients influence updates.

\item Momentum:
\begin{itemize}
    \item Accelerates movement in consistent directions
    \item Reduces oscillations in narrow valleys
\end{itemize}
\end{enumerate}

\textbf{Insight.}\\ Momentum smooths optimization trajectories.
\end{frame}

\subsection{Exercise 6: Implicit Regularization of SGD}
\begin{frame}[t,allowframebreaks]{Exercise 6: Implicit Regularization of SGD}
\small
\textbf{Question.}\\ Why might SGD generalize better than deterministic gradient descent?

\textbf{Discussion points.}\\ \begin{itemize}
    \item Noise helps escape sharp minima
    \item Prefers flatter regions of the loss landscape
    \item Acts as implicit regularization
\end{itemize}

\textbf{Insight.}\\ Optimization dynamics influence the solutions found.
\end{frame}

\subsection{Exercise 7: Early Stopping as Regularization}
\begin{frame}[t,allowframebreaks]{Exercise 7: Early Stopping as Regularization}
\small
\textbf{Tasks.}\\ \begin{enumerate}
    \item Why can training too long lead to overfitting?
    \item How does early stopping help?
\end{enumerate}

\textbf{Solution.}\\ \begin{enumerate}
\item Long training:
\begin{itemize}
    \item Model fits noise
\end{itemize}

\item Early stopping:
\begin{itemize}
    \item Limits effective complexity
\end{itemize}
\end{enumerate}

\textbf{Insight.}\\ Training time controls model complexity.
\end{frame}

\subsection{Exercise 8 (Discussion): Flat vs Sharp Minima}
\begin{frame}[t,allowframebreaks]{Exercise 8 (Discussion): Flat vs Sharp Minima}
\small
\textbf{Question.}\\ Why are flat minima often associated with better generalization?

\textbf{Discussion points.}\\ \begin{itemize}
    \item Robustness to perturbations
    \item Stability under noise
    \item Relation to SGD dynamics
\end{itemize}

\textbf{Insight.}\\ Geometry of the loss landscape affects generalization.
\end{frame}

\subsection{Summary}
\begin{frame}[t,allowframebreaks]{Summary}
\small
\begin{itemize}
    \item Optimization behavior depends on learning rate, batch size, and conditioning
    \item SGD introduces noise that affects both convergence and generalization
    \item Momentum and adaptive methods improve optimization efficiency
    \item Optimization dynamics act as implicit regularization
\end{itemize}

\textbf{Preparation for next lecture.}\\ We now study numerical and statistical stability, examining how scaling and noise affect learning behavior.
\end{frame}

\section{Lecture 13: Why Deep Learning?}
\SectionDivider{Lecture 13: Why Deep Learning?}
\begin{frame}[t,allowframebreaks]{Outline}
\small
\begin{itemize}
  \item Motivation
  \item Limitations of Shallow Models
  \item Compositional Structure of Functions
  \item Motivation for Deep Learning
  \item Transition to Learned Representations
  \item A Unifying Perspective
  \item Outlook
\end{itemize}
\end{frame}

\subsection{Motivation}
\begin{frame}[t,allowframebreaks]{Motivation}
\small
So far, we have studied:

\begin{itemize}
    \item linear and logistic regression,
    \item kernel methods and support vector machines,
    \item optimization and generalization principles.
\end{itemize}

These methods are powerful, but they rely heavily on \emph{manually designed features}.

\textbf{Key question.}\\ Can we automatically learn representations from data?

\textbf{Guiding idea.}\\ Deep learning combines representation learning with function approximation.
\end{frame}

\subsection{Limitations of Shallow Models}
\begin{frame}[t,allowframebreaks]{Limitations of Shallow Models}
\small
Classical models typically take the form:
\[
f(x) = w^\top \phi(x),
\]
where $\phi(x)$ is a fixed feature map.

\textbf{Limitations.}\\ \begin{itemize}
    \item \textbf{Feature engineering:}\\ \begin{itemize}
        \item Requires domain expertise
        \item Difficult to scale
    \end{itemize}

    \item \textbf{Expressivity constraints:}\\ \begin{itemize}
        \item May require very high-dimensional feature spaces
    \end{itemize}

    \item \textbf{Lack of hierarchy:}\\ \begin{itemize}
        \item Cannot naturally capture compositional structure
    \end{itemize}
\end{itemize}

\textbf{Insight.}\\ Shallow models separate representation from learning.
\end{frame}

\subsection{Compositional Structure of Functions}
\begin{frame}[t,allowframebreaks]{Compositional Structure of Functions}
\small
Many real-world functions exhibit hierarchical structure.

\textbf{Example.}\\ \begin{itemize}
    \item Images:
    \begin{itemize}
        \item edges $\rightarrow$ shapes $\rightarrow$ objects
    \end{itemize}

    \item Language:
    \begin{itemize}
        \item characters $\rightarrow$ words $\rightarrow$ sentences
    \end{itemize}
\end{itemize}

\textbf{Mathematical structure.}\\ \[
f(x) = f_L(f_{L-1}(\cdots f_1(x)\cdots)).
\]

\textbf{Key idea.}\\ \begin{itemize}
    \item Complex functions can be expressed as compositions of simpler functions
\end{itemize}

\textbf{Insight.}\\ Hierarchical structure suggests using layered models.
\end{frame}

\subsection{Motivation for Deep Learning}
\begin{frame}[t,allowframebreaks]{Motivation for Deep Learning}
\small
Deep learning models take the form:
\[
f(x) = f_L \circ f_{L-1} \circ \cdots \circ f_1(x).
\]

\textbf{Advantages.}\\ \begin{itemize}
    \item \textbf{Representation learning:}\\ \begin{itemize}
        \item Features are learned automatically
    \end{itemize}

    \item \textbf{Expressivity:}\\ \begin{itemize}
        \item Deep models can represent complex functions efficiently
    \end{itemize}

    \item \textbf{Parameter sharing:}\\ \begin{itemize}
        \item Reduces complexity in structured problems
    \end{itemize}
\end{itemize}

\textbf{Comparison with kernel methods.}\\ \begin{itemize}
    \item Kernels:
    \begin{itemize}
        \item Fixed feature space
    \end{itemize}

    \item Deep learning:
    \begin{itemize}
        \item Learned feature space
    \end{itemize}
\end{itemize}

\textbf{Insight.}\\ Deep learning replaces manual feature design with learned representations.
\end{frame}

\subsection{Transition to Learned Representations}
\begin{frame}[t,allowframebreaks]{Transition to Learned Representations}
\small
We reinterpret learning as:

\[
f(x) = w^\top \phi(x)
\quad \longrightarrow \quad
f(x) = w^\top \phi_\theta(x),
\]
where $\phi_\theta$ is learned from data.

\textbf{Key shift.}\\ \begin{itemize}
    \item Classical ML:
    \begin{itemize}
        \item Fixed $\phi(x)$
    \end{itemize}

    \item Deep learning:
    \begin{itemize}
        \item Learn $\phi_\theta(x)$ jointly with $w$
    \end{itemize}
\end{itemize}

\textbf{Optimization problem.}\\ \[
\min_{\theta, w} \; \sum_{i=1}^n \ell(w^\top \phi_\theta(x_i), y_i).
\]

\textbf{Challenge.}\\ \begin{itemize}
    \item Highly non-convex optimization
\end{itemize}

\textbf{Opportunity.}\\ \begin{itemize}
    \item Rich, task-specific representations
\end{itemize}

\textbf{Insight.}\\ Learning representations and predictors jointly increases modeling power.
\end{frame}

\subsection{A Unifying Perspective}
\begin{frame}[t,allowframebreaks]{A Unifying Perspective}
\small
Deep learning builds on earlier ideas:

\begin{itemize}
    \item \textbf{Function approximation:}\\ complex hypothesis spaces
    \item \textbf{Optimization:}\\ gradient-based training
    \item \textbf{Generalization:}\\ controlled through implicit and explicit regularization
    \item \textbf{Representation:}\\ learned feature maps
\end{itemize}

\textbf{Core idea.}\\ Deep learning integrates representation and prediction into a single optimization problem.

\textbf{Key insight.}\\ The power of deep learning comes from composing simple transformations into complex structures.
\end{frame}

\subsection{Outlook}
\begin{frame}[t,allowframebreaks]{Outlook}
\small
In this lecture, we examined:

\begin{itemize}
    \item limitations of shallow models,
    \item compositional structure of real-world functions,
    \item motivation for deep learning,
    \item transition to learned representations.
\end{itemize}

We now begin the study of neural networks.

In the next lecture, we introduce neural networks as compositional models and study their expressive power.

\textbf{Guiding principle.}\\ Learning good representations is as important as learning the final predictor.
\end{frame}

\section{Lecture 14: Neural Networks as Compositional Models}
\SectionDivider{Lecture 14: Neural Networks as Compositional Models}
\begin{frame}[t,allowframebreaks]{Outline}
\small
\begin{itemize}
  \item Motivation
  \item Layered Architectures
  \item Expressivity and Depth
  \item Universal Approximation (Insight)
  \item Representation Learning
  \item A Unifying Perspective
  \item Outlook
\end{itemize}
\end{frame}

\subsection{Motivation}
\begin{frame}[t,allowframebreaks]{Motivation}
\small
In the previous lecture, we argued that many real-world functions exhibit compositional structure and that learning representations is essential.

\textbf{Key question.}\\ How can we construct models that explicitly capture compositional structure?

\textbf{Answer.}\\ Neural networks implement function composition through layered architectures.
\end{frame}

\subsection{Layered Architectures}
\begin{frame}[t,allowframebreaks]{Layered Architectures}
\small
A neural network is defined as a composition of layers:
\[
f(x) = f_L \circ f_{L-1} \circ \cdots \circ f_1(x).
\]

Each layer has the form:
\[
f_\ell(h) = \sigma(W_\ell h + b_\ell),
\]
where:
\begin{itemize}
    \item $W_\ell$ is a weight matrix,
    \item $b_\ell$ is a bias vector,
    \item $\sigma$ is a nonlinear activation function.
\end{itemize}

\textbf{Forward pass.}\\ \[
h_1 = \sigma(W_1 x + b_1),
\]
\[
h_2 = \sigma(W_2 h_1 + b_2),
\]
\[
\vdots
\]
\[
f(x) = W_L h_{L-1} + b_L.
\]

\textbf{Interpretation.}\\ \begin{itemize}
    \item Each layer transforms its input into a new representation
    \item Representations become progressively more abstract
\end{itemize}

\textbf{Insight.}\\ Neural networks are structured compositions of simple functions.
\end{frame}

\subsection{Expressivity and Depth}
\begin{frame}[t,allowframebreaks]{Expressivity and Depth}
\small
\textbf{Key question.}\\ Why use multiple layers instead of a single large layer?

\textbf{Observation.}\\ \begin{itemize}
    \item Shallow networks can approximate many functions
    \item But may require exponentially many units
\end{itemize}

\textbf{Depth advantage.}\\ \begin{itemize}
    \item Deep networks can represent certain functions more efficiently
    \item Composition allows reuse of intermediate computations
\end{itemize}

\textbf{Example intuition.}\\ \begin{itemize}
    \item Repeated composition builds hierarchical features
\end{itemize}

\textbf{Tradeoff.}\\ \begin{itemize}
    \item Depth increases expressivity
    \item Also increases optimization complexity
\end{itemize}

\textbf{Insight.}\\ Depth enables efficient representation of complex functions.
\end{frame}

\subsection{Universal Approximation (Insight)}
\begin{frame}[t,allowframebreaks]{Universal Approximation (Insight)}
\small
\textbf{Statement (informal).}\\ A neural network with a single hidden layer and sufficient width can approximate any continuous function on a compact domain.

\textbf{Implications.}\\ \begin{itemize}
    \item Neural networks are highly expressive
    \item Depth is not required for expressivity in principle
\end{itemize}

\textbf{Limitation.}\\ \begin{itemize}
    \item Approximation may require very large width
\end{itemize}

\textbf{Key insight.}\\ \begin{itemize}
    \item Depth is about \emph{efficiency}, not just expressivity
\end{itemize}
\end{frame}

\subsection{Representation Learning}
\begin{frame}[t,allowframebreaks]{Representation Learning}
\small
Neural networks learn feature maps automatically:
\[
\phi_\theta(x) = h_{L-1}.
\]

\textbf{Model.}\\ \[
f(x) = w^\top \phi_\theta(x).
\]

\textbf{Interpretation.}\\ \begin{itemize}
    \item Early layers learn low-level features
    \item Later layers learn high-level abstractions
\end{itemize}

\textbf{Contrast with classical methods.}\\ \begin{itemize}
    \item Classical ML:
    \begin{itemize}
        \item Fixed features
    \end{itemize}
    \item Neural networks:
    \begin{itemize}
        \item Learned features
    \end{itemize}
\end{itemize}

\textbf{Insight.}\\ Representation learning is a central advantage of deep learning.
\end{frame}

\subsection{A Unifying Perspective}
\begin{frame}[t,allowframebreaks]{A Unifying Perspective}
\small
Neural networks unify key ideas from earlier lectures:

\begin{itemize}
    \item \textbf{Function approximation:}\\ large hypothesis spaces
    \item \textbf{Optimization:}\\ gradient-based training
    \item \textbf{Representation:}\\ learned feature maps
    \item \textbf{Composition:}\\ hierarchical structure
\end{itemize}

\textbf{Core idea.}\\ A neural network is a parameterized composition of functions trained end-to-end.

\textbf{Key insight.}\\ The power of neural networks lies in combining representation learning with compositional structure.
\end{frame}

\subsection{Outlook}
\begin{frame}[t,allowframebreaks]{Outlook}
\small
In this lecture, we introduced:

\begin{itemize}
    \item layered neural network architectures,
    \item expressivity and the role of depth,
    \item the universal approximation principle,
    \item representation learning.
\end{itemize}

To train neural networks, we need efficient methods to compute gradients.

In the next lecture, we study backpropagation and automatic differentiation, which enable scalable training of deep models.

\textbf{Guiding principle.}\\ Deep models learn powerful representations through composition and optimization.
\end{frame}

\section{Tutorial 7: Building an MLP and Backpropagation}
\SectionDivider{Tutorial 7: Building an MLP and Backpropagation}
\begin{frame}[t,allowframebreaks]{Outline}
\small
\begin{itemize}
  \item Objectives
  \item Exercise 1: A Simple Neural Network
  \item Exercise 2: Forward Computation (Concrete Example)
  \item Exercise 3: Loss Function
  \item Exercise 4: Gradient with Respect to Output Layer
  \item Exercise 5: Backpropagation to Hidden Layer
  \item Exercise 6: Gradients for First Layer
  \item Exercise 7: Computational Perspective
  \item Exercise 8 (Discussion): What Backpropagation Does
  \item Summary
\end{itemize}
\end{frame}

\subsection{Objectives}
\begin{frame}[t,allowframebreaks]{Objectives}
\small
This tutorial introduces neural networks as concrete computational systems:

\begin{itemize}
    \item Construct a simple multi-layer perceptron (MLP)
    \item Perform forward computation step by step
    \item Derive gradients manually for a small network
    \item Understand the need for backpropagation
\end{itemize}

\textbf{Focus.}\\ Understand how neural networks compute and how gradients flow through them.
\end{frame}

\subsection{Exercise 1: A Simple Neural Network}
\begin{frame}[t,allowframebreaks]{Exercise 1: A Simple Neural Network}
\small
Consider a two-layer neural network:

\[
h = \sigma(W_1 x + b_1), \quad
\hat{y} = W_2 h + b_2.
\]

\textbf{Tasks.}\\ \begin{enumerate}
    \item Specify the dimensions of all variables.
    \item Describe the forward pass.
\end{enumerate}

\textbf{Solution.}\\ Let:
\begin{itemize}
    \item $x \in \mathbb{R}^d$
    \item $W_1 \in \mathbb{R}^{m \times d}$, $b_1 \in \mathbb{R}^m$
    \item $h \in \mathbb{R}^m$
    \item $W_2 \in \mathbb{R}^{1 \times m}$, $b_2 \in \mathbb{R}$
\end{itemize}

Forward pass:
\begin{itemize}
    \item Linear transformation: $z_1 = W_1 x + b_1$
    \item Nonlinearity: $h = \sigma(z_1)$
    \item Output: $\hat{y} = W_2 h + b_2$
\end{itemize}

\textbf{Key takeaway.}\\ Neural networks alternate between linear transformations and nonlinearities.
\end{frame}

\subsection{Exercise 2: Forward Computation (Concrete Example)}
\begin{frame}[t,allowframebreaks]{Exercise 2: Forward Computation (Concrete Example)}
\small
Let:
\[
x = \begin{pmatrix}1 \\ 2\end{pmatrix}, \quad
W_1 =
\begin{pmatrix}
1 & -1 \\
0 & 2
\end{pmatrix}, \quad
b_1 = \begin{pmatrix}0 \\ 1\end{pmatrix}.
\]

Use ReLU activation: $\sigma(z) = \max(0,z)$.

\textbf{Tasks.}\\ \begin{enumerate}
    \item Compute $z_1$
    \item Compute $h$
\end{enumerate}

\textbf{Solution.}\\ \[
z_1 =
\begin{pmatrix}
1\cdot1 + (-1)\cdot2 \\
0\cdot1 + 2\cdot2
\end{pmatrix}
+
\begin{pmatrix}0 \\ 1\end{pmatrix}
=
\begin{pmatrix}-1 \\ 5\end{pmatrix}.
\]

\[
h = \begin{pmatrix}0 \\ 5\end{pmatrix}.
\]

\textbf{Insight.}\\ Nonlinearity introduces piecewise behavior.
\end{frame}

\subsection{Exercise 3: Loss Function}
\begin{frame}[t,allowframebreaks]{Exercise 3: Loss Function}
\small
Let the target be $y$ and define squared loss:
\[
L = \frac{1}{2} (\hat{y} - y)^2.
\]

\textbf{Task.}\\ \begin{enumerate}
    \item Express $L$ in terms of network parameters.
\end{enumerate}

\textbf{Solution.}\\ \[
L = \frac{1}{2} \left(W_2 \sigma(W_1 x + b_1) + b_2 - y\right)^2.
\]

\textbf{Insight.}\\ The loss is a nested composition of functions.
\end{frame}

\subsection{Exercise 4: Gradient with Respect to Output Layer}
\begin{frame}[t,allowframebreaks]{Exercise 4: Gradient with Respect to Output Layer}
\small
\textbf{Tasks.}\\ \begin{enumerate}
    \item Compute $\frac{\partial L}{\partial W_2}$
    \item Compute $\frac{\partial L}{\partial b_2}$
\end{enumerate}

\textbf{Solution.}\\ Let $\hat{y} = W_2 h + b_2$.

\[
\frac{\partial L}{\partial \hat{y}} = \hat{y} - y.
\]

\[
\frac{\partial L}{\partial W_2} = (\hat{y} - y) h^\top.
\]

\[
\frac{\partial L}{\partial b_2} = \hat{y} - y.
\]

\textbf{Insight.}\\ Gradients at the output layer are straightforward.
\end{frame}

\subsection{Exercise 5: Backpropagation to Hidden Layer}
\begin{frame}[t,allowframebreaks]{Exercise 5: Backpropagation to Hidden Layer}
\small
\textbf{Tasks.}\\ \begin{enumerate}
    \item Compute $\frac{\partial L}{\partial h}$
    \item Compute $\frac{\partial L}{\partial z_1}$
\end{enumerate}

\textbf{Solution.}\\ \[
\frac{\partial L}{\partial h} = (\hat{y} - y) W_2^\top.
\]

For ReLU:
\[
\frac{\partial h}{\partial z_1} =
\begin{cases}
1 & z_1 > 0 \\
0 & z_1 \leq 0
\end{cases}
\]

\[
\frac{\partial L}{\partial z_1} =
\frac{\partial L}{\partial h} \odot \sigma'(z_1).
\]

\textbf{Insight.}\\ Gradients are modulated by activation derivatives.
\end{frame}

\subsection{Exercise 6: Gradients for First Layer}
\begin{frame}[t,allowframebreaks]{Exercise 6: Gradients for First Layer}
\small
\textbf{Tasks.}\\ \begin{enumerate}
    \item Compute $\frac{\partial L}{\partial W_1}$
    \item Compute $\frac{\partial L}{\partial b_1}$
\end{enumerate}

\textbf{Solution.}\\ \[
\frac{\partial L}{\partial W_1} =
\frac{\partial L}{\partial z_1} x^\top.
\]

\[
\frac{\partial L}{\partial b_1} =
\frac{\partial L}{\partial z_1}.
\]

\textbf{Insight.}\\ Gradients propagate backward through layers via the chain rule.
\end{frame}

\subsection{Exercise 7: Computational Perspective}
\begin{frame}[t,allowframebreaks]{Exercise 7: Computational Perspective}
\small
\textbf{Question.}\\ Why is naive differentiation inefficient for deep networks?

\textbf{Discussion.}\\ \begin{itemize}
    \item Repeated computation of intermediate derivatives
    \item Exponential growth in computation if not reused
\end{itemize}

\textbf{Key idea.}\\ Backpropagation reuses intermediate results efficiently.
\end{frame}

\subsection{Exercise 8 (Discussion): What Backpropagation Does}
\begin{frame}[t,allowframebreaks]{Exercise 8 (Discussion): What Backpropagation Does}
\small
\textbf{Question.}\\ What is the core idea of backpropagation?

\textbf{Discussion points.}\\ \begin{itemize}
    \item Apply chain rule systematically
    \item Propagate gradients from output to input
    \item Reuse computations
\end{itemize}

\textbf{Insight.}\\ Backpropagation is dynamic programming for gradient computation.
\end{frame}

\subsection{Summary}
\begin{frame}[t,allowframebreaks]{Summary}
\small
\begin{itemize}
    \item Neural networks compute via layered transformations
    \item Gradients are computed using the chain rule
    \item Backpropagation efficiently propagates gradients
    \item Training neural networks requires structured gradient computation
\end{itemize}

\textbf{Preparation for next lecture.}\\ We now formalize backpropagation and automatic differentiation, enabling efficient training of deep networks.
\end{frame}

\section{Lecture 15: Backpropagation and Automatic Differentiation}
\SectionDivider{Lecture 15: Backpropagation and Automatic Differentiation}
\begin{frame}[t,allowframebreaks]{Outline}
\small
\begin{itemize}
  \item Motivation
  \item Computational Graphs
  \item Chain Rule and Reverse-Mode Differentiation
  \item Efficient Gradient Computation
  \item Backpropagation in Neural Networks
  \item Implementation Perspectives
  \item A Unifying Perspective
  \item Outlook
\end{itemize}
\end{frame}

\subsection{Motivation}
\begin{frame}[t,allowframebreaks]{Motivation}
\small
In the previous tutorial, we manually computed gradients in a small neural network using the chain rule.

However, modern neural networks may contain:
\begin{itemize}
    \item millions or billions of parameters,
    \item deeply nested compositions of functions.
\end{itemize}

\textbf{Key question.}\\ How can we compute gradients efficiently in such systems?

\textbf{Answer.}\\ Backpropagation provides a systematic and efficient method for computing gradients.
\end{frame}

\subsection{Computational Graphs}
\begin{frame}[t,allowframebreaks]{Computational Graphs}
\small
A computational graph represents a function as a composition of elementary operations.

\textbf{Example.}\\ \[
z = f(x,y) = (x \cdot y + y)^2.
\]

We can represent this as a graph:
\begin{itemize}
    \item $a = x \cdot y$
    \item $b = a + y$
    \item $z = b^2$
\end{itemize}

\textbf{Structure.}\\ \begin{itemize}
    \item Nodes: intermediate variables
    \item Edges: functional dependencies
\end{itemize}

\textbf{Insight.}\\ Complex functions can be decomposed into simple operations.
\end{frame}

\subsection{Chain Rule and Reverse-Mode Differentiation}
\begin{frame}[t,allowframebreaks]{Chain Rule and Reverse-Mode Differentiation}
\small
\textbf{Chain rule.}\\ For composed functions:
\[
z = f(g(h(x))),
\]
we have:
\[
\frac{dz}{dx} = \frac{dz}{df} \cdot \frac{df}{dg} \cdot \frac{dg}{dh} \cdot \frac{dh}{dx}.
\]

\textbf{Forward vs reverse mode.}\\ \begin{itemize}
    \item Forward mode:
    \begin{itemize}
        \item propagate derivatives from inputs to outputs
    \end{itemize}

    \item Reverse mode (backpropagation):
    \begin{itemize}
        \item propagate gradients from outputs to inputs
    \end{itemize}
\end{itemize}

\textbf{Key observation.}\\ \begin{itemize}
    \item Neural networks typically have many parameters and one output
    \item Reverse mode is more efficient in this setting
\end{itemize}

\textbf{Insight.}\\ Backpropagation is reverse-mode automatic differentiation.
\end{frame}

\subsection{Efficient Gradient Computation}
\begin{frame}[t,allowframebreaks]{Efficient Gradient Computation}
\small
Let $L$ be the loss function and $\theta$ the parameters.

\textbf{Goal.}\\ \[
\nabla_\theta L.
\]

\textbf{Backpropagation procedure.}\\ \begin{enumerate}
    \item \textbf{Forward pass:}\\ \begin{itemize}
        \item Compute all intermediate variables
    \end{itemize}

    \item \textbf{Backward pass:}\\ \begin{itemize}
        \item Compute gradients starting from output
        \item Use chain rule to propagate backward
    \end{itemize}
\end{enumerate}

\textbf{Local gradients.}\\ \begin{itemize}
    \item Each node computes partial derivatives with respect to its inputs
\end{itemize}

\textbf{Gradient accumulation.}\\ \[
\frac{\partial L}{\partial x} = \sum_{\text{paths}} \prod \text{local derivatives}.
\]

\textbf{Efficiency.}\\ \begin{itemize}
    \item Each intermediate result is reused
    \item Computational cost comparable to forward pass
\end{itemize}

\textbf{Insight.}\\ Backpropagation avoids redundant computation through dynamic programming.
\end{frame}

\subsection{Backpropagation in Neural Networks}
\begin{frame}[t,allowframebreaks]{Backpropagation in Neural Networks}
\small
Consider a layer:
\[
z_\ell = W_\ell h_{\ell-1} + b_\ell, \quad h_\ell = \sigma(z_\ell).
\]

\textbf{Backward pass.}\\ Let $\delta_\ell = \frac{\partial L}{\partial z_\ell}$.

\textbf{Recursion.}\\ \[
\delta_\ell = (W_{\ell+1}^\top \delta_{\ell+1}) \odot \sigma'(z_\ell).
\]

\textbf{Gradients.}\\ \[
\frac{\partial L}{\partial W_\ell} = \delta_\ell h_{\ell-1}^\top,
\]
\[
\frac{\partial L}{\partial b_\ell} = \delta_\ell.
\]

\textbf{Interpretation.}\\ \begin{itemize}
    \item Gradients flow backward through layers
    \item Each layer transforms gradients using local derivatives
\end{itemize}

\textbf{Insight.}\\ Backpropagation is a structured application of the chain rule across layers.
\end{frame}

\subsection{Implementation Perspectives}
\begin{frame}[t,allowframebreaks]{Implementation Perspectives}
\small
Modern machine learning frameworks implement automatic differentiation.

\textbf{Key ideas.}\\ \begin{itemize}
    \item Build computational graph during forward pass
    \item Store intermediate values
    \item Automatically compute gradients in backward pass
\end{itemize}

\textbf{Advantages.}\\ \begin{itemize}
    \item No need to manually derive gradients
    \item Works for arbitrary compositions of functions
\end{itemize}

\textbf{Tradeoffs.}\\ \begin{itemize}
    \item Memory cost for storing intermediate values
    \item Numerical stability considerations
\end{itemize}

\textbf{Insight.}\\ Automatic differentiation turns gradient computation into a programmable operation.
\end{frame}

\subsection{A Unifying Perspective}
\begin{frame}[t,allowframebreaks]{A Unifying Perspective}
\small
Backpropagation unifies several ideas:

\begin{itemize}
    \item \textbf{Chain rule:}\\ foundation of gradient computation
    \item \textbf{Computation graphs:}\\ structure of functions
    \item \textbf{Dynamic programming:}\\ efficient reuse of computations
    \item \textbf{Optimization:}\\ enables gradient-based learning
\end{itemize}

\textbf{Core idea.}\\ Gradients of complex functions can be computed efficiently by propagating local derivatives backward.

\textbf{Key insight.}\\ Backpropagation makes training deep neural networks computationally feasible.
\end{frame}

\subsection{Outlook}
\begin{frame}[t,allowframebreaks]{Outlook}
\small
In this lecture, we introduced:

\begin{itemize}
    \item computational graphs,
    \item reverse-mode differentiation,
    \item efficient gradient computation,
    \item implementation via automatic differentiation.
\end{itemize}

We now turn to the practical challenges of training deep networks.

In the next lecture, we study initialization, vanishing and exploding gradients, and optimization difficulties in deep learning.

\textbf{Guiding principle.}\\ Efficient gradient computation is the foundation of modern deep learning.
\end{frame}

\section{Lecture 16: Training Deep Networks}
\SectionDivider{Lecture 16: Training Deep Networks}
\begin{frame}[t,allowframebreaks]{Outline}
\small
\begin{itemize}
  \item Motivation
  \item Initialization Strategies
  \item Vanishing and Exploding Gradients
  \item Optimization Challenges
  \item Empirical Training Behavior
  \item Practical Remedies (Preview)
  \item A Unifying Perspective
  \item Outlook
\end{itemize}
\end{frame}

\subsection{Motivation}
\begin{frame}[t,allowframebreaks]{Motivation}
\small
In the previous lecture, we introduced backpropagation as an efficient method for computing gradients.

However, in practice, training deep neural networks presents significant challenges.

\textbf{Key question.}\\ Why is training deep networks difficult, and how can we address these challenges?

\textbf{Guiding idea.}\\ Training behavior depends on initialization, gradient dynamics, and optimization structure.
\end{frame}

\subsection{Initialization Strategies}
\begin{frame}[t,allowframebreaks]{Initialization Strategies}
\small
Neural networks are typically initialized randomly:
\[
W_\ell \sim \text{random distribution}.
\]

\textbf{Why initialization matters.}\\ \begin{itemize}
    \item Affects scale of activations
    \item Influences gradient flow
    \item Determines training stability
\end{itemize}

\textbf{Poor initialization.}\\ \begin{itemize}
    \item Too small:
    \begin{itemize}
        \item Signals vanish across layers
    \end{itemize}

    \item Too large:
    \begin{itemize}
        \item Activations and gradients explode
    \end{itemize}
\end{itemize}

\textbf{Principle of variance preservation.}\\ \begin{itemize}
    \item Choose initialization so that:
    \begin{itemize}
        \item activations maintain stable variance across layers
        \item gradients do not shrink or explode
    \end{itemize}
\end{itemize}

\textbf{Insight.}\\ Initialization sets the starting geometry of optimization.
\end{frame}

\subsection{Vanishing and Exploding Gradients}
\begin{frame}[t,allowframebreaks]{Vanishing and Exploding Gradients}
\small
During backpropagation, gradients are multiplied repeatedly:

\[
\frac{\partial L}{\partial x} =
\prod_{\ell} W_\ell^\top \cdot \sigma'(z_\ell).
\]

\textbf{Vanishing gradients.}\\ \begin{itemize}
    \item Gradients shrink exponentially
    \item Early layers learn very slowly
\end{itemize}

\textbf{Exploding gradients.}\\ \begin{itemize}
    \item Gradients grow exponentially
    \item Training becomes unstable
\end{itemize}

\textbf{Causes.}\\ \begin{itemize}
    \item Repeated matrix multiplications
    \item Activation functions with small or large derivatives
\end{itemize}

\textbf{Consequences.}\\ \begin{itemize}
    \item Difficulty training deep networks
    \item Sensitivity to hyperparameters
\end{itemize}

\textbf{Insight.}\\ Gradient flow degrades as depth increases.
\end{frame}

\subsection{Optimization Challenges}
\begin{frame}[t,allowframebreaks]{Optimization Challenges}
\small
Deep learning optimization differs from classical settings.

\textbf{Non-convexity.}\\ \begin{itemize}
    \item Many local minima and saddle points
\end{itemize}

\textbf{Ill-conditioning.}\\ \begin{itemize}
    \item Different directions have different curvature
    \item Slows convergence
\end{itemize}

\textbf{Coupled parameters.}\\ \begin{itemize}
    \item Layers interact in complex ways
\end{itemize}

\textbf{Empirical observation.}\\ \begin{itemize}
    \item Gradient-based methods often work well despite theoretical difficulty
\end{itemize}

\textbf{Insight.}\\ Optimization landscapes are complex but structured.
\end{frame}

\subsection{Empirical Training Behavior}
\begin{frame}[t,allowframebreaks]{Empirical Training Behavior}
\small
In practice, several patterns emerge:

\textbf{Rapid initial progress.}\\ \begin{itemize}
    \item Loss decreases quickly early in training
\end{itemize}

\textbf{Plateaus.}\\ \begin{itemize}
    \item Slow progress due to flat regions or saddle points
\end{itemize}

\textbf{Generalization behavior.}\\ \begin{itemize}
    \item Training loss may continue decreasing
    \item Test performance may peak earlier (overfitting)
\end{itemize}

\textbf{Role of SGD.}\\ \begin{itemize}
    \item Noise helps escape poor regions
    \item Encourages flatter solutions
\end{itemize}

\textbf{Insight.}\\ Training dynamics are influenced by both optimization and stochasticity.
\end{frame}

\subsection{Practical Remedies (Preview)}
\begin{frame}[t,allowframebreaks]{Practical Remedies (Preview)}
\small
Several techniques address training challenges:

\begin{itemize}
    \item Careful initialization
    \item Choice of activation functions
    \item Normalization methods
    \item Architectural design (e.g., residual connections)
\end{itemize}

These will be studied in subsequent lectures.
\end{frame}

\subsection{A Unifying Perspective}
\begin{frame}[t,allowframebreaks]{A Unifying Perspective}
\small
Training deep networks involves:

\begin{itemize}
    \item \textbf{Initialization:}\\ setting the starting point
    \item \textbf{Gradient flow:}\\ propagation of learning signals
    \item \textbf{Optimization dynamics:}\\ navigating complex landscapes
    \item \textbf{Stochasticity:}\\ shaping generalization
\end{itemize}

\textbf{Core idea.}\\ Effective training requires controlling both signal propagation and optimization dynamics.

\textbf{Key insight.}\\ Depth introduces both expressive power and optimization difficulty.
\end{frame}

\subsection{Outlook}
\begin{frame}[t,allowframebreaks]{Outlook}
\small
In this lecture, we examined:

\begin{itemize}
    \item initialization strategies,
    \item vanishing and exploding gradients,
    \item optimization challenges,
    \item empirical training behavior.
\end{itemize}

To address these issues, we introduce techniques that stabilize training.

In the next lecture, we study activation functions and architectural design choices that improve performance.

\textbf{Guiding principle.}\\ Successful deep learning requires both expressive models and stable training dynamics.
\end{frame}

\section{Tutorial 8: Gradient Flow and Instability}
\SectionDivider{Tutorial 8: Gradient Flow and Instability}
\begin{frame}[t,allowframebreaks]{Outline}
\small
\begin{itemize}
  \item Objectives
  \item Exercise 1: Gradient Through a Deep Linear Network
  \item Exercise 2: Scalar Case Intuition
  \item Exercise 3: Effect of Activation Functions
  \item Exercise 4: Variance Propagation (Conceptual)
  \item Exercise 5: Gradient Flow as a Dynamical System
  \item Exercise 6: Practical Diagnosis
  \item Exercise 7 (Discussion): Remedies
  \item Summary
\end{itemize}
\end{frame}

\subsection{Objectives}
\begin{frame}[t,allowframebreaks]{Objectives}
\small
This tutorial develops intuition for gradient propagation in deep networks:

\begin{itemize}
    \item Understand how gradients evolve across layers
    \item Analyze vanishing and exploding gradients
    \item Study the role of activation functions
    \item Examine the effect of initialization
\end{itemize}

\textbf{Focus.}\\ Understand gradient flow as a dynamical process across layers.
\end{frame}

\subsection{Exercise 1: Gradient Through a Deep Linear Network}
\begin{frame}[t,allowframebreaks]{Exercise 1: Gradient Through a Deep Linear Network}
\small
Consider a deep linear network:
\[
f(x) = W_L W_{L-1} \cdots W_1 x.
\]

\textbf{Tasks.}\\ \begin{enumerate}
    \item Write the gradient $\frac{\partial L}{\partial x}$.
    \item Explain how depth affects gradient magnitude.
\end{enumerate}

\textbf{Solution.}\\ \[
\frac{\partial L}{\partial x}
= W_1^\top W_2^\top \cdots W_L^\top \frac{\partial L}{\partial f}.
\]

\begin{itemize}
    \item Gradient is a product of matrices
    \item Norm scales roughly as product of operator norms
\end{itemize}

\textbf{Insight.}\\ Repeated multiplication leads to exponential growth or decay.
\end{frame}

\subsection{Exercise 2: Scalar Case Intuition}
\begin{frame}[t,allowframebreaks]{Exercise 2: Scalar Case Intuition}
\small
Consider:
\[
f(x) = a^L x.
\]

\textbf{Tasks.}\\ \begin{enumerate}
    \item Compute $\frac{df}{dx}$.
    \item Analyze behavior for:
    \begin{itemize}
        \item $|a| < 1$
        \item $|a| > 1$
    \end{itemize}
\end{enumerate}

\textbf{Solution.}\\ \[
\frac{df}{dx} = a^L.
\]

\begin{itemize}
    \item $|a| < 1$ $\Rightarrow$ vanishing gradients
    \item $|a| > 1$ $\Rightarrow$ exploding gradients
\end{itemize}

\textbf{Insight.}\\ Depth amplifies small deviations multiplicatively.
\end{frame}

\subsection{Exercise 3: Effect of Activation Functions}
\begin{frame}[t,allowframebreaks]{Exercise 3: Effect of Activation Functions}
\small
Consider activations:

\begin{itemize}
    \item Sigmoid: $\sigma'(z) \leq 0.25$
    \item ReLU: $\sigma'(z) \in \{0,1\}$
\end{itemize}

\textbf{Tasks.}\\ \begin{enumerate}
    \item Why do sigmoid networks suffer from vanishing gradients?
    \item Why does ReLU mitigate this problem?
\end{enumerate}

\textbf{Solution.}\\ \begin{enumerate}
\item Sigmoid:
\begin{itemize}
    \item Derivatives are small ($<1$)
    \item Repeated multiplication $\Rightarrow$ vanishing gradients
\end{itemize}

\item ReLU:
\begin{itemize}
    \item Derivative is $1$ in active region
    \item Preserves gradient magnitude
\end{itemize}
\end{enumerate}

\textbf{Insight.}\\ Activation choice directly affects gradient flow.
\end{frame}

\subsection{Exercise 4: Variance Propagation (Conceptual)}
\begin{frame}[t,allowframebreaks]{Exercise 4: Variance Propagation (Conceptual)}
\small
Assume:
\[
h_\ell = W_\ell h_{\ell-1},
\]
with random weights.

\textbf{Tasks.}\\ \begin{enumerate}
    \item What happens if weights are too large?
    \item What happens if weights are too small?
\end{enumerate}

\textbf{Solution.}\\ \begin{itemize}
    \item Large weights:
    \begin{itemize}
        \item Activations grow exponentially
        \item Exploding gradients
    \end{itemize}

    \item Small weights:
    \begin{itemize}
        \item Activations shrink
        \item Vanishing gradients
    \end{itemize}
\end{itemize}

\textbf{Insight.}\\ Initialization controls signal propagation.
\end{frame}

\subsection{Exercise 5: Gradient Flow as a Dynamical System}
\begin{frame}[t,allowframebreaks]{Exercise 5: Gradient Flow as a Dynamical System}
\small
\textbf{Question.}\\ Why can we think of gradient propagation as a dynamical system?

\textbf{Discussion.}\\ \begin{itemize}
    \item Each layer transforms gradients
    \item Repeated transformations resemble iterative dynamics
\end{itemize}

\textbf{Insight.}\\ Training stability depends on the stability of this dynamical system.
\end{frame}

\subsection{Exercise 6: Practical Diagnosis}
\begin{frame}[t,allowframebreaks]{Exercise 6: Practical Diagnosis}
\small
\textbf{Tasks.}\\ \begin{enumerate}
    \item How can you detect vanishing gradients during training?
    \item How can you detect exploding gradients?
\end{enumerate}

\textbf{Solution.}\\ \begin{itemize}
    \item Vanishing:
    \begin{itemize}
        \item Near-zero gradients in early layers
        \item Slow learning
    \end{itemize}

    \item Exploding:
    \begin{itemize}
        \item Very large gradients
        \item Unstable loss or divergence
    \end{itemize}
\end{itemize}

\textbf{Insight.}\\ Monitoring gradients is essential in practice.
\end{frame}

\subsection{Exercise 7 (Discussion): Remedies}
\begin{frame}[t,allowframebreaks]{Exercise 7 (Discussion): Remedies}
\small
\textbf{Question.}\\ What techniques can mitigate gradient instability?

\textbf{Discussion points.}\\ \begin{itemize}
    \item Proper initialization
    \item ReLU-type activations
    \item Normalization methods
    \item Residual connections
\end{itemize}

\textbf{Insight.}\\ Modern deep learning is largely about stabilizing gradient flow.
\end{frame}

\subsection{Summary}
\begin{frame}[t,allowframebreaks]{Summary}
\small
\begin{itemize}
    \item Gradient flow involves repeated transformations across layers
    \item Vanishing and exploding gradients arise from multiplicative effects
    \item Activation functions and initialization strongly influence stability
    \item Training deep networks requires controlling gradient dynamics
\end{itemize}

\textbf{Preparation for next lecture.}\\ We now study activation functions and architectural design choices that improve gradient flow and expressivity.
\end{frame}

\section{Lecture 17: Activation Functions and Depth}
\SectionDivider{Lecture 17: Activation Functions and Depth}
\begin{frame}[t,allowframebreaks]{Outline}
\small
\begin{itemize}
  \item Motivation
  \item Nonlinear Activations
  \item Architectural Design Patterns
  \item Depth versus Width
  \item Expressivity Considerations
  \item A Unifying Perspective
  \item Outlook
\end{itemize}
\end{frame}

\subsection{Motivation}
\begin{frame}[t,allowframebreaks]{Motivation}
\small
In the previous lecture, we saw that training deep networks is difficult due to unstable gradient flow.

\textbf{Key question.}\\ How can we design networks that both:
\begin{itemize}
    \item are expressive,
    \item and are trainable?
\end{itemize}

\textbf{Guiding idea.}\\ Activation functions and architectural design play a central role in controlling both expressivity and gradient flow.
\end{frame}

\subsection{Nonlinear Activations}
\begin{frame}[t,allowframebreaks]{Nonlinear Activations}
\small
A neural network layer has the form:
\[
h = \sigma(Wx + b),
\]
where $\sigma$ is a nonlinear activation function.

\textbf{Why nonlinearity is necessary.}\\ \begin{itemize}
    \item Without $\sigma$, composition of layers is linear:
    \[
    W_L \cdots W_1 x = W x
    \]
    \item Depth would not increase expressivity
\end{itemize}

\textbf{Common activation functions.}\\ \begin{itemize}
    \item \textbf{Sigmoid:}\\ \[
    \sigma(z) = \frac{1}{1 + e^{-z}}
    \]

    \item \textbf{Tanh:}\\ \[
    \tanh(z)
    \]

    \item \textbf{ReLU:}\\ \[
    \text{ReLU}(z) = \max(0, z)
    \]
\end{itemize}

\textbf{Comparison.}\\ \begin{itemize}
    \item Sigmoid / tanh:
    \begin{itemize}
        \item smooth
        \item saturate $\Rightarrow$ vanishing gradients
    \end{itemize}

    \item ReLU:
    \begin{itemize}
        \item non-saturating for $z > 0$
        \item preserves gradient magnitude
    \end{itemize}
\end{itemize}

\textbf{Insight.}\\ Activation functions control both expressivity and gradient flow.
\end{frame}

\subsection{Architectural Design Patterns}
\begin{frame}[t,allowframebreaks]{Architectural Design Patterns}
\small
Beyond individual layers, architecture matters.

\textbf{Basic pattern.}\\ \begin{itemize}
    \item Linear transformation $\rightarrow$ nonlinearity
\end{itemize}

\textbf{Stacking layers.}\\ \begin{itemize}
    \item Builds hierarchical representations
\end{itemize}

\textbf{Key considerations.}\\ \begin{itemize}
    \item Depth (number of layers)
    \item Width (number of units per layer)
    \item Activation choice
\end{itemize}

\textbf{Design principle.}\\ Structure should reflect the underlying structure of the data.
\end{frame}

\subsection{Depth versus Width}
\begin{frame}[t,allowframebreaks]{Depth versus Width}
\small
\textbf{Two ways to increase model capacity:}\\ \begin{itemize}
    \item Increase width (more neurons per layer)
    \item Increase depth (more layers)
\end{itemize}

\textbf{Width.}\\ \begin{itemize}
    \item Can approximate complex functions (universal approximation)
    \item May require exponentially many units
\end{itemize}

\textbf{Depth.}\\ \begin{itemize}
    \item Enables compositional structure
    \item More parameter-efficient
\end{itemize}

\textbf{Key intuition.}\\ \begin{itemize}
    \item Depth allows reuse of intermediate features
    \item Represents hierarchical computations
\end{itemize}

\textbf{Tradeoff.}\\ \begin{itemize}
    \item Deeper networks are harder to train
\end{itemize}

\textbf{Insight.}\\ Depth provides efficient expressivity but introduces optimization challenges.
\end{frame}

\subsection{Expressivity Considerations}
\begin{frame}[t,allowframebreaks]{Expressivity Considerations}
\small
\textbf{Piecewise linear perspective (ReLU networks).}\\ \begin{itemize}
    \item ReLU networks partition input space into regions
    \item Each region corresponds to a linear function
\end{itemize}

\textbf{Effect of depth.}\\ \begin{itemize}
    \item Number of linear regions grows rapidly with depth
    \item Enables highly complex decision boundaries
\end{itemize}

\textbf{Representation viewpoint.}\\ \begin{itemize}
    \item Early layers: simple features
    \item Deeper layers: abstract features
\end{itemize}

\textbf{Connection to kernels.}\\ \begin{itemize}
    \item Kernel methods use fixed feature maps
    \item Deep networks learn feature maps dynamically
\end{itemize}

\textbf{Insight.}\\ Expressivity depends on both architecture and learned representation.
\end{frame}

\subsection{A Unifying Perspective}
\begin{frame}[t,allowframebreaks]{A Unifying Perspective}
\small
Activation functions and depth together determine:

\begin{itemize}
    \item \textbf{Expressivity:}\\ ability to represent complex functions
    \item \textbf{Trainability:}\\ stability of gradient flow
    \item \textbf{Inductive bias:}\\ structure imposed on solutions
\end{itemize}

\textbf{Core idea.}\\ Deep networks balance expressivity and trainability through architectural design.

\textbf{Key insight.}\\ The success of deep learning depends as much on architecture as on optimization.
\end{frame}

\subsection{Outlook}
\begin{frame}[t,allowframebreaks]{Outlook}
\small
In this lecture, we studied:

\begin{itemize}
    \item nonlinear activation functions,
    \item architectural design patterns,
    \item depth versus width tradeoffs,
    \item expressivity of neural networks.
\end{itemize}

However, activation functions alone do not fully solve training instability.

In the next lecture, we study normalization methods, which stabilize training and enable deeper networks.

\textbf{Guiding principle.}\\ Designing deep networks requires balancing expressivity with stable gradient propagation.
\end{frame}

\section{Lecture 18: Normalization Methods}
\SectionDivider{Lecture 18: Normalization Methods}
\begin{frame}[t,allowframebreaks]{Outline}
\small
\begin{itemize}
  \item Motivation
  \item Batch Normalization
  \item Layer and Instance Normalization
  \item Stabilization of Training
  \item Theoretical and Practical Insights
  \item A Unifying Perspective
  \item Outlook
\end{itemize}
\end{frame}

\subsection{Motivation}
\begin{frame}[t,allowframebreaks]{Motivation}
\small
In the previous lectures, we examined two key aspects of training deep networks:

\begin{itemize}
    \item gradient instability (vanishing and exploding gradients),
    \item architectural choices such as activation functions and depth.
\end{itemize}

Even with careful initialization and suitable activations, training deep networks can remain unstable.

\textbf{Key question.}\\ How can we stabilize the distribution of activations during training?

\textbf{Guiding idea.}\\ Normalization methods control the scale and distribution of intermediate representations.
\end{frame}

\subsection{Batch Normalization}
\begin{frame}[t,allowframebreaks]{Batch Normalization}
\small
Batch normalization is one of the most influential techniques in modern deep learning.

Consider activations in a layer:
\[
z = W x + b.
\]

Batch normalization normalizes these activations across a mini-batch.

For a batch $\{z_1,\dots,z_m\}$:

\[
\mu_B = \frac{1}{m}\sum_{i=1}^m z_i
\]

\[
\sigma_B^2 = \frac{1}{m}\sum_{i=1}^m (z_i-\mu_B)^2
\]

Normalized activations:

\[
\hat{z}_i = \frac{z_i - \mu_B}{\sqrt{\sigma_B^2 + \epsilon}}
\]

To preserve representational flexibility:

\[
y_i = \gamma \hat{z}_i + \beta
\]

where $\gamma$ and $\beta$ are learnable parameters.

\textbf{Interpretation}\\ \begin{itemize}
\item normalize mean and variance
\item then allow the network to re-scale if needed
\end{itemize}

\textbf{Benefits}\\ \begin{itemize}
\item stabilizes activation distributions
\item accelerates training
\item reduces sensitivity to initialization
\end{itemize}
\end{frame}

\subsection{Layer and Instance Normalization}
\begin{frame}[t,allowframebreaks]{Layer and Instance Normalization}
\small
Batch normalization relies on mini-batch statistics, which may not always be appropriate.

Alternative normalization schemes include:

\textbf{Layer normalization}\\ Normalize across features of a single example:

\[
\mu = \frac{1}{d}\sum_{j=1}^{d} z_j
\]

\[
\sigma^2 = \frac{1}{d}\sum_{j=1}^{d}(z_j-\mu)^2
\]

\textbf{Instance normalization}\\ Normalize each example independently across spatial dimensions (commonly used in vision tasks).

\textbf{Comparison}\\ \begin{itemize}
\item Batch normalization: uses batch statistics
\item Layer normalization: uses feature statistics
\item Instance normalization: uses per-sample spatial statistics
\end{itemize}

\textbf{Practical observation}\\ Layer normalization is commonly used in transformer architectures.
\end{frame}

\subsection{Stabilization of Training}
\begin{frame}[t,allowframebreaks]{Stabilization of Training}
\small
Normalization improves training dynamics in several ways.

\textbf{Gradient stabilization}\\ \begin{itemize}
\item keeps activations within a reasonable range
\item prevents gradients from exploding or vanishing
\end{itemize}

\textbf{Improved conditioning}\\ \begin{itemize}
\item reduces sensitivity to parameter scaling
\item improves optimization geometry
\end{itemize}

\textbf{Regularization effect}\\ \begin{itemize}
\item stochasticity from batch statistics can act as implicit regularization
\end{itemize}

\textbf{Insight}\\ Normalization reshapes the optimization landscape.
\end{frame}

\subsection{Theoretical and Practical Insights}
\begin{frame}[t,allowframebreaks]{Theoretical and Practical Insights}
\small
Although batch normalization was originally motivated heuristically, several explanations have been proposed.

\textbf{Variance control}\\ Normalization prevents signal magnitudes from growing or shrinking across layers.

\textbf{Optimization smoothing}\\ Empirical studies suggest normalization makes the loss landscape smoother and easier to optimize.

\textbf{Gradient flow}\\ Normalized activations help maintain stable gradient magnitudes during backpropagation.

\textbf{Architectural interaction}\\ Normalization often works together with other design choices:

\begin{itemize}
\item residual connections
\item ReLU activations
\item careful initialization
\end{itemize}
\end{frame}

\subsection{A Unifying Perspective}
\begin{frame}[t,allowframebreaks]{A Unifying Perspective}
\small
Normalization methods address a central challenge of deep learning:

\begin{itemize}
\item controlling the propagation of signals through deep networks.
\end{itemize}

They interact with:

\begin{itemize}
\item activation functions,
\item optimization methods,
\item architectural design.
\end{itemize}

\textbf{Core idea}\\ Stable training requires controlling both forward activations and backward gradients.
\end{frame}

\subsection{Outlook}
\begin{frame}[t,allowframebreaks]{Outlook}
\small
In this lecture we introduced:

\begin{itemize}
\item batch normalization,
\item layer and instance normalization,
\item the role of normalization in training stability.
\end{itemize}

These techniques allow much deeper networks to be trained successfully.

In the next lecture, we study regularization methods in deep learning, including dropout and data augmentation, which improve generalization in large neural networks.

\textbf{Guiding principle}\\ Deep learning success depends on carefully balancing expressivity, optimization, and regularization.
\end{frame}

\section{Tutorial 9: Gradient Flow Experiments}
\SectionDivider{Tutorial 9: Gradient Flow Experiments}
\begin{frame}[t,allowframebreaks]{Outline}
\small
\begin{itemize}
  \item Objectives
  \item Exercise 1: Gradient Norm Across Depth
  \item Exercise 2: Activation Function Comparison
  \item Exercise 3: Initialization Effects
  \item Exercise 4: Effect of Depth
  \item Exercise 5: Normalization Effects
  \item Exercise 6: Practical Experiment (Conceptual)
  \item Exercise 7 (Discussion): Diagnosing Training Failures
  \item Summary
\end{itemize}
\end{frame}

\subsection{Objectives}
\begin{frame}[t,allowframebreaks]{Objectives}
\small
This tutorial provides hands-on and conceptual understanding of gradient flow in deep networks:

\begin{itemize}
    \item Observe vanishing and exploding gradients
    \item Compare activation functions
    \item Study the effect of initialization
    \item Understand the role of normalization
\end{itemize}

\textbf{Focus.}\\ Turn theoretical insights about gradient flow into observable behavior.
\end{frame}

\subsection{Exercise 1: Gradient Norm Across Depth}
\begin{frame}[t,allowframebreaks]{Exercise 1: Gradient Norm Across Depth}
\small
Consider a deep neural network with $L$ layers.

\textbf{Task.}\\ \begin{enumerate}
    \item For each layer $\ell$, define:
    \[
    g_\ell = \left\|\frac{\partial L}{\partial h_\ell}\right\|.
    \]
    \item Predict how $g_\ell$ behaves as $\ell$ decreases (toward input).
\end{enumerate}

\textbf{Discussion.}\\ \begin{itemize}
    \item If gradients vanish:
    \[
    g_1 \ll g_L
    \]

    \item If gradients explode:
    \[
    g_1 \gg g_L
    \]
\end{itemize}

\textbf{Insight.}\\ Gradient norms reveal how learning signals propagate.
\end{frame}

\subsection{Exercise 2: Activation Function Comparison}
\begin{frame}[t,allowframebreaks]{Exercise 2: Activation Function Comparison}
\small
Compare sigmoid and ReLU networks of depth $L$.

\textbf{Tasks.}\\ \begin{enumerate}
    \item Explain why sigmoid leads to vanishing gradients.
    \item Explain why ReLU improves gradient flow.
\end{enumerate}

\textbf{Solution.}\\ \begin{itemize}
    \item Sigmoid:
    \begin{itemize}
        \item $\sigma'(z) \leq 0.25$
        \item repeated multiplication $\Rightarrow$ exponential decay
    \end{itemize}

    \item ReLU:
    \begin{itemize}
        \item derivative is $1$ in active region
        \item preserves gradient magnitude
    \end{itemize}
\end{itemize}

\textbf{Insight.}\\ Activation choice directly affects gradient dynamics.
\end{frame}

\subsection{Exercise 3: Initialization Effects}
\begin{frame}[t,allowframebreaks]{Exercise 3: Initialization Effects}
\small
Assume weights are initialized as:
\[
W_{ij} \sim \mathcal{N}(0, \sigma^2).
\]

\textbf{Tasks.}\\ \begin{enumerate}
    \item What happens if $\sigma$ is too small?
    \item What happens if $\sigma$ is too large?
\end{enumerate}

\textbf{Solution.}\\ \begin{itemize}
    \item Small $\sigma$:
    \begin{itemize}
        \item activations shrink
        \item gradients vanish
    \end{itemize}

    \item Large $\sigma$:
    \begin{itemize}
        \item activations grow
        \item gradients explode
    \end{itemize}
\end{itemize}

\textbf{Insight.}\\ Initialization controls both forward and backward signal propagation.
\end{frame}

\subsection{Exercise 4: Effect of Depth}
\begin{frame}[t,allowframebreaks]{Exercise 4: Effect of Depth}
\small
Consider increasing depth $L$.

\textbf{Tasks.}\\ \begin{enumerate}
    \item Describe how gradient stability changes with depth.
    \item Why are shallow networks easier to train?
\end{enumerate}

\textbf{Solution.}\\ \begin{itemize}
    \item Gradient instability increases with depth
    \item Fewer multiplicative terms in shallow networks
\end{itemize}

\textbf{Insight.}\\ Depth amplifies both expressivity and instability.
\end{frame}

\subsection{Exercise 5: Normalization Effects}
\begin{frame}[t,allowframebreaks]{Exercise 5: Normalization Effects}
\small
Consider inserting batch normalization after each layer.

\textbf{Tasks.}\\ \begin{enumerate}
    \item How does normalization affect activations?
    \item How does it affect gradients?
\end{enumerate}

\textbf{Solution.}\\ \begin{itemize}
    \item Activations are rescaled to stable distributions
    \item Gradients become more stable and less sensitive to scale
\end{itemize}

\textbf{Insight.}\\ Normalization stabilizes both forward and backward passes.
\end{frame}

\subsection{Exercise 6: Practical Experiment (Conceptual)}
\begin{frame}[t,allowframebreaks]{Exercise 6: Practical Experiment (Conceptual)}
\small
\textbf{Experiment design.}\\ Train three networks with identical architecture:

\begin{itemize}
    \item Model A: sigmoid activation
    \item Model B: ReLU activation
    \item Model C: ReLU + batch normalization
\end{itemize}

\textbf{Tasks.}\\ \begin{enumerate}
    \item Predict training behavior of each model.
    \item Rank them by training speed and stability.
\end{enumerate}

\textbf{Expected outcome.}\\ \begin{itemize}
    \item Model A:
    \begin{itemize}
        \item slow training, vanishing gradients
    \end{itemize}

    \item Model B:
    \begin{itemize}
        \item faster, but may still be unstable
    \end{itemize}

    \item Model C:
    \begin{itemize}
        \item fastest and most stable
    \end{itemize}
\end{itemize}

\textbf{Insight.}\\ Modern deep learning relies on combining multiple stabilization techniques.
\end{frame}

\subsection{Exercise 7 (Discussion): Diagnosing Training Failures}
\begin{frame}[t,allowframebreaks]{Exercise 7 (Discussion): Diagnosing Training Failures}
\small
\textbf{Question.}\\ How can you diagnose gradient-related issues in practice?

\textbf{Discussion points.}\\ \begin{itemize}
    \item monitor gradient norms per layer
    \item observe training loss behavior
    \item check activation distributions
\end{itemize}

\textbf{Insight.}\\ Understanding gradient flow is essential for debugging deep models.
\end{frame}

\subsection{Summary}
\begin{frame}[t,allowframebreaks]{Summary}
\small
\begin{itemize}
    \item Gradient flow determines whether deep networks can learn
    \item Activation functions, initialization, and normalization all affect it
    \item Instability grows with depth
    \item Practical success depends on controlling these dynamics
\end{itemize}

\textbf{Preparation for next tutorial.}\\ We next study how architectural changes (e.g., residual connections) further improve stability and enable very deep networks.
\end{frame}

\section{Lecture 19: Regularization in Deep Learning}
\SectionDivider{Lecture 19: Regularization in Deep Learning}
\begin{frame}[t,allowframebreaks]{Outline}
\small
\begin{itemize}
  \item Motivation
  \item Dropout and Noise Injection
  \item Weight Decay
  \item Data Augmentation
  \item Implicit Bias of Optimization
  \item A Unifying Perspective
  \item Outlook
\end{itemize}
\end{frame}

\subsection{Motivation}
\begin{frame}[t,allowframebreaks]{Motivation}
\small
Deep neural networks are highly expressive and often overparameterized.

As a result:
\begin{itemize}
    \item they can fit training data extremely well,
    \item but may overfit and generalize poorly.
\end{itemize}

\textbf{Key question.}\\ How can we control model complexity and improve generalization in deep learning?

\textbf{Guiding idea.}\\ Regularization techniques constrain or perturb learning to favor simpler and more robust solutions.
\end{frame}

\subsection{Dropout and Noise Injection}
\begin{frame}[t,allowframebreaks]{Dropout and Noise Injection}
\small
\textbf{Dropout.}\\ During training, randomly set activations to zero:

\[
h_i \mapsto h_i \cdot \xi_i, \quad \xi_i \sim \text{Bernoulli}(p).
\]

\textbf{Effect.}\\ \begin{itemize}
    \item Prevents co-adaptation of neurons
    \item Encourages redundancy and robustness
\end{itemize}

\textbf{Inference.}\\ \begin{itemize}
    \item Use full network with scaled activations
\end{itemize}

\textbf{Interpretation.}\\ \begin{itemize}
    \item Approximate model averaging over subnetworks
\end{itemize}

\textbf{Noise injection.}\\ More generally, we can inject noise into:

\begin{itemize}
    \item inputs,
    \item activations,
    \item gradients.
\end{itemize}

\textbf{Insight.}\\ Noise prevents over-reliance on specific features and improves robustness.
\end{frame}

\subsection{Weight Decay}
\begin{frame}[t,allowframebreaks]{Weight Decay}
\small
Weight decay adds an $\ell_2$ penalty:

\[
\min_\theta \; L(\theta) + \lambda \|\theta\|^2.
\]

\textbf{Interpretation.}\\ \begin{itemize}
    \item Penalizes large weights
    \item Encourages simpler models
\end{itemize}

\textbf{Connection to optimization.}\\ \begin{itemize}
    \item Equivalent to shrinking parameters during updates
\end{itemize}

\textbf{Geometric view.}\\ \begin{itemize}
    \item Constrains solution to lie in a bounded region
\end{itemize}

\textbf{Insight.}\\ Weight decay controls model capacity through parameter magnitude.
\end{frame}

\subsection{Data Augmentation}
\begin{frame}[t,allowframebreaks]{Data Augmentation}
\small
\textbf{Idea.}\\ Generate additional training data by applying transformations.

\textbf{Examples.}\\ \begin{itemize}
    \item Image transformations:
    \begin{itemize}
        \item rotation, cropping, flipping
    \end{itemize}

    \item Text transformations:
    \begin{itemize}
        \item synonym replacement, paraphrasing
    \end{itemize}
\end{itemize}

\textbf{Interpretation.}\\ \begin{itemize}
    \item Encodes invariances in the data
\end{itemize}

\textbf{Effect.}\\ \begin{itemize}
    \item Expands effective dataset
    \item Improves generalization
\end{itemize}

\textbf{Insight.}\\ Data augmentation injects prior knowledge into learning.
\end{frame}

\subsection{Implicit Bias of Optimization}
\begin{frame}[t,allowframebreaks]{Implicit Bias of Optimization}
\small
Even without explicit regularization, optimization itself induces bias.

\textbf{Examples.}\\ \begin{itemize}
    \item Gradient descent in linear models:
    \begin{itemize}
        \item finds minimum-norm solution
    \end{itemize}

    \item SGD in deep networks:
    \begin{itemize}
        \item prefers flatter minima
    \end{itemize}
\end{itemize}

\textbf{Flat vs sharp minima.}\\ \begin{itemize}
    \item Flat minima:
    \begin{itemize}
        \item robust to perturbations
        \item better generalization
    \end{itemize}

    \item Sharp minima:
    \begin{itemize}
        \item sensitive to noise
        \item poorer generalization
    \end{itemize}
\end{itemize}

\textbf{Role of SGD.}\\ \begin{itemize}
    \item Noise in updates biases toward flat regions
\end{itemize}

\textbf{Insight.}\\ Optimization dynamics act as an implicit regularizer.
\end{frame}

\subsection{A Unifying Perspective}
\begin{frame}[t,allowframebreaks]{A Unifying Perspective}
\small
Regularization in deep learning arises from multiple sources:

\begin{itemize}
    \item \textbf{Explicit:}\\ \begin{itemize}
        \item weight decay, dropout
    \end{itemize}

    \item \textbf{Data-driven:}\\ \begin{itemize}
        \item augmentation
    \end{itemize}

    \item \textbf{Implicit:}\\ \begin{itemize}
        \item optimization dynamics
    \end{itemize}
\end{itemize}

\textbf{Core idea.}\\ Generalization is controlled not only by model size, but by how the model is trained.

\textbf{Key insight.}\\ Modern deep learning relies heavily on implicit and data-driven regularization.
\end{frame}

\subsection{Outlook}
\begin{frame}[t,allowframebreaks]{Outlook}
\small
In this lecture, we studied:

\begin{itemize}
    \item dropout and noise injection,
    \item weight decay,
    \item data augmentation,
    \item implicit bias of optimization.
\end{itemize}

These techniques improve generalization in deep networks.

In the next lecture, we study residual networks and skip connections, which enable training of very deep architectures.

\textbf{Guiding principle.}\\ Effective deep learning balances expressivity, stability, and generalization.
\end{frame}

\section{Lecture 20: Residual Networks and Skip Connections}
\SectionDivider{Lecture 20: Residual Networks and Skip Connections}
\begin{frame}[t,allowframebreaks]{Outline}
\small
\begin{itemize}
  \item Motivation
  \item The Degradation Problem
  \item Residual Learning
  \item Gradient Flow in Residual Networks
  \item Dynamical Systems Perspective
  \item Very Deep Architectures
  \item A Unifying Perspective
  \item Outlook
\end{itemize}
\end{frame}

\subsection{Motivation}
\begin{frame}[t,allowframebreaks]{Motivation}
\small
In previous lectures, we studied:

\begin{itemize}
    \item gradient instability (Lecture 16),
    \item activation functions and depth (Lecture 17),
    \item normalization methods (Lecture 18),
    \item regularization (Lecture 19).
\end{itemize}

These techniques improve training, yet a fundamental problem remains.

\textbf{Observation.}\\ As networks become deeper, performance may degrade.

\textbf{Key question.}\\ Why does increasing depth sometimes hurt performance, and how can we fix it?
\end{frame}

\subsection{The Degradation Problem}
\begin{frame}[t,allowframebreaks]{The Degradation Problem}
\small
Empirically, deeper networks can have:

\begin{itemize}
    \item higher training error,
    \item worse optimization behavior,
\end{itemize}

even though they are more expressive.

\textbf{Paradox.}\\ \begin{itemize}
    \item A deeper model should be able to represent a shallow model
    \item Yet optimization fails to find such solutions
\end{itemize}

\textbf{Insight.}\\ The issue is not expressivity, but optimization difficulty.
\end{frame}

\subsection{Residual Learning}
\begin{frame}[t,allowframebreaks]{Residual Learning}
\small
Residual networks introduce skip connections.

Instead of learning:
\[
H(x),
\]
we learn a residual function:
\[
F(x) = H(x) - x.
\]

Thus:
\[
H(x) = x + F(x).
\]

\textbf{Residual block.}\\ \[
y = x + F(x; \theta).
\]

\textbf{Interpretation.}\\ \begin{itemize}
    \item The network learns a correction to the identity
    \item If identity is optimal, set $F(x) = 0$
\end{itemize}

\textbf{Benefit.}\\ \begin{itemize}
    \item Easier optimization
    \item Improved gradient flow
\end{itemize}

\textbf{Insight.}\\ Learning residuals is easier than learning full transformations.
\end{frame}

\subsection{Gradient Flow in Residual Networks}
\begin{frame}[t,allowframebreaks]{Gradient Flow in Residual Networks}
\small
Consider:
\[
y = x + F(x).
\]

\textbf{Gradient.}\\ \[
\frac{\partial L}{\partial x}
= \frac{\partial L}{\partial y}
\left( I + \frac{\partial F}{\partial x} \right).
\]

\textbf{Key observation.}\\ \begin{itemize}
    \item The identity term ensures gradients can flow directly
    \item Avoids vanishing gradients
\end{itemize}

\textbf{Insight.}\\ Skip connections create direct gradient pathways.
\end{frame}

\subsection{Dynamical Systems Perspective}
\begin{frame}[t,allowframebreaks]{Dynamical Systems Perspective}
\small
Residual networks can be interpreted as discretizations of continuous systems.

\textbf{Form.}\\ \[
h_{k+1} = h_k + F(h_k).
\]

This resembles:
\[
\frac{dh}{dt} = F(h),
\]

which is a differential equation.

\textbf{Interpretation.}\\ \begin{itemize}
    \item Depth corresponds to time steps
    \item Network computes a flow in representation space
\end{itemize}

\textbf{Insight.}\\ Deep networks can be viewed as discrete dynamical systems.
\end{frame}

\subsection{Very Deep Architectures}
\begin{frame}[t,allowframebreaks]{Very Deep Architectures}
\small
Residual connections enable training of very deep networks:

\begin{itemize}
    \item hundreds or even thousands of layers
\end{itemize}

\textbf{Key components.}\\ \begin{itemize}
    \item skip connections
    \item normalization (e.g., batch normalization)
    \item suitable activations (e.g., ReLU)
\end{itemize}

\textbf{Empirical observation.}\\ \begin{itemize}
    \item deeper residual networks often perform better
    \item training becomes more stable
\end{itemize}

\textbf{Connection to modern models.}\\ \begin{itemize}
    \item residual connections are used in:
    \begin{itemize}
        \item convolutional networks
        \item transformers
    \end{itemize}
\end{itemize}
\end{frame}

\subsection{A Unifying Perspective}
\begin{frame}[t,allowframebreaks]{A Unifying Perspective}
\small
Residual networks address a fundamental challenge:

\begin{itemize}
    \item enabling depth without sacrificing trainability.
\end{itemize}

They combine:

\begin{itemize}
    \item architectural design (skip connections),
    \item normalization,
    \item activation functions.
\end{itemize}

\textbf{Core idea.}\\ Introduce identity mappings to stabilize both forward and backward signal propagation.

\textbf{Key insight.}\\ Depth becomes useful only when it is trainable.
\end{frame}

\subsection{Outlook}
\begin{frame}[t,allowframebreaks]{Outlook}
\small
In this lecture, we studied:

\begin{itemize}
    \item the degradation problem,
    \item residual learning,
    \item gradient flow in skip connections,
    \item dynamical systems interpretation.
\end{itemize}

Residual networks mark a turning point in deep learning, enabling very deep architectures.

In the next lectures, we study specialized architectures that exploit structure in data:

\begin{itemize}
    \item convolutional networks for spatial data,
    \item sequence models and attention for sequential data.
\end{itemize}

\textbf{Guiding principle.}\\ Architectural design is central to making deep learning effective at scale.
\end{frame}

\section{Tutorial 10: Training Ablations (BatchNorm, Dropout, Residuals)}
\SectionDivider{Tutorial 10: Training Ablations (BatchNorm, Dropout, Residuals)}
\begin{frame}[t,allowframebreaks]{Outline}
\small
\begin{itemize}
  \item Objectives
  \item Exercise 1: What is an Ablation Study?
  \item Exercise 2: Baseline Model
  \item Exercise 3: Effect of Batch Normalization
  \item Exercise 4: Effect of Dropout
  \item Exercise 5: Effect of Residual Connections
  \item Exercise 6: Combined Model
  \item Exercise 7: Comparative Table
  \item Exercise 8 (Discussion): Design Principles
  \item Summary
\end{itemize}
\end{frame}

\subsection{Objectives}
\begin{frame}[t,allowframebreaks]{Objectives}
\small
This tutorial studies how architectural and training choices affect learning:

\begin{itemize}
    \item Understand the role of normalization, regularization, and residual connections
    \item Compare training dynamics under different configurations
    \item Develop intuition through controlled ablation experiments
\end{itemize}

\textbf{Focus.}\\ Identify which components make deep networks trainable and generalize well.
\end{frame}

\subsection{Exercise 1: What is an Ablation Study?}
\begin{frame}[t,allowframebreaks]{Exercise 1: What is an Ablation Study?}
\small
\textbf{Question.}\\ What is the purpose of an ablation study in machine learning?

\textbf{Discussion.}\\ \begin{itemize}
    \item Systematically remove or modify components
    \item Measure their effect on performance
\end{itemize}

\textbf{Insight.}\\ Ablation reveals causal contributions of design choices.
\end{frame}

\subsection{Exercise 2: Baseline Model}
\begin{frame}[t,allowframebreaks]{Exercise 2: Baseline Model}
\small
Consider a deep fully connected network:

\begin{itemize}
    \item depth: $L = 10$
    \item activation: ReLU
    \item optimizer: SGD
\end{itemize}

\textbf{Task.}\\ Predict training behavior without:
\begin{itemize}
    \item normalization
    \item residual connections
\end{itemize}

\textbf{Expected outcome.}\\ \begin{itemize}
    \item unstable gradients
    \item slow or failed convergence
\end{itemize}

\textbf{Insight.}\\ Depth alone does not guarantee learnability.
\end{frame}

\subsection{Exercise 3: Effect of Batch Normalization}
\begin{frame}[t,allowframebreaks]{Exercise 3: Effect of Batch Normalization}
\small
Modify the baseline:

\begin{itemize}
    \item add batch normalization after each layer
\end{itemize}

\textbf{Tasks.}\\ \begin{enumerate}
    \item Predict changes in training speed
    \item Predict effect on gradient stability
\end{enumerate}

\textbf{Expected outcome.}\\ \begin{itemize}
    \item faster convergence
    \item more stable gradients
\end{itemize}

\textbf{Insight.}\\ Normalization stabilizes optimization.
\end{frame}

\subsection{Exercise 4: Effect of Dropout}
\begin{frame}[t,allowframebreaks]{Exercise 4: Effect of Dropout}
\small
Modify the baseline:

\begin{itemize}
    \item add dropout with rate $p$
\end{itemize}

\textbf{Tasks.}\\ \begin{enumerate}
    \item Predict effect on training loss
    \item Predict effect on test performance
\end{enumerate}

\textbf{Expected outcome.}\\ \begin{itemize}
    \item slower training (higher training loss)
    \item improved generalization
\end{itemize}

\textbf{Insight.}\\ Regularization trades training performance for robustness.
\end{frame}

\subsection{Exercise 5: Effect of Residual Connections}
\begin{frame}[t,allowframebreaks]{Exercise 5: Effect of Residual Connections}
\small
Modify the baseline:

\begin{itemize}
    \item replace layers with residual blocks:
    \[
    h_{k+1} = h_k + F(h_k)
    \]
\end{itemize}

\textbf{Tasks.}\\ \begin{enumerate}
    \item Predict training behavior
    \item Compare to plain deep network
\end{enumerate}

\textbf{Expected outcome.}\\ \begin{itemize}
    \item easier optimization
    \item deeper networks become trainable
\end{itemize}

\textbf{Insight.}\\ Residual connections improve gradient flow structurally.
\end{frame}

\subsection{Exercise 6: Combined Model}
\begin{frame}[t,allowframebreaks]{Exercise 6: Combined Model}
\small
Consider a model with:

\begin{itemize}
    \item ReLU activation
    \item batch normalization
    \item residual connections
    \item optional dropout
\end{itemize}

\textbf{Tasks.}\\ \begin{enumerate}
    \item Predict overall training behavior
    \item Explain why this combination is effective
\end{enumerate}

\textbf{Expected outcome.}\\ \begin{itemize}
    \item fast convergence
    \item stable gradients
    \item good generalization
\end{itemize}

\textbf{Insight.}\\ Modern deep learning relies on combining multiple techniques.
\end{frame}

\subsection{Exercise 7: Comparative Table}
\begin{frame}[t,allowframebreaks]{Exercise 7: Comparative Table}
\small
Fill in the following table:

\[
\begin{array}{c|c|c|c}
\text{Model} & \text{Training Stability} & \text{Speed} & \text{Generalization} \\
\hline
\text{Baseline} & & & \\
\text{+ BN} & & & \\
\text{+ Dropout} & & & \\
\text{+ Residual} & & & \\
\text{Combined} & & & \\
\end{array}
\]

\textbf{Discussion.}\\ Students should observe:

\begin{itemize}
    \item BN $\rightarrow$ stability + speed
    \item Dropout $\rightarrow$ generalization
    \item Residual $\rightarrow$ depth scalability
\end{itemize}
\end{frame}

\subsection{Exercise 8 (Discussion): Design Principles}
\begin{frame}[t,allowframebreaks]{Exercise 8 (Discussion): Design Principles}
\small
\textbf{Question.}\\ What general principles can we extract from these experiments?

\textbf{Discussion points.}\\ \begin{itemize}
    \item Control signal propagation (BN, residuals)
    \item Prevent overfitting (dropout)
    \item Combine techniques for best performance
\end{itemize}

\textbf{Insight.}\\ Successful deep learning is the result of carefully engineered systems.
\end{frame}

\subsection{Summary}
\begin{frame}[t,allowframebreaks]{Summary}
\small
\begin{itemize}
    \item Batch normalization stabilizes training
    \item Dropout improves generalization
    \item Residual connections enable depth
    \item Combined techniques yield robust models
\end{itemize}

\textbf{Preparation for next lecture.}\\ We now study convolutional neural networks, which incorporate structural assumptions about data into architecture design.
\end{frame}

\section{Lecture 21: Convolutional Neural Networks}
\SectionDivider{Lecture 21: Convolutional Neural Networks}
\begin{frame}[t,allowframebreaks]{Outline}
\small
\begin{itemize}
  \item Motivation
  \item Convolution as a Structured Linear Operator
  \item Translation Invariance
  \item Pooling and Hierarchies
  \item Architectural Components
  \item Applications in Vision
  \item Fourier Perspective on Convolution (Advanced)
  \item Spectral View of CNN Layers
  \item Connection to Linear Operators
  \item Role of Nonlinearity
  \item Broader Connections
  \item A Unifying Perspective
  \item Outlook
\end{itemize}
\end{frame}

\subsection{Motivation}
\begin{frame}[t,allowframebreaks]{Motivation}
\small
So far, we have studied neural networks as general function approximators.

However, many real-world data types exhibit structure. In particular:

\begin{itemize}
    \item images have spatial structure,
    \item nearby pixels are highly correlated,
    \item patterns repeat across locations.
\end{itemize}

\textbf{Key question.}\\ How can we design neural networks that exploit such structure?

\textbf{Guiding idea.}\\ Convolutional neural networks (CNNs) incorporate prior knowledge about spatial locality and translation invariance.
\end{frame}

\subsection{Convolution as a Structured Linear Operator}
\begin{frame}[t,allowframebreaks]{Convolution as a Structured Linear Operator}
\small
A convolutional layer applies a filter (kernel) across the input.

For a 1D signal:
\[
y[i] = \sum_{k} w[k] \, x[i-k].
\]

For images (2D):
\[
y[i,j] = \sum_{u,v} w[u,v] \, x[i-u, j-v].
\]

\textbf{Interpretation.}\\ \begin{itemize}
    \item Local connectivity: each output depends on a small neighborhood
    \item Weight sharing: the same filter is applied everywhere
\end{itemize}

\textbf{Matrix view.}\\ \begin{itemize}
    \item Convolution corresponds to a structured (Toeplitz) matrix
    \item Far fewer parameters than fully connected layers
\end{itemize}

\textbf{Insight.}\\ Convolution imposes strong structure on linear transformations.
\end{frame}

\subsection{Translation Invariance}
\begin{frame}[t,allowframebreaks]{Translation Invariance}
\small
A key property of convolution:

\textbf{Translation equivariance.}\\ \begin{itemize}
    \item Shifting the input shifts the output
\end{itemize}

\textbf{Why this matters.}\\ \begin{itemize}
    \item Objects may appear anywhere in an image
    \item Model should detect features regardless of location
\end{itemize}

\textbf{Approximate invariance.}\\ \begin{itemize}
    \item Pooling and deeper layers help achieve invariance
\end{itemize}

\textbf{Insight.}\\ CNNs encode translation symmetry as an inductive bias.
\end{frame}

\subsection{Pooling and Hierarchies}
\begin{frame}[t,allowframebreaks]{Pooling and Hierarchies}
\small
\textbf{Pooling layers.}\\ \begin{itemize}
    \item Reduce spatial resolution
    \item Aggregate local information
\end{itemize}

\textbf{Examples.}\\ \begin{itemize}
    \item max pooling
    \item average pooling
\end{itemize}

\textbf{Effect.}\\ \begin{itemize}
    \item increases receptive field
    \item reduces sensitivity to small shifts
\end{itemize}

\textbf{Hierarchical representations.}\\ \begin{itemize}
    \item early layers: edges, textures
    \item deeper layers: shapes, objects
\end{itemize}

\textbf{Insight.}\\ CNNs build representations from local to global structure.
\end{frame}

\subsection{Architectural Components}
\begin{frame}[t,allowframebreaks]{Architectural Components}
\small
Modern CNNs combine several elements:

\begin{itemize}
    \item convolution layers
    \item nonlinear activations (e.g., ReLU)
    \item normalization (e.g., batch normalization)
    \item pooling or strided convolution
    \item residual connections
\end{itemize}

\textbf{Design pattern.}\\ \[
\text{Conv} \rightarrow \text{BN} \rightarrow \text{ReLU}
\]

\textbf{Insight.}\\ CNNs integrate multiple principles developed in earlier lectures.
\end{frame}

\subsection{Applications in Vision}
\begin{frame}[t,allowframebreaks]{Applications in Vision}
\small
CNNs are widely used in:

\begin{itemize}
    \item image classification
    \item object detection
    \item segmentation
\end{itemize}

\textbf{Why CNNs work well.}\\ \begin{itemize}
    \item exploit spatial locality
    \item reduce number of parameters
    \item encode useful invariances
\end{itemize}

\textbf{Limitations.}\\ \begin{itemize}
    \item limited ability to model long-range dependencies
    \item fixed receptive field growth
\end{itemize}

\textbf{Insight.}\\ CNNs are powerful but tailored to specific data structures.
\end{frame}

\subsection{Fourier Perspective on Convolution (Advanced)}
\begin{frame}[t,allowframebreaks]{Fourier Perspective on Convolution (Advanced)}
\small
Convolution has a fundamental interpretation in the frequency domain.

\textbf{Fourier transform.}\\ Let $\mathcal{F}$ denote the Fourier transform. For a signal $x$:
\[
\hat{x}(\omega) = \mathcal{F}[x](\omega).
\]

\textbf{Convolution theorem.}\\ For signals $x$ and $w$:
\[
x * w \;\longleftrightarrow\; \hat{x}(\omega)\,\hat{w}(\omega),
\]
i.e., convolution in the spatial domain corresponds to pointwise multiplication in the frequency domain.

\textbf{Interpretation.}\\ \begin{itemize}
    \item Convolution acts as a \emph{filter} in frequency space
    \item The kernel $w$ determines which frequencies are amplified or suppressed
\end{itemize}

\textbf{Examples.}\\ \begin{itemize}
    \item Smooth filters $\Rightarrow$ low-pass (remove high-frequency noise)
    \item Edge detectors $\Rightarrow$ emphasize high-frequency components
\end{itemize}

\textbf{Insight.}\\ Convolutional layers learn frequency-selective filters.
\end{frame}

\subsection{Spectral View of CNN Layers}
\begin{frame}[t,allowframebreaks]{Spectral View of CNN Layers}
\small
Consider a convolutional layer:
\[
y = w * x.
\]

In the Fourier domain:
\[
\hat{y}(\omega) = \hat{w}(\omega)\,\hat{x}(\omega).
\]

\textbf{Implications.}\\ \begin{itemize}
    \item Each layer modulates the spectrum of the signal
    \item Deep CNNs apply successive spectral transformations
\end{itemize}

\textbf{Insight.}\\ CNNs can be interpreted as hierarchical spectral filters.
\end{frame}

\subsection{Connection to Linear Operators}
\begin{frame}[t,allowframebreaks]{Connection to Linear Operators}
\small
From linear algebra:

\begin{itemize}
    \item Convolution corresponds to a structured matrix (Toeplitz / circulant)
    \item Fourier transform approximately diagonalizes convolution
\end{itemize}

\textbf{Key idea.}\\ \[
\text{Convolution} \;\approx\; \text{diagonal operator in Fourier basis}.
\]

\textbf{Insight.}\\ The Fourier basis reveals the intrinsic structure of convolutional operators.
\end{frame}

\subsection{Role of Nonlinearity}
\begin{frame}[t,allowframebreaks]{Role of Nonlinearity}
\small
CNNs are not purely linear:

\begin{itemize}
    \item nonlinear activations break simple spectral interpretation
    \item interactions between frequencies emerge across layers
\end{itemize}

\textbf{Interpretation.}\\ \begin{itemize}
    \item linear part: frequency filtering
    \item nonlinear part: mixing and recombining frequencies
\end{itemize}

\textbf{Insight.}\\ Deep CNNs combine spectral filtering with nonlinear feature interactions.
\end{frame}

\subsection{Broader Connections}
\begin{frame}[t,allowframebreaks]{Broader Connections}
\small
This perspective connects CNNs to broader mathematical frameworks:

\begin{itemize}
    \item \textbf{Signal processing:}\\ classical filtering theory
    \item \textbf{Harmonic analysis:}\\ decomposition into frequency components
    \item \textbf{Operator theory:}\\ diagonalization of structured operators
\end{itemize}

\textbf{Modern developments.}\\ \begin{itemize}
    \item spectral bias in neural networks
    \item frequency-domain analysis of generalization
\end{itemize}

\textbf{Key insight.}\\ Convolutional neural networks can be understood as learnable spectral filters acting on structured data.
\end{frame}

\subsection{A Unifying Perspective}
\begin{frame}[t,allowframebreaks]{A Unifying Perspective}
\small
CNNs illustrate a central theme in machine learning:

\begin{itemize}
    \item incorporating structure into models improves efficiency and performance.
\end{itemize}

They combine:

\begin{itemize}
    \item inductive bias (locality, translation symmetry),
    \item deep representations,
    \item efficient parameterization.
\end{itemize}

\textbf{Core idea.}\\ Architectural design encodes prior knowledge about the data.

\textbf{Key insight.}\\ Good architectures reduce the burden on learning by building in the right structure.
\end{frame}

\subsection{Outlook}
\begin{frame}[t,allowframebreaks]{Outlook}
\small
In this lecture, we introduced:

\begin{itemize}
    \item convolution as a structured linear operator,
    \item translation invariance,
    \item pooling and hierarchical representations,
    \item applications in vision.
\end{itemize}

Next, we study sequence models, which address data with temporal or sequential structure.

\textbf{Guiding principle.}\\ Different data modalities require different architectural biases.
\end{frame}

\section{Lecture 22: Sequence Models (RNN, LSTM, GRU)}
\SectionDivider{Lecture 22: Sequence Models (RNN, LSTM, GRU)}
\begin{frame}[t,allowframebreaks]{Outline}
\small
\begin{itemize}
  \item Motivation
  \item Recurrent Neural Networks
  \item Long-Term Dependencies
  \item Gated Architectures: LSTM
  \item Gated Recurrent Units (GRU)
  \item Limitations of Recurrence
  \item Continuous-Time and Dynamical Systems Perspective (Advanced)
  \item Connection to Discrete Dynamical Systems
  \item Continuous-Time Limit
  \item Gradient Flow and Stability
  \item Gating as Controlled Dynamics
  \item Connection to Residual Networks
  \item Broader Connections
  \item A Unifying Perspective
  \item Outlook
\end{itemize}
\end{frame}

\subsection{Motivation}
\begin{frame}[t,allowframebreaks]{Motivation}
\small
Many real-world data are sequential:

\begin{itemize}
    \item text (sentences, documents),
    \item speech signals,
    \item time series data.
\end{itemize}

\textbf{Key question.}\\ How can we model data where order and temporal dependencies matter?

\textbf{Guiding idea.}\\ Sequence models process inputs step-by-step, maintaining a hidden state that summarizes past information.
\end{frame}

\subsection{Recurrent Neural Networks}
\begin{frame}[t,allowframebreaks]{Recurrent Neural Networks}
\small
A recurrent neural network (RNN) updates a hidden state over time:

\[
h_t = \sigma(W_h h_{t-1} + W_x x_t + b),
\]

\[
y_t = W_y h_t.
\]

\textbf{Interpretation.}\\ \begin{itemize}
    \item $h_t$ stores information about the past
    \item The same parameters are reused across time steps
\end{itemize}

\textbf{Unrolled view.}\\ \begin{itemize}
    \item RNN can be viewed as a deep network with shared weights
    \item Depth corresponds to sequence length
\end{itemize}

\textbf{Insight.}\\ RNNs extend neural networks to sequential data via parameter sharing over time.
\end{frame}

\subsection{Long-Term Dependencies}
\begin{frame}[t,allowframebreaks]{Long-Term Dependencies}
\small
\textbf{Problem.}\\ \begin{itemize}
    \item Information from early time steps may be lost
    \item Gradients decay or explode over long sequences
\end{itemize}

\textbf{Mathematical intuition.}\\ \[
\frac{\partial L}{\partial h_t}
= \prod_{k=t}^{T} \frac{\partial h_{k+1}}{\partial h_k}.
\]

\textbf{Observation.}\\ \begin{itemize}
    \item Similar to deep networks in space, but now in time
    \item Suffers from vanishing/exploding gradients
\end{itemize}

\textbf{Insight.}\\ Sequence length plays the role of depth in RNNs.
\end{frame}

\subsection{Gated Architectures: LSTM}
\begin{frame}[t,allowframebreaks]{Gated Architectures: LSTM}
\small
Long Short-Term Memory (LSTM) networks introduce gates to control information flow.

\textbf{Core components.}\\ \begin{itemize}
    \item \textbf{Cell state}\\ $c_t$: long-term memory
    \item \textbf{Gates:}\\ \begin{itemize}
        \item input gate
        \item forget gate
        \item output gate
    \end{itemize}
\end{itemize}

\textbf{Update equations (conceptual).}\\ \begin{itemize}
    \item Forget:
    \[
    f_t = \sigma(W_f x_t + U_f h_{t-1})
    \]

    \item Input:
    \[
    i_t = \sigma(W_i x_t + U_i h_{t-1})
    \]

    \item Candidate:
    \[
    \tilde{c}_t = \tanh(W_c x_t + U_c h_{t-1})
    \]

    \item Cell update:
    \[
    c_t = f_t \odot c_{t-1} + i_t \odot \tilde{c}_t
    \]
\end{itemize}

\textbf{Key idea.}\\ \begin{itemize}
    \item gates regulate information flow
    \item prevent uncontrolled decay or explosion
\end{itemize}

\textbf{Insight.}\\ LSTM introduces controlled memory dynamics.
\end{frame}

\subsection{Gated Recurrent Units (GRU)}
\begin{frame}[t,allowframebreaks]{Gated Recurrent Units (GRU)}
\small
GRU is a simplified gated architecture.

\textbf{Key components.}\\ \begin{itemize}
    \item update gate
    \item reset gate
\end{itemize}

\textbf{Characteristics.}\\ \begin{itemize}
    \item fewer parameters than LSTM
    \item similar performance in many tasks
\end{itemize}

\textbf{Insight.}\\ GRU simplifies gating while retaining effectiveness.
\end{frame}

\subsection{Limitations of Recurrence}
\begin{frame}[t,allowframebreaks]{Limitations of Recurrence}
\small
Despite improvements, RNN-based models have limitations.

\textbf{Sequential computation.}\\ \begin{itemize}
    \item cannot fully parallelize across time steps
\end{itemize}

\textbf{Long-range dependencies.}\\ \begin{itemize}
    \item still difficult to capture very long dependencies
\end{itemize}

\textbf{Training difficulty.}\\ \begin{itemize}
    \item optimization remains challenging
\end{itemize}

\textbf{Insight.}\\ Recurrence introduces inherent computational and optimization constraints.
\end{frame}

\subsection{Continuous-Time and Dynamical Systems Perspective (Advanced)}
\begin{frame}[t,allowframebreaks]{Continuous-Time and Dynamical Systems Perspective (Advanced)}
\small
Recurrent neural networks can be interpreted as discrete-time dynamical systems.

Recall the RNN update:
\[
h_t = \sigma(W_h h_{t-1} + W_x x_t + b).
\]

\textbf{Dynamical systems view.}\\ \begin{itemize}
    \item $h_t$ is the state of the system at time $t$
    \item the update defines a nonlinear map:
    \[
    h_t = F(h_{t-1}, x_t)
    \]
\end{itemize}

\textbf{Insight.}\\ RNNs evolve a hidden state through time via repeated application of a nonlinear transformation.
\end{frame}

\subsection{Connection to Discrete Dynamical Systems}
\begin{frame}[t,allowframebreaks]{Connection to Discrete Dynamical Systems}
\small
If inputs are fixed or absent, the system becomes autonomous:
\[
h_t = F(h_{t-1}).
\]

\textbf{Behavior.}\\ \begin{itemize}
    \item fixed points: $h^* = F(h^*)$
    \item stable vs unstable dynamics
    \item attractors in state space
\end{itemize}

\textbf{Interpretation.}\\ RNNs can store information in stable states of the dynamics.
\end{frame}

\subsection{Continuous-Time Limit}
\begin{frame}[t,allowframebreaks]{Continuous-Time Limit}
\small
We can interpret the recurrence as a discretization of a continuous-time system.

Consider:
\[
h_{t+1} = h_t + \Delta t \, F(h_t, x_t).
\]

In the limit $\Delta t \to 0$, this becomes an ordinary differential equation:
\[
\frac{dh(t)}{dt} = F(h(t), x(t)).
\]

\textbf{Interpretation.}\\ \begin{itemize}
    \item hidden state evolves continuously over time
    \item RNN is a discretized flow in state space
\end{itemize}

\textbf{Insight.}\\ RNNs approximate continuous-time dynamical systems.
\end{frame}

\subsection{Gradient Flow and Stability}
\begin{frame}[t,allowframebreaks]{Gradient Flow and Stability}
\small
Backpropagation through time corresponds to sensitivity analysis of the dynamical system.

\textbf{Gradient evolution.}\\ \[
\frac{\partial L}{\partial h_t}
= \left( \prod_{k=t}^{T} \frac{\partial h_{k+1}}{\partial h_k} \right)
\frac{\partial L}{\partial h_T}.
\]

\textbf{Connection to dynamics.}\\ \begin{itemize}
    \item gradients depend on Jacobians of the system
    \item stability of dynamics determines gradient behavior
\end{itemize}

\textbf{Lyapunov intuition.}\\ \begin{itemize}
    \item contracting dynamics $\Rightarrow$ vanishing gradients
    \item expanding dynamics $\Rightarrow$ exploding gradients
\end{itemize}

\textbf{Insight.}\\ Gradient instability reflects instability of the underlying dynamical system.
\end{frame}

\subsection{Gating as Controlled Dynamics}
\begin{frame}[t,allowframebreaks]{Gating as Controlled Dynamics}
\small
LSTM and GRU can be interpreted as modifying the dynamics.

\textbf{Example (LSTM cell).}\\ \[
c_t = f_t \odot c_{t-1} + i_t \odot \tilde{c}_t.
\]

\textbf{Interpretation.}\\ \begin{itemize}
    \item $f_t$ controls memory retention
    \item $i_t$ controls input injection
\end{itemize}

\textbf{Insight.}\\ Gates act as adaptive coefficients controlling system stability.
\end{frame}

\subsection{Connection to Residual Networks}
\begin{frame}[t,allowframebreaks]{Connection to Residual Networks}
\small
Recall residual networks:
\[
h_{k+1} = h_k + F(h_k).
\]

\textbf{Parallel.}\\ \begin{itemize}
    \item ResNet: depth as time (spatial layers)
    \item RNN: time as time (temporal steps)
\end{itemize}

\textbf{Unified view.}\\ \begin{itemize}
    \item both are discretizations of continuous flows
\end{itemize}

\textbf{Insight.}\\ Deep networks and sequence models share a common dynamical systems foundation.
\end{frame}

\subsection{Broader Connections}
\begin{frame}[t,allowframebreaks]{Broader Connections}
\small
This perspective connects to:

\begin{itemize}
    \item \textbf{Neural ODEs:}\\ continuous-depth models
    \item \textbf{Control theory:}\\ stability and feedback
    \item \textbf{Dynamical systems:}\\ attractors and trajectories
\end{itemize}

\textbf{Key insight.}\\ Learning in sequence models can be viewed as learning a dynamical system that evolves representations over time.
\end{frame}

\subsection{A Unifying Perspective}
\begin{frame}[t,allowframebreaks]{A Unifying Perspective}
\small
Sequence models extend neural networks by:

\begin{itemize}
    \item introducing temporal structure,
    \item sharing parameters across time,
    \item maintaining a hidden state.
\end{itemize}

\textbf{Core idea.}\\ Learning involves propagating information through time as well as depth.

\textbf{Key insight.}\\ RNNs treat time as another dimension of computation.
\end{frame}

\subsection{Outlook}
\begin{frame}[t,allowframebreaks]{Outlook}
\small
In this lecture, we introduced:

\begin{itemize}
    \item recurrent neural networks,
    \item long-term dependency problems,
    \item gated architectures (LSTM, GRU),
    \item limitations of recurrence.
\end{itemize}

These limitations motivate alternative architectures.

In the next lecture, we study attention mechanisms and transformers, which overcome many limitations of recurrent models.

\textbf{Guiding principle.}\\ Efficient sequence modeling requires balancing memory, computation, and parallelism.
\end{frame}

\section{Tutorial 11: CNN and Sequence Modeling Intuition}
\SectionDivider{Tutorial 11: CNN and Sequence Modeling Intuition}
\begin{frame}[t,allowframebreaks]{Outline}
\small
\begin{itemize}
  \item Objectives
  \item Exercise 1: Fully Connected vs Convolution
  \item Exercise 2: Local Receptive Fields
  \item Exercise 3: Weight Sharing
  \item Exercise 4: Hierarchical Features
  \item Exercise 5: Sequential Data Processing
  \item Exercise 6: Information Bottleneck in RNNs
  \item Exercise 7: CNN vs RNN Comparison
  \item Exercise 8: When to Use Which?
  \item Exercise 9 (Discussion): Limitations
  \item Summary
\end{itemize}
\end{frame}

\subsection{Objectives}
\begin{frame}[t,allowframebreaks]{Objectives}
\small
This tutorial develops intuition for structure-aware neural networks:

\begin{itemize}
    \item Understand how CNNs process spatial data
    \item Understand how RNNs process sequential data
    \item Compare architectural inductive biases
    \item Relate structure to efficiency and performance
\end{itemize}

\textbf{Focus.}\\ How architectural design reflects the structure of data.
\end{frame}

\subsection{Exercise 1: Fully Connected vs Convolution}
\begin{frame}[t,allowframebreaks]{Exercise 1: Fully Connected vs Convolution}
\small
\textbf{Question.}\\ Why is a fully connected network inefficient for image data?

\textbf{Discussion.}\\ \begin{itemize}
    \item ignores spatial locality
    \item requires many parameters
    \item does not exploit repeating patterns
\end{itemize}

\textbf{Insight.}\\ Ignoring structure leads to inefficiency and poor generalization.
\end{frame}

\subsection{Exercise 2: Local Receptive Fields}
\begin{frame}[t,allowframebreaks]{Exercise 2: Local Receptive Fields}
\small
Consider a convolutional filter applied to an image.

\textbf{Tasks.}\\ \begin{enumerate}
    \item What does a $3 \times 3$ filter capture?
    \item How does stacking layers increase receptive field?
\end{enumerate}

\textbf{Solution.}\\ \begin{itemize}
    \item $3 \times 3$ filter captures local patterns (edges, textures)
    \item stacking layers aggregates larger spatial context
\end{itemize}

\textbf{Insight.}\\ CNNs build global understanding from local computations.
\end{frame}

\subsection{Exercise 3: Weight Sharing}
\begin{frame}[t,allowframebreaks]{Exercise 3: Weight Sharing}
\small
\textbf{Question.}\\ Why is weight sharing important in CNNs?

\textbf{Discussion.}\\ \begin{itemize}
    \item same feature can appear anywhere
    \item reduces number of parameters
    \item improves generalization
\end{itemize}

\textbf{Insight.}\\ Weight sharing encodes translation symmetry.
\end{frame}

\subsection{Exercise 4: Hierarchical Features}
\begin{frame}[t,allowframebreaks]{Exercise 4: Hierarchical Features}
\small
\textbf{Task.}\\ Describe what different CNN layers learn.

\textbf{Discussion.}\\ \begin{itemize}
    \item early layers: edges and textures
    \item middle layers: shapes and parts
    \item deep layers: objects and semantics
\end{itemize}

\textbf{Insight.}\\ Deep networks build hierarchical representations.
\end{frame}

\subsection{Exercise 5: Sequential Data Processing}
\begin{frame}[t,allowframebreaks]{Exercise 5: Sequential Data Processing}
\small
Consider a sequence $(x_1, x_2, \dots, x_T)$.

\textbf{Task.}\\ Explain how an RNN processes this sequence.

\textbf{Solution.}\\ \begin{itemize}
    \item processes inputs one step at a time
    \item maintains hidden state $h_t$
    \item summarizes past information in $h_t$
\end{itemize}

\textbf{Insight.}\\ RNNs compress past information into a state.
\end{frame}

\subsection{Exercise 6: Information Bottleneck in RNNs}
\begin{frame}[t,allowframebreaks]{Exercise 6: Information Bottleneck in RNNs}
\small
\textbf{Question.}\\ Why is it difficult for RNNs to capture long-term dependencies?

\textbf{Discussion.}\\ \begin{itemize}
    \item all past information must pass through $h_t$
    \item gradients degrade over long sequences
\end{itemize}

\textbf{Insight.}\\ The hidden state acts as a bottleneck.
\end{frame}

\subsection{Exercise 7: CNN vs RNN Comparison}
\begin{frame}[t,allowframebreaks]{Exercise 7: CNN vs RNN Comparison}
\small
Fill in the table:

\[
\begin{array}{c|c|c}
\text{Aspect} & \text{CNN} & \text{RNN} \\
\hline
\text{Structure} & & \\
\text{Computation} & & \\
\text{Parallelism} & & \\
\text{Inductive bias} & & \\
\end{array}
\]

\textbf{Discussion.}\\ \begin{itemize}
    \item CNN:
    \begin{itemize}
        \item spatial structure
        \item parallel computation
        \item locality + translation invariance
    \end{itemize}

    \item RNN:
    \begin{itemize}
        \item temporal structure
        \item sequential computation
        \item memory over time
    \end{itemize}
\end{itemize}

\textbf{Insight.}\\ Different architectures encode different structural assumptions.
\end{frame}

\subsection{Exercise 8: When to Use Which?}
\begin{frame}[t,allowframebreaks]{Exercise 8: When to Use Which?}
\small
\textbf{Task.}\\ Match model types to data:

\begin{itemize}
    \item images
    \item language
    \item time series
\end{itemize}

\textbf{Solution.}\\ \begin{itemize}
    \item images $\rightarrow$ CNN
    \item language $\rightarrow$ RNN / Transformer
    \item time series $\rightarrow$ RNN / hybrid models
\end{itemize}

\textbf{Insight.}\\ Model choice should reflect data structure.
\end{frame}

\subsection{Exercise 9 (Discussion): Limitations}
\begin{frame}[t,allowframebreaks]{Exercise 9 (Discussion): Limitations}
\small
\textbf{Question.}\\ What are the limitations of CNNs and RNNs?

\textbf{Discussion.}\\ \begin{itemize}
    \item CNN:
    \begin{itemize}
        \item limited long-range interaction
    \end{itemize}

    \item RNN:
    \begin{itemize}
        \item sequential bottleneck
        \item difficult to parallelize
    \end{itemize}
\end{itemize}

\textbf{Insight.}\\ Limitations motivate new architectures.
\end{frame}

\subsection{Summary}
\begin{frame}[t,allowframebreaks]{Summary}
\small
\begin{itemize}
    \item CNNs exploit spatial locality and symmetry
    \item RNNs model temporal dependencies via hidden states
    \item Both rely on structural inductive bias
    \item Limitations lead to attention-based models
\end{itemize}

\textbf{Preparation for next lecture.}\\ We next study attention mechanisms, which overcome limitations of both CNNs and RNNs by modeling global interactions directly.
\end{frame}

\section{Lecture 23: Attention Mechanisms}
\SectionDivider{Lecture 23: Attention Mechanisms}
\begin{frame}[t,allowframebreaks]{Outline}
\small
\begin{itemize}
  \item Motivation
  \item Basic Attention Mechanism
  \item Self-Attention
  \item Scaled Dot-Product Attention
  \item Multi-Head Attention
  \item Comparison with CNNs and RNNs
  \item Computational Considerations
  \item Geometric Interpretation of Attention
  \item Matrix Formulation of Attention
  \item Attention as a Kernel Operator
  \item Continuous / Integral Operator View
  \item Unifying Interpretation
  \item A Unifying Perspective
  \item Outlook
\end{itemize}
\end{frame}

\subsection{Motivation}
\begin{frame}[t,allowframebreaks]{Motivation}
\small
In the previous lecture, we studied sequence models based on recurrence.

While recurrent networks can model sequential data, they suffer from key limitations:

\begin{itemize}
    \item difficulty capturing long-range dependencies,
    \item sequential computation limits parallelism,
    \item information is compressed into a single hidden state.
\end{itemize}

\textbf{Key question.}\\ Can we design models that directly relate all elements of a sequence?

\textbf{Guiding idea.}\\ Attention mechanisms allow each element of a sequence to interact with all others.
\end{frame}

\subsection{Basic Attention Mechanism}
\begin{frame}[t,allowframebreaks]{Basic Attention Mechanism}
\small
Attention computes a weighted combination of values.

Given:
\begin{itemize}
    \item queries $q$,
    \item keys $k_i$,
    \item values $v_i$,
\end{itemize}

define attention weights:
\[
\alpha_i = \frac{\exp(q \cdot k_i)}{\sum_j \exp(q \cdot k_j)}.
\]

Output:
\[
\text{Attention}(q, K, V) = \sum_i \alpha_i v_i.
\]

\textbf{Interpretation.}\\ \begin{itemize}
    \item keys determine relevance
    \item values carry information
    \item query selects what to focus on
\end{itemize}

\textbf{Insight.}\\ Attention dynamically weights information based on relevance.
\end{frame}

\subsection{Self-Attention}
\begin{frame}[t,allowframebreaks]{Self-Attention}
\small
In self-attention, queries, keys, and values are derived from the same sequence.

Given inputs $\{x_1, \dots, x_T\}$:

\[
q_i = W_Q x_i, \quad k_i = W_K x_i, \quad v_i = W_V x_i.
\]

Each position attends to all others:
\[
y_i = \sum_{j=1}^{T} \alpha_{ij} v_j.
\]

\textbf{Key properties.}\\ \begin{itemize}
    \item captures global dependencies
    \item fully parallelizable
\end{itemize}

\textbf{Insight.}\\ Self-attention replaces sequential processing with global interaction.
\end{frame}

\subsection{Scaled Dot-Product Attention}
\begin{frame}[t,allowframebreaks]{Scaled Dot-Product Attention}
\small
To improve stability, attention uses scaling:

\[
\alpha_{ij} =
\frac{\exp\left(\frac{q_i \cdot k_j}{\sqrt{d}}\right)}
{\sum_{k} \exp\left(\frac{q_i \cdot k_k}{\sqrt{d}}\right)}.
\]

\textbf{Reason.}\\ \begin{itemize}
    \item prevents large dot products
    \item stabilizes gradients
\end{itemize}

\textbf{Insight.}\\ Scaling ensures numerically stable similarity computation.
\end{frame}

\subsection{Multi-Head Attention}
\begin{frame}[t,allowframebreaks]{Multi-Head Attention}
\small
Instead of a single attention mechanism, we use multiple heads:

\[
\text{MultiHead}(Q,K,V) =
\text{Concat}(\text{head}_1, \dots, \text{head}_h) W_O.
\]

\textbf{Interpretation.}\\ \begin{itemize}
    \item each head learns different relationships
    \item captures multiple types of dependencies
\end{itemize}

\textbf{Insight.}\\ Multi-head attention enriches representation by combining diverse interactions.
\end{frame}

\subsection{Comparison with CNNs and RNNs}
\begin{frame}[t,allowframebreaks]{Comparison with CNNs and RNNs}
\small
\textbf{CNNs:}\\ \begin{itemize}
    \item local receptive fields
    \item limited long-range interaction
\end{itemize}

\textbf{RNNs:}\\ \begin{itemize}
    \item sequential processing
    \item information bottleneck
\end{itemize}

\textbf{Attention:}\\ \begin{itemize}
    \item global receptive field
    \item direct interactions between all positions
\end{itemize}

\textbf{Insight.}\\ Attention removes structural constraints of locality and recurrence.
\end{frame}

\subsection{Computational Considerations}
\begin{frame}[t,allowframebreaks]{Computational Considerations}
\small
\textbf{Complexity.}\\ \begin{itemize}
    \item attention requires $O(T^2)$ operations
\end{itemize}

\textbf{Tradeoffs.}\\ \begin{itemize}
    \item more expressive interactions
    \item higher computational cost
\end{itemize}

\textbf{Insight.}\\ Expressivity comes at a computational cost.
\end{frame}

\subsection{Geometric Interpretation of Attention}
\begin{frame}[t,allowframebreaks]{Geometric Interpretation of Attention}
\small
Attention can be interpreted geometrically as measuring similarity in a feature space.

Recall that attention weights are computed via dot products:
\[
\alpha_{ij} \propto \exp\left( \frac{q_i \cdot k_j}{\sqrt{d}} \right).
\]

\textbf{Geometric view.}\\ \begin{itemize}
    \item $q_i$ and $k_j$ are vectors in $\mathbb{R}^d$
    \item $q_i \cdot k_j$ measures alignment (similarity)
\end{itemize}

\textbf{Normalized interpretation.}\\ If vectors are normalized:
\[
q_i \cdot k_j = \|q_i\|\|k_j\|\cos\theta_{ij},
\]
so attention depends on the angle between vectors.

\textbf{Interpretation.}\\ \begin{itemize}
    \item large weight $\Rightarrow$ vectors aligned
    \item small weight $\Rightarrow$ vectors orthogonal or opposed
\end{itemize}

\textbf{Insight.}\\ Attention selects information based on geometric similarity in representation space.
\end{frame}

\subsection{Matrix Formulation of Attention}
\begin{frame}[t,allowframebreaks]{Matrix Formulation of Attention}
\small
Let:
\[
Q = [q_1, \dots, q_T]^\top, \quad
K = [k_1, \dots, k_T]^\top, \quad
V = [v_1, \dots, v_T]^\top.
\]

\textbf{Attention matrix.}\\ \[
A = \text{softmax}\left( \frac{QK^\top}{\sqrt{d}} \right).
\]

\textbf{Output.}\\ \[
Y = A V.
\]

\textbf{Interpretation.}\\ \begin{itemize}
    \item $QK^\top$ computes all pairwise similarities
    \item softmax produces a row-stochastic matrix
    \item each row of $A$ defines a weighted average over values
\end{itemize}

\textbf{Key observation.}\\ \begin{itemize}
    \item attention is a data-dependent linear operator
    \item weights depend on the input itself
\end{itemize}

\textbf{Insight.}\\ Self-attention performs adaptive, input-dependent linear transformations.
\end{frame}

\subsection{Attention as a Kernel Operator}
\begin{frame}[t,allowframebreaks]{Attention as a Kernel Operator}
\small
The attention mechanism resembles kernel methods.

Define a similarity function:
\[
k(x_i, x_j) = \exp\left( \frac{q_i \cdot k_j}{\sqrt{d}} \right).
\]

\textbf{Interpretation.}\\ \begin{itemize}
    \item attention computes a normalized kernel
    \item similar to kernel regression or smoothing
\end{itemize}

\textbf{Connection.}\\ \[
y_i = \sum_j \alpha_{ij} v_j
\]
resembles:
\[
f(x_i) = \sum_j k(x_i, x_j) y_j.
\]

\textbf{Key difference.}\\ \begin{itemize}
    \item classical kernels are fixed
    \item attention kernels are learned and data-dependent
\end{itemize}

\textbf{Insight.}\\ Attention generalizes kernel methods by learning the similarity function.
\end{frame}

\subsection{Continuous / Integral Operator View}
\begin{frame}[t,allowframebreaks]{Continuous / Integral Operator View}
\small
In the limit of large sequences, attention can be viewed as an integral operator.

Let $x(t)$ be a continuous input and define:
\[
y(t) = \int k(t,s) v(s) \, ds.
\]

\textbf{Interpretation.}\\ \begin{itemize}
    \item $k(t,s)$ measures interaction between positions $t$ and $s$
    \item $y(t)$ aggregates information across the domain
\end{itemize}

\textbf{Connection.}\\ \begin{itemize}
    \item discrete attention $\rightarrow$ finite-dimensional approximation
    \item continuous limit $\rightarrow$ operator on function spaces
\end{itemize}

\textbf{Relation to functional analysis.}\\ \begin{itemize}
    \item attention acts as a linear integral operator
    \item kernel defines the operator
\end{itemize}

\textbf{Insight.}\\ Attention can be viewed as a learned operator acting on functions.
\end{frame}

\subsection{Unifying Interpretation}
\begin{frame}[t,allowframebreaks]{Unifying Interpretation}
\small
From these perspectives, attention can be understood as:

\begin{itemize}
    \item \textbf{Geometric:}\\ similarity in feature space
    \item \textbf{Algebraic:}\\ matrix-based linear transformation
    \item \textbf{Kernel-based:}\\ adaptive similarity operator
    \item \textbf{Analytical:}\\ integral operator on functions
\end{itemize}

\textbf{Core idea.}\\ Attention is a flexible mechanism for modeling interactions across inputs.

\textbf{Key insight.}\\ Modern architectures can be interpreted through multiple mathematical lenses, revealing deeper structure behind empirical success.
\end{frame}

\subsection{A Unifying Perspective}
\begin{frame}[t,allowframebreaks]{A Unifying Perspective}
\small
Attention mechanisms provide a new paradigm:

\begin{itemize}
    \item representations are computed through interactions,
    \item relationships between elements are learned dynamically.
\end{itemize}

\textbf{Core idea.}\\ Learning involves selecting and combining relevant information across the entire input.

\textbf{Key insight.}\\ Attention transforms representation learning into interaction modeling.
\end{frame}

\subsection{Outlook}
\begin{frame}[t,allowframebreaks]{Outlook}
\small
In this lecture, we introduced:

\begin{itemize}
    \item attention mechanisms,
    \item self-attention,
    \item scaled dot-product attention,
    \item multi-head attention.
\end{itemize}

These ideas form the foundation of transformer architectures.

In the next lecture, we study transformers, which build entirely on attention mechanisms and enable large-scale models.

\textbf{Guiding principle.}\\ Modeling interactions directly can overcome limitations of traditional architectures.
\end{frame}

\section{Lecture 24: Transformers}
\SectionDivider{Lecture 24: Transformers}
\begin{frame}[t,allowframebreaks]{Outline}
\small
\begin{itemize}
  \item Motivation
  \item Transformer Architecture
  \item Positional Encoding
  \item Residual Connections and Normalization
  \item Computational Characteristics
  \item Scaling Laws
  \item Transformers in Practice
  \item Connection to Previous Concepts
  \item A Unifying Perspective
  \item Outlook
\end{itemize}
\end{frame}

\subsection{Motivation}
\begin{frame}[t,allowframebreaks]{Motivation}
\small
In the previous lecture, we introduced attention as a mechanism for modeling interactions across sequences.

\textbf{Key question.}\\ Can we build an entire architecture based solely on attention?

\textbf{Answer.}\\ Transformers replace recurrence and convolution with attention, enabling scalable and highly expressive models.

\textbf{Guiding idea.}\\ Representations are computed through repeated layers of attention and simple transformations.
\end{frame}

\subsection{Transformer Architecture}
\begin{frame}[t,allowframebreaks]{Transformer Architecture}
\small
A transformer is composed of stacked layers, each containing:

\begin{itemize}
    \item multi-head self-attention
    \item feedforward network
    \item residual connections
    \item normalization
\end{itemize}

\textbf{Layer structure.}\\ Given input $x$:

\[
z = \text{SelfAttention}(x),
\quad
x' = x + z
\]

\[
\tilde{x} = \text{LayerNorm}(x')
\]

\[
u = \text{FFN}(\tilde{x}),
\quad
y = \tilde{x} + u
\]

\[
\text{output} = \text{LayerNorm}(y)
\]

\textbf{Feedforward network.}\\ \[
\text{FFN}(x) = \sigma(xW_1 + b_1) W_2 + b_2.
\]

\textbf{Insight.}\\ Transformers alternate between global interaction (attention) and local transformation (FFN).
\end{frame}

\subsection{Positional Encoding}
\begin{frame}[t,allowframebreaks]{Positional Encoding}
\small
Attention is permutation-invariant, so positional information must be added.

\textbf{Sinusoidal encoding.}\\ \[
\text{PE}(t, 2k) = \sin\left(\frac{t}{10000^{2k/d}}\right), \quad
\text{PE}(t, 2k+1) = \cos\left(\frac{t}{10000^{2k/d}}\right).
\]

\textbf{Purpose.}\\ \begin{itemize}
    \item encode order information
    \item allow model to distinguish positions
\end{itemize}

\textbf{Insight.}\\ Transformers separate content (tokens) from position (encoding).
\end{frame}

\subsection{Residual Connections and Normalization}
\begin{frame}[t,allowframebreaks]{Residual Connections and Normalization}
\small
Transformers rely heavily on:

\begin{itemize}
    \item residual connections
    \item layer normalization
\end{itemize}

\textbf{Role.}\\ \begin{itemize}
    \item stabilize deep architectures
    \item ensure gradient flow
\end{itemize}

\textbf{Connection to earlier lectures.}\\ \begin{itemize}
    \item residuals $\rightarrow$ enable depth (Lecture 20)
    \item normalization $\rightarrow$ stabilize training (Lecture 18)
\end{itemize}

\textbf{Insight.}\\ Transformers integrate multiple stabilization techniques.
\end{frame}

\subsection{Computational Characteristics}
\begin{frame}[t,allowframebreaks]{Computational Characteristics}
\small
\textbf{Advantages.}\\ \begin{itemize}
    \item full parallelism across sequence
    \item direct modeling of long-range dependencies
\end{itemize}

\textbf{Cost.}\\ \begin{itemize}
    \item attention requires $O(T^2)$ computation
\end{itemize}

\textbf{Tradeoff.}\\ \begin{itemize}
    \item higher expressivity vs computational cost
\end{itemize}

\textbf{Insight.}\\ Transformers trade efficiency for expressive global interaction.
\end{frame}

\subsection{Scaling Laws}
\begin{frame}[t,allowframebreaks]{Scaling Laws}
\small
Modern transformer-based models exhibit predictable scaling behavior.

\textbf{Empirical observation.}\\ \begin{itemize}
    \item performance improves with:
    \begin{itemize}
        \item model size (parameters)
        \item dataset size
        \item compute
    \end{itemize}
\end{itemize}

\textbf{Scaling law (informal).}\\ \[
\text{Loss} \propto N^{-\alpha},
\]

where $N$ represents scale (data, parameters, compute).

\textbf{Implication.}\\ \begin{itemize}
    \item larger models trained on more data perform better
\end{itemize}

\textbf{Insight.}\\ Performance follows smooth scaling trends rather than abrupt improvements.
\end{frame}

\subsection{Transformers in Practice}
\begin{frame}[t,allowframebreaks]{Transformers in Practice}
\small
Transformers are the foundation of modern AI systems:

\begin{itemize}
    \item large language models (LLMs)
    \item vision transformers
    \item multimodal models
\end{itemize}

\textbf{Key property.}\\ \begin{itemize}
    \item flexible across domains
    \item unified architecture for different data types
\end{itemize}

\textbf{Insight.}\\ Transformers provide a general-purpose architecture for learning representations.
\end{frame}

\subsection{Connection to Previous Concepts}
\begin{frame}[t,allowframebreaks]{Connection to Previous Concepts}
\small
Transformers unify ideas developed throughout the course:

\begin{itemize}
    \item kernel methods $\rightarrow$ attention as learned similarity
    \item deep networks $\rightarrow$ compositional structure
    \item residual connections $\rightarrow$ stable depth
    \item normalization $\rightarrow$ controlled training dynamics
\end{itemize}

\textbf{Insight.}\\ Transformers synthesize multiple strands of machine learning into a single framework.
\end{frame}

\subsection{A Unifying Perspective}
\begin{frame}[t,allowframebreaks]{A Unifying Perspective}
\small
Transformers can be viewed as:

\begin{itemize}
    \item \textbf{Interaction networks:}\\ modeling pairwise relationships
    \item \textbf{Adaptive operators:}\\ data-dependent transformations
    \item \textbf{Scalable architectures:}\\ benefiting from scale
\end{itemize}

\textbf{Core idea.}\\ Learning is achieved by iteratively refining representations through global interactions.

\textbf{Key insight.}\\ The success of modern AI arises from combining expressive architectures with large-scale optimization.
\end{frame}

\subsection{Outlook}
\begin{frame}[t,allowframebreaks]{Outlook}
\small
In this lecture, we studied:

\begin{itemize}
    \item transformer architecture,
    \item positional encoding,
    \item residual and normalization mechanisms,
    \item scaling laws and modern applications.
\end{itemize}

Transformers represent a major milestone in machine learning.

In the following part of the course, we connect these ideas to generative modeling, large-scale systems, and theoretical perspectives.

\textbf{Guiding principle.}\\ Modern machine learning systems are shaped by both architectural design and scale.
\end{frame}

\section{Tutorial 12: Attention Visualization}
\SectionDivider{Tutorial 12: Attention Visualization}
\begin{frame}[t,allowframebreaks]{Outline}
\small
\begin{itemize}
  \item Objectives
  \item Exercise 1: Attention Matrix as a Map of Interactions
  \item Exercise 2: Interpreting a Simple Attention Pattern
  \item Exercise 3: Global Attention Pattern
  \item Exercise 4: Multi-Head Attention Interpretation
  \item Exercise 5: Attention and Positional Structure
  \item Exercise 6: Comparison with RNNs and CNNs
  \item Exercise 7: Attention as a Kernel Matrix
  \item Exercise 8: Visualization Thought Experiment
  \item Exercise 9 (Discussion): Are Attention Weights Explanations?
  \item Summary
\end{itemize}
\end{frame}

\subsection{Objectives}
\begin{frame}[t,allowframebreaks]{Objectives}
\small
This tutorial develops intuition for how attention behaves inside transformer models:

\begin{itemize}
    \item Understand attention weights as interaction patterns
    \item Interpret self-attention matrices visually
    \item Relate attention patterns to linguistic or structural roles
    \item Connect local and global dependencies to model behavior
\end{itemize}

\textbf{Focus.}\\ Attention is not just a formula; it is a structured map of interactions between elements.
\end{frame}

\subsection{Exercise 1: Attention Matrix as a Map of Interactions}
\begin{frame}[t,allowframebreaks]{Exercise 1: Attention Matrix as a Map of Interactions}
\small
Consider a sequence of length $T$ with self-attention matrix:
\[
A = \text{softmax}\left(\frac{QK^\top}{\sqrt{d}}\right).
\]

\textbf{Tasks.}\\ \begin{enumerate}
    \item What are the dimensions of $A$?
    \item What does the entry $A_{ij}$ represent?
    \item Why does each row sum to $1$?
\end{enumerate}

\textbf{Solution.}\\ \begin{enumerate}
    \item
    \[
    A \in \mathbb{R}^{T \times T}.
    \]

    \item $A_{ij}$ is the attention weight from position $i$ to position $j$; it measures how much token $i$ attends to token $j$.

    \item Each row is the output of a softmax, so:
    \[
    \sum_{j=1}^T A_{ij} = 1.
    \]
\end{enumerate}

\textbf{Key takeaway.}\\ Attention matrices are row-stochastic interaction maps.
\end{frame}

\subsection{Exercise 2: Interpreting a Simple Attention Pattern}
\begin{frame}[t,allowframebreaks]{Exercise 2: Interpreting a Simple Attention Pattern}
\small
Suppose a sentence has 5 tokens, and one attention head produces:

\[
A =
\begin{pmatrix}
0.7 & 0.2 & 0.1 & 0.0 & 0.0 \\
0.1 & 0.6 & 0.2 & 0.1 & 0.0 \\
0.0 & 0.2 & 0.5 & 0.2 & 0.1 \\
0.0 & 0.1 & 0.2 & 0.6 & 0.1 \\
0.0 & 0.0 & 0.1 & 0.2 & 0.7
\end{pmatrix}.
\]

\textbf{Tasks.}\\ \begin{enumerate}
    \item Describe the qualitative pattern.
    \item What kind of dependency does this head emphasize?
\end{enumerate}

\textbf{Solution.}\\ \begin{enumerate}
    \item Most of the mass is near the diagonal.

    \item This head emphasizes local context: each token mostly attends to itself and nearby tokens.
\end{enumerate}

\textbf{Insight.}\\ Some attention heads behave like local filters, similar to convolutional receptive fields.
\end{frame}

\subsection{Exercise 3: Global Attention Pattern}
\begin{frame}[t,allowframebreaks]{Exercise 3: Global Attention Pattern}
\small
Now consider:

\[
A =
\begin{pmatrix}
0.1 & 0.1 & 0.1 & 0.6 & 0.1 \\
0.0 & 0.2 & 0.1 & 0.6 & 0.1 \\
0.1 & 0.1 & 0.1 & 0.6 & 0.1 \\
0.1 & 0.1 & 0.1 & 0.5 & 0.2 \\
0.1 & 0.1 & 0.1 & 0.6 & 0.1
\end{pmatrix}.
\]

\textbf{Tasks.}\\ \begin{enumerate}
    \item What is the dominant attended position?
    \item What might this represent in language?
\end{enumerate}

\textbf{Solution.}\\ \begin{enumerate}
    \item Position 4 is dominant.

    \item This may correspond to a syntactically or semantically central token, such as a main verb, subject, or separator token.
\end{enumerate}

\textbf{Insight.}\\ Attention can capture global structure, not only local neighborhoods.
\end{frame}

\subsection{Exercise 4: Multi-Head Attention Interpretation}
\begin{frame}[t,allowframebreaks]{Exercise 4: Multi-Head Attention Interpretation}
\small
Transformers use multiple attention heads:
\[
\text{MultiHead}(Q,K,V) = \text{Concat}(\text{head}_1, \dots, \text{head}_h)W_O.
\]

\textbf{Question.}\\ Why is multi-head attention useful?

\textbf{Discussion.}\\ \begin{itemize}
    \item Different heads can specialize in different patterns
    \item One head may capture local context
    \item Another may capture long-range dependencies
    \item Another may track punctuation or structural boundaries
\end{itemize}

\textbf{Key takeaway.}\\ Multi-head attention enables diverse relational representations.
\end{frame}

\subsection{Exercise 5: Attention and Positional Structure}
\begin{frame}[t,allowframebreaks]{Exercise 5: Attention and Positional Structure}
\small
Since self-attention alone is permutation-invariant, transformers add positional encodings.

\textbf{Tasks.}\\ \begin{enumerate}
    \item Why can attention alone not distinguish order?
    \item How does positional encoding fix this?
\end{enumerate}

\textbf{Solution.}\\ \begin{enumerate}
    \item Without positional information, the model only sees a set of vectors and their pairwise similarities.

    \item Positional encodings inject order information into token representations, so attention can depend on both content and position.
\end{enumerate}

\textbf{Insight.}\\ Attention captures interactions; positional encoding tells the model where those interactions occur.
\end{frame}

\subsection{Exercise 6: Comparison with RNNs and CNNs}
\begin{frame}[t,allowframebreaks]{Exercise 6: Comparison with RNNs and CNNs}
\small
\textbf{Question.}\\ How does attention differ from CNNs and RNNs in how it processes structure?

\textbf{Discussion.}\\ \begin{itemize}
    \item CNN:
    \begin{itemize}
        \item local interactions
        \item fixed receptive structure
    \end{itemize}

    \item RNN:
    \begin{itemize}
        \item sequential memory
        \item hidden-state bottleneck
    \end{itemize}

    \item Attention:
    \begin{itemize}
        \item global interactions
        \item data-dependent connectivity
    \end{itemize}
\end{itemize}

\textbf{Insight.}\\ Attention replaces fixed or sequential structure with adaptive interaction patterns.
\end{frame}

\subsection{Exercise 7: Attention as a Kernel Matrix}
\begin{frame}[t,allowframebreaks]{Exercise 7: Attention as a Kernel Matrix}
\small
Recall the matrix form:
\[
A = \text{softmax}\left(\frac{QK^\top}{\sqrt{d}}\right).
\]

\textbf{Task.}\\ Explain why $QK^\top$ resembles a kernel or similarity matrix.

\textbf{Solution.}\\ \begin{itemize}
    \item Each entry is a dot-product similarity
    \item Similar to Gram matrices in kernel methods
    \item Difference: the representation itself is learned and input-dependent
\end{itemize}

\textbf{Insight.}\\ Attention can be viewed as a learned, adaptive kernel operator.
\end{frame}

\subsection{Exercise 8: Visualization Thought Experiment}
\begin{frame}[t,allowframebreaks]{Exercise 8: Visualization Thought Experiment}
\small
Suppose you plot attention matrices from different layers.

\textbf{Question.}\\ What patterns might you expect from shallow vs deep layers?

\textbf{Discussion.}\\ \begin{itemize}
    \item Early layers:
    \begin{itemize}
        \item more local or lexical patterns
    \end{itemize}

    \item Later layers:
    \begin{itemize}
        \item more global, semantic, or task-specific interactions
    \end{itemize}
\end{itemize}

\textbf{Insight.}\\ Attention patterns evolve hierarchically across depth.
\end{frame}

\subsection{Exercise 9 (Discussion): Are Attention Weights Explanations?}
\begin{frame}[t,allowframebreaks]{Exercise 9 (Discussion): Are Attention Weights Explanations?}
\small
\textbf{Question.}\\ Can attention maps be interpreted as explanations of model decisions?

\textbf{Discussion.}\\ \begin{itemize}
    \item They show interaction patterns
    \item But they do not fully explain the model’s computation
    \item Feedforward layers and residual paths also matter
\end{itemize}

\textbf{Insight.}\\ Attention visualization is informative, but not a complete explanation.
\end{frame}

\subsection{Summary}
\begin{frame}[t,allowframebreaks]{Summary}
\small
\begin{itemize}
    \item Attention matrices visualize interactions between tokens
    \item Different heads can capture different structural roles
    \item Attention can express both local and global dependencies
    \item Transformers combine interaction modeling with positional information
\end{itemize}

\textbf{Preparation for the next part of the course.}\\ We now move from architectures to generative modeling and large-scale AI systems, where attention and transformers play a central role.
\end{frame}

\end{document}
